%% Generated by Sphinx.
\def\sphinxdocclass{report}
\IfFileExists{luatex85.sty}
 {\RequirePackage{luatex85}}
 {\ifdefined\luatexversion\ifnum\luatexversion>84\relax
  \PackageError{sphinx}
  {** With this LuaTeX (\the\luatexversion),Sphinx requires luatex85.sty **}
  {** Add the LaTeX package luatex85 to your TeX installation, and try again **}
  \endinput\fi\fi}
\documentclass[letterpaper,10pt,spanish]{sphinxmanual}
\ifdefined\pdfpxdimen
   \let\sphinxpxdimen\pdfpxdimen\else\newdimen\sphinxpxdimen
\fi \sphinxpxdimen=.75bp\relax
\ifdefined\pdfimageresolution
    \pdfimageresolution= \numexpr \dimexpr1in\relax/\sphinxpxdimen\relax
\fi
%% let collapsible pdf bookmarks panel have high depth per default
\PassOptionsToPackage{bookmarksdepth=5}{hyperref}
%% turn off hyperref patch of \index as sphinx.xdy xindy module takes care of
%% suitable \hyperpage mark-up, working around hyperref-xindy incompatibility
\PassOptionsToPackage{hyperindex=false}{hyperref}
%% memoir class requires extra handling
\makeatletter\@ifclassloaded{memoir}
{\ifdefined\memhyperindexfalse\memhyperindexfalse\fi}{}\makeatother

\PassOptionsToPackage{borderless}{sphinx}
\PassOptionsToPackage{colorrows}{sphinx}

\PassOptionsToPackage{warn}{textcomp}

\catcode`^^^^00a0\active\protected\def^^^^00a0{\leavevmode\nobreak\ }
\usepackage{cmap}
\usepackage{fontspec}
\defaultfontfeatures[\rmfamily,\sffamily,\ttfamily]{}
\usepackage{amsmath,amssymb,amstext}
\usepackage{polyglossia}
\setmainlanguage{spanish}



\setmainfont{FreeSerif}[
  Extension      = .otf,
  UprightFont    = *,
  ItalicFont     = *Italic,
  BoldFont       = *Bold,
  BoldItalicFont = *BoldItalic
]
\setsansfont{FreeSans}[
  Extension      = .otf,
  UprightFont    = *,
  ItalicFont     = *Oblique,
  BoldFont       = *Bold,
  BoldItalicFont = *BoldOblique,
]
\setmonofont{FreeMono}[Scale=0.9,
  Extension      = .otf,
  UprightFont    = *,
  ItalicFont     = *Oblique,
  BoldFont       = *Bold,
  BoldItalicFont = *BoldOblique,
]



\usepackage[Sonny]{fncychap}
\ChNameVar{\Large\normalfont\sffamily}
\ChTitleVar{\Large\normalfont\sffamily}
\usepackage{sphinx}

\fvset{fontsize=auto}
\usepackage{geometry}


% Include hyperref last.
\usepackage{hyperref}
% Fix anchor placement for figures with captions.
\usepackage{hypcap}% it must be loaded after hyperref.
% Set up styles of URL: it should be placed after hyperref.
\urlstyle{same}


\usepackage{sphinxmessages}
\setcounter{tocdepth}{1}



\title{Libro didáctico para el estudiante}
\date{18 de mayo de 2026}
\release{}
\author{Julian Camilo Fonseca Romero}
\newcommand{\sphinxlogo}{\sphinxincludegraphics{logo_light.png}\par}
\renewcommand{\releasename}{}
\makeindex
\begin{document}

\pagestyle{empty}
\sphinxmaketitle
\pagestyle{plain}
\sphinxtableofcontents
\pagestyle{normal}
\phantomsection\label{\detokenize{index::doc}}


\sphinxAtStartPar
\sphinxincludegraphics{{aula115}.png}
\begin{quote}

\sphinxAtStartPar
El presente libro didáctico digital, \sphinxstylestrong{aula 115}, es una herramienta fundamental para estudiantes de educación básica y media de la \sphinxstyleemphasis{Institución Educativa John F. Kennedy}. Su propósito es potenciar el uso estratégico de herramientas ofimáticas libres, promoviendo no solo la productividad en el manejo de la información, sino también la autonomía e independencia tecnológica necesaria en los campos laboral y académico del siglo XXI.
\end{quote}


\begin{savenotes}\sphinxattablestart
\sphinxthistablewithglobalstyle
\centering
\begin{tabular}[t]{*{2}{\X{1}{2}}}
\sphinxtoprule
\sphinxtableatstartofbodyhook

\begin{savenotes}\sphinxattablestart
\sphinxthistablewithglobalstyle
\centering
\begin{tabulary}{\linewidth}[t]{T}
\sphinxtoprule
\sphinxtableatstartofbodyhook

\begin{sphinxfigure-in-table}
\centering

\noindent\sphinxincludegraphics[width=180\sphinxpxdimen]{{camilo.fonseca}.png}
\end{sphinxfigure-in-table}\relax
\\
\sphinxbottomrule
\end{tabulary}
\sphinxtableafterendhook\par
\sphinxattableend\end{savenotes}
&

\begin{savenotes}\sphinxattablestart
\sphinxthistablewithglobalstyle
\centering
\begin{tabular}[t]{*{1}{\X{1}{1}}}
\sphinxtoprule
\sphinxtableatstartofbodyhook
\sphinxAtStartPar
\sphinxstylestrong{DOCENTE:} Julian Camilo Fonseca Romero

\sphinxAtStartPar
\sphinxstylestrong{email:} camilo.fonseca@jfkennedy.edu.co
\begin{itemize}
\item {} 
\sphinxAtStartPar
Ingeniero en Mecatrónica (2012)

\item {} 
\sphinxAtStartPar
Especialista en Seguridad de la Información (2015)

\item {} 
\sphinxAtStartPar
Magister en Ingeniería de Software y Sistemas Informáticos (2017)

\end{itemize}
\\
\sphinxbottomrule
\end{tabular}
\sphinxtableafterendhook\par
\sphinxattableend\end{savenotes}
\\
\sphinxbottomrule
\end{tabular}
\sphinxtableafterendhook\par
\sphinxattableend\end{savenotes}

\begin{sphinxadmonition}{note}{PRINCIPIOS DE LA CLASE}
\begin{enumerate}
\sphinxsetlistlabels{\arabic}{enumi}{enumii}{}{.}%
\item {} 
\sphinxAtStartPar
“Se aprende haciendo”.

\item {} 
\sphinxAtStartPar
“Se aprende enseñando”.

\item {} 
\sphinxAtStartPar
“Todo se sustenta”.

\end{enumerate}
\end{sphinxadmonition}

\begin{sphinxadmonition}{note}{SECUENCIA DIDÁCTICA (MODELO ACTIVO)}
\begin{enumerate}
\sphinxsetlistlabels{\arabic}{enumi}{enumii}{}{.}%
\item {} 
\sphinxAtStartPar
Exploración.

\item {} 
\sphinxAtStartPar
Estructuración.

\item {} 
\sphinxAtStartPar
Actividad práctica.

\item {} 
\sphinxAtStartPar
Transferencia.

\end{enumerate}
\end{sphinxadmonition}

\sphinxstepscope


\chapter{Información}
\label{\detokenize{informacion:informacion}}\label{\detokenize{informacion::doc}}
\sphinxAtStartPar
La siguiente información es útil para el estudiante.


\bigskip\hrule\bigskip



\section{Útiles escolares}
\label{\detokenize{informacion:utiles-escolares}}\begin{itemize}
\item {} 
\sphinxAtStartPar
Cuaderno cocido (100 hojas y cuadriculado).

\item {} 
\sphinxAtStartPar
Lapiz, Borrador, Tajalapiz.

\item {} 
\sphinxAtStartPar
Esferos azul, negro.

\item {} 
\sphinxAtStartPar
Corrector.

\end{itemize}


\bigskip\hrule\bigskip



\section{Porcentaje de notas \%}
\label{\detokenize{informacion:porcentaje-de-notas}}
\begin{sphinxadmonition}{note}{50\% + 40\% + 10\% = 100\%}
\begin{itemize}
\item {} 
\sphinxAtStartPar
\sphinxstylestrong{50\% COGNITIVA:} Guías

\item {} 
\sphinxAtStartPar
\sphinxstylestrong{40\% PROCEDIMENTAL:} Resultado general prueba INSTRUIMOS

\item {} 
\sphinxAtStartPar
\sphinxstylestrong{10\% ACTITUDINAL} = 5\% PERFIL Kennediano \sphinxstylestrong{+} 5\% Trabajo en equipo

\end{itemize}
\end{sphinxadmonition}


\bigskip\hrule\bigskip



\section{Equivalencia calificación INSTRUIMOS}
\label{\detokenize{informacion:equivalencia-calificacion-instruimos}}
\sphinxAtStartPar
La siguiente tabla muestra el equivalente del puntaje de nota Instruimos (1\sphinxhyphen{}100) respecto al valor de la nota en la Institución Educativa (1\sphinxhyphen{}5).


\begin{savenotes}\sphinxattablestart
\sphinxthistablewithglobalstyle
\centering
\begin{tabulary}{\linewidth}[t]{TT}
\sphinxtoprule
\sphinxstyletheadfamily 
\sphinxAtStartPar
Puntaje INSTRUIMOS
&\sphinxstyletheadfamily 
\sphinxAtStartPar
Nota IEJFK
\\
\sphinxmidrule
\sphinxtableatstartofbodyhook
\sphinxAtStartPar
10\sphinxhyphen{}15
&
\sphinxAtStartPar
10
\\
\sphinxhline
\sphinxAtStartPar
16\sphinxhyphen{}25
&
\sphinxAtStartPar
15
\\
\sphinxhline
\sphinxAtStartPar
26\sphinxhyphen{}35
&
\sphinxAtStartPar
20
\\
\sphinxhline
\sphinxAtStartPar
36\sphinxhyphen{}40
&
\sphinxAtStartPar
25
\\
\sphinxhline
\sphinxAtStartPar
41\sphinxhyphen{}45
&
\sphinxAtStartPar
29
\\
\sphinxhline
\sphinxAtStartPar
46\sphinxhyphen{}55
&
\sphinxAtStartPar
30
\\
\sphinxhline
\sphinxAtStartPar
56\sphinxhyphen{}65
&
\sphinxAtStartPar
35
\\
\sphinxhline
\sphinxAtStartPar
66\sphinxhyphen{}75
&
\sphinxAtStartPar
40
\\
\sphinxhline
\sphinxAtStartPar
76\sphinxhyphen{}85
&
\sphinxAtStartPar
45
\\
\sphinxhline
\sphinxAtStartPar
86\sphinxhyphen{}95
&
\sphinxAtStartPar
48
\\
\sphinxhline
\sphinxAtStartPar
96\sphinxhyphen{}100
&
\sphinxAtStartPar
50
\\
\sphinxbottomrule
\end{tabulary}
\sphinxtableafterendhook\par
\sphinxattableend\end{savenotes}


\bigskip\hrule\bigskip



\section{Sello docente}
\label{\detokenize{informacion:sello-docente}}\begin{itemize}
\item {} 
\sphinxAtStartPar
Estudiante coloca la fecha.

\item {} 
\sphinxAtStartPar
La \sphinxstylestrong{nota} puede ser cualitativa o cuantitativa.

\item {} 
\sphinxAtStartPar
Por lo general, cada sello se marca después de terminar una guía.

\end{itemize}


\begin{savenotes}\sphinxattablestart
\sphinxthistablewithglobalstyle
\centering
\begin{tabulary}{\linewidth}[t]{T}
\sphinxtoprule
\sphinxtableatstartofbodyhook

\begin{sphinxfigure-in-table}
\centering
\capstart
\noindent\sphinxincludegraphics[width=400\sphinxpxdimen]{{sello}.png}
\sphinxfigcaption{Sello docente}\label{\detokenize{informacion:id1}}\end{sphinxfigure-in-table}\relax
\\
\sphinxbottomrule
\end{tabulary}
\sphinxtableafterendhook\par
\sphinxattableend\end{savenotes}


\bigskip\hrule\bigskip



\section{Acceso al portal dentro y fuera del aula}
\label{\detokenize{informacion:acceso-al-portal-dentro-y-fuera-del-aula}}
\sphinxAtStartPar
Los \sphinxstylestrong{estudiantes} pueden acceder al \sphinxstylestrong{portal web} desde cualquier lugar siempre y cuando \sphinxstylestrong{tengan acceso a INTERNET}.

\sphinxAtStartPar
Si los estudiantes también puede acceder al \sphinxstylestrong{portal web} por medio de la INTRANET siempre y cuando estén conectados a la red del aula.
\begin{itemize}
\item {} 
\sphinxAtStartPar
INTRANET: \sphinxurl{http://192.168.115.100}

\item {} 
\sphinxAtStartPar
INTERNET: \sphinxurl{https://aula115.readthedocs.io}

\end{itemize}


\bigskip\hrule\bigskip



\section{Modelo de almacenamiento Google WorkSpace®}
\label{\detokenize{informacion:modelo-de-almacenamiento-google-workspace}}
\sphinxAtStartPar
La siguiente tabla muestra los \sphinxstylestrong{límites de almacenamiento} de la cuenta institucional que cada integrante tiene (según su rol) en la \sphinxhref{https://www.jfkennedy.edu.co}{Institución Educativa John F. Kennedy}.


\begin{savenotes}\sphinxattablestart
\sphinxthistablewithglobalstyle
\centering
\begin{tabulary}{\linewidth}[t]{TTTTT}
\sphinxtoprule
\sphinxstyletheadfamily 
\sphinxAtStartPar
rol
&\sphinxstyletheadfamily 
\sphinxAtStartPar
ACTIVO
&\sphinxstyletheadfamily 
\sphinxAtStartPar
EGRESADO
&\sphinxstyletheadfamily 
\sphinxAtStartPar
RETIRADO
&\sphinxstyletheadfamily 
\sphinxAtStartPar
PENSIONADO
\\
\sphinxmidrule
\sphinxtableatstartofbodyhook
\sphinxAtStartPar
\sphinxstylestrong{ADMINISTRATIVO}
&
\sphinxAtStartPar
1TB
&
\sphinxAtStartPar
\sphinxstylestrong{N/A}
&
\sphinxAtStartPar
16GB
&
\sphinxAtStartPar
16GB
\\
\sphinxhline
\sphinxAtStartPar
\sphinxstylestrong{DOCENTE}
&
\sphinxAtStartPar
256GB
&
\sphinxAtStartPar
\sphinxstylestrong{N/A}
&
\sphinxAtStartPar
16GB
&
\sphinxAtStartPar
16GB
\\
\sphinxhline
\sphinxAtStartPar
\sphinxstylestrong{ESTUDIANTE}
&
\sphinxAtStartPar
16GB
&
\sphinxAtStartPar
16GB
&
\sphinxAtStartPar
\sphinxguilabel{X}
&
\sphinxAtStartPar
\sphinxstylestrong{N/A}
\\
\sphinxbottomrule
\end{tabulary}
\sphinxtableafterendhook\par
\sphinxattableend\end{savenotes}

\sphinxAtStartPar
\sphinxguilabel{X} = Suspensión de la cuenta


\bigskip\hrule\bigskip



\section{Acceso al aula 115 en contrajornada}
\label{\detokenize{informacion:acceso-al-aula-115-en-contrajornada}}
\sphinxAtStartPar
Para \sphinxstylestrong{asegurar y mantener la correcta funcionalidad de los equipos}, se debe tener un control y seguimiento en el aula. Es por eso que los estudiantes que deseen \sphinxstyleemphasis{\sphinxstylestrong{hacer uso del aula de informática 115 en contra\sphinxhyphen{}jornada}}, deben solicitar el acceso enviando un correo a:
\begin{quote}\begin{description}
\sphinxlineitem{Cuenta}
\sphinxAtStartPar
camilo.fonseca@jfkennedy.edu.co

\sphinxlineitem{Asunto}
\sphinxAtStartPar
Autorización acceso aula 115 en contra\sphinxhyphen{}jornada

\end{description}\end{quote}

\sphinxAtStartPar
Los horarios en contrajornada son:
\begin{itemize}
\item {} 
\sphinxAtStartPar
\sphinxstylestrong{Mañana:} 10AM\sphinxhyphen{}12PM

\item {} 
\sphinxAtStartPar
\sphinxstylestrong{Tarde:} 01PM\sphinxhyphen{}03PM

\end{itemize}


\bigskip\hrule\bigskip



\section{Acceso de estudiantes al PC}
\label{\detokenize{informacion:acceso-de-estudiantes-al-pc}}

\begin{savenotes}\sphinxattablestart
\sphinxthistablewithglobalstyle
\centering
\begin{tabular}[t]{*{1}{\X{1}{1}}}
\sphinxtoprule
\sphinxtableatstartofbodyhook
\begin{sphinxVerbatimintable}[commandchars=\\\{\},numbers=left,firstnumber=1,stepnumber=1]
\PYG{n}{cursoActual} \PYG{o}{=} \PYG{n+nb}{input}\PYG{p}{(}\PYG{l+s+s2}{\PYGZdq{}}\PYG{l+s+s2}{CURSO: }\PYG{l+s+s2}{\PYGZdq{}}\PYG{p}{)}

\PYG{k}{if} \PYG{p}{(}\PYG{n}{cursoActual} \PYG{o}{==} \PYG{l+s+s2}{\PYGZdq{}}\PYG{l+s+s2}{8A}\PYG{l+s+s2}{\PYGZdq{}} \PYG{o+ow}{or} \PYG{n}{cursoActual} \PYG{o}{==} \PYG{l+s+s2}{\PYGZdq{}}\PYG{l+s+s2}{8B}\PYG{l+s+s2}{\PYGZdq{}}\PYG{p}{)}\PYG{p}{:}
    \PYG{n}{usuarioEstudiante} \PYG{o}{=} \PYG{l+s+s2}{\PYGZdq{}}\PYG{l+s+s2}{usuario}\PYG{l+s+s2}{\PYGZdq{}}

\PYG{k}{elif} \PYG{p}{(}\PYG{n}{cursoActual} \PYG{o}{==} \PYG{l+s+s2}{\PYGZdq{}}\PYG{l+s+s2}{10A}\PYG{l+s+s2}{\PYGZdq{}}\PYG{p}{)}\PYG{p}{:}
    \PYG{n}{usuarioEstudiante} \PYG{o}{=} \PYG{l+s+s2}{\PYGZdq{}}\PYG{l+s+s2}{10A}\PYG{l+s+s2}{\PYGZdq{}}
    \PYG{n+nb}{print}\PYG{p}{(}\PYG{l+s+s2}{\PYGZdq{}}\PYG{l+s+s2}{dar contraseña de acceso a 10A}\PYG{l+s+s2}{\PYGZdq{}}\PYG{p}{)}

\PYG{k}{elif} \PYG{p}{(}\PYG{n}{cursoActual} \PYG{o}{==} \PYG{l+s+s2}{\PYGZdq{}}\PYG{l+s+s2}{10B}\PYG{l+s+s2}{\PYGZdq{}}\PYG{p}{)}\PYG{p}{:}
    \PYG{n}{usuarioEstudiante} \PYG{o}{=} \PYG{l+s+s2}{\PYGZdq{}}\PYG{l+s+s2}{10B}\PYG{l+s+s2}{\PYGZdq{}}
    \PYG{n+nb}{print}\PYG{p}{(}\PYG{l+s+s2}{\PYGZdq{}}\PYG{l+s+s2}{dar contraseña de acceso a 10B}\PYG{l+s+s2}{\PYGZdq{}}\PYG{p}{)}

\PYG{k}{elif} \PYG{p}{(}\PYG{n}{cursoActual} \PYG{o}{==} \PYG{l+s+s2}{\PYGZdq{}}\PYG{l+s+s2}{11A}\PYG{l+s+s2}{\PYGZdq{}}\PYG{p}{)}\PYG{p}{:}
    \PYG{n}{usuarioEstudiante} \PYG{o}{=} \PYG{l+s+s2}{\PYGZdq{}}\PYG{l+s+s2}{11A}\PYG{l+s+s2}{\PYGZdq{}}
    \PYG{n+nb}{print}\PYG{p}{(}\PYG{l+s+s2}{\PYGZdq{}}\PYG{l+s+s2}{dar contraseña de acceso a 11A}\PYG{l+s+s2}{\PYGZdq{}}\PYG{p}{)}

\PYG{k}{elif} \PYG{p}{(}\PYG{n}{cursoActual} \PYG{o}{==} \PYG{l+s+s2}{\PYGZdq{}}\PYG{l+s+s2}{11B}\PYG{l+s+s2}{\PYGZdq{}}\PYG{p}{)}\PYG{p}{:}
    \PYG{n}{usuarioEstudiante} \PYG{o}{=} \PYG{l+s+s2}{\PYGZdq{}}\PYG{l+s+s2}{11B}\PYG{l+s+s2}{\PYGZdq{}}
    \PYG{n+nb}{print}\PYG{p}{(}\PYG{l+s+s2}{\PYGZdq{}}\PYG{l+s+s2}{dar contraseña de acceso a 11B}\PYG{l+s+s2}{\PYGZdq{}}\PYG{p}{)}

\PYG{k}{else}\PYG{p}{:}
    \PYG{n}{usuarioEstudiante} \PYG{o}{=} \PYG{l+s+s2}{\PYGZdq{}}\PYG{l+s+s2}{NO INFO}\PYG{l+s+s2}{\PYGZdq{}}
\end{sphinxVerbatimintable}
\\
\sphinxbottomrule
\end{tabular}
\sphinxtableafterendhook\par
\sphinxattableend\end{savenotes}


\bigskip\hrule\bigskip



\section{Estructura carpeta de evidencias}
\label{\detokenize{informacion:estructura-carpeta-de-evidencias}}
\sphinxAtStartPar
Los \sphinxstylestrong{estudiantes} tienen la siguiente estructura de \sphinxstylestrong{carpeta de evidencias}.

\sphinxAtStartPar
Los \sphinxstylestrong{archivos} que se trabajen en clase, deben ser clasificados en las carpetas según el periodo y la asignatura.


\begin{savenotes}\sphinxattablestart
\sphinxthistablewithglobalstyle
\centering
\begin{tabulary}{\linewidth}[t]{T}
\sphinxtoprule
\sphinxtableatstartofbodyhook

\begin{sphinxfigure-in-table}
\centering
\capstart
\noindent\sphinxincludegraphics[width=200\sphinxpxdimen]{{image017}.png}
\sphinxfigcaption{Carpeta de evidencias}\label{\detokenize{informacion:id2}}\end{sphinxfigure-in-table}\relax
\\
\sphinxbottomrule
\end{tabulary}
\sphinxtableafterendhook\par
\sphinxattableend\end{savenotes}


\bigskip\hrule\bigskip



\section{Actividades inicio de año escolar}
\label{\detokenize{informacion:actividades-inicio-de-ano-escolar}}\begin{itemize}
\item {} 
\sphinxAtStartPar
 ASIGNAR equipos de pareja (hombre y mujer)

\item {} 
\sphinxAtStartPar
 DAR hoja para que anoten sus nombres, correo institucional y equipo asignado por pareja.

\item {} 
\sphinxAtStartPar
 PRESENTACIÓN del docente.

\item {} 
\sphinxAtStartPar
 Explicación \sphinxstylestrong{útiles escolares}.

\item {} 
\sphinxAtStartPar
 Explicación de \sphinxstylestrong{porcentaje de notas \%}.

\item {} 
\sphinxAtStartPar
 Explicación \sphinxstylestrong{equivalencia de notas INSTRUIMOS}.

\item {} 
\sphinxAtStartPar
 Explicación \sphinxstylestrong{sello docente}.

\item {} 
\sphinxAtStartPar
 Explicación \sphinxstylestrong{acceso al portal aula 115 (Internet e Intranet)}.

\item {} 
\sphinxAtStartPar
 Explicación \sphinxstylestrong{modelo de almacenamiento}.

\item {} 
\sphinxAtStartPar
 Explicación \sphinxstylestrong{acceso al aula 115 en contrajornada}.

\item {} 
\sphinxAtStartPar
 \sphinxstylestrong{Reglas del aula} (cada estudiante lee una regla y escriben las reglas en su cuaderno).

\item {} 
\sphinxAtStartPar
 FIRMA DE COMPROMISOS en el cuaderno (Desbloqueo de equipo).

\item {} 
\sphinxAtStartPar
 Acceso de estudiantes al PC.

\item {} 
\sphinxAtStartPar
 Explicación estructura carpeta de evidencias.

\end{itemize}

\sphinxstepscope


\chapter{Reglas del aula}
\label{\detokenize{reglas:reglas-del-aula}}\label{\detokenize{reglas::doc}}\begin{enumerate}
\sphinxsetlistlabels{\arabic}{enumi}{enumii}{}{.}%
\item {} 
\sphinxAtStartPar
Las personas que hagan uso del aula 115, deben atender y obedecer las instrucciones del profesor (o encargado del aula), así como mantener respeto, buen comportamiento y disciplina.

\end{enumerate}
\begin{enumerate}
\sphinxsetlistlabels{\arabic}{enumi}{enumii}{}{.}%
\setcounter{enumi}{1}
\item {} 
\sphinxAtStartPar
Los estudiantes deben acomodarse de la siguiente manera:
\begin{itemize}
\item {} 
\sphinxAtStartPar
Bolsos en la parte de atrás de la silla.

\item {} 
\sphinxAtStartPar
Cuadernos y útiles en el soporte frontal.

\item {} 
\sphinxAtStartPar
Si hay evaluación, los bolsos deben estar en el tablero.

\end{itemize}

\end{enumerate}
\begin{enumerate}
\sphinxsetlistlabels{\arabic}{enumi}{enumii}{}{.}%
\setcounter{enumi}{2}
\item {} 
\sphinxAtStartPar
Para el aseo al final de la jornada escolar:
\begin{itemize}
\item {} 
\sphinxAtStartPar
Primero se dejan las sillas organizadas en la esquina del aula.

\item {} 
\sphinxAtStartPar
Posteriormente se procede a barrer y luego trapear.

\item {} 
\sphinxAtStartPar
Se debe botar la basura y el agua sucia.

\item {} 
\sphinxAtStartPar
También se debe reciclar las botellas plásticas.

\end{itemize}

\end{enumerate}
\begin{enumerate}
\sphinxsetlistlabels{\arabic}{enumi}{enumii}{}{.}%
\setcounter{enumi}{3}
\item {} 
\sphinxAtStartPar
Cualquier duda sobre el uso del aula, preguntar al docente o al encargado.

\end{enumerate}
\begin{enumerate}
\sphinxsetlistlabels{\arabic}{enumi}{enumii}{}{.}%
\setcounter{enumi}{4}
\item {} 
\sphinxAtStartPar
Cumplimiento de la hora establecida:
\begin{itemize}
\item {} 
\sphinxAtStartPar
De entrada los estudiantes deben formar dos filas. Una fila de hombres y una fila de mujeres.

\item {} 
\sphinxAtStartPar
Para permitir la entrada los estudiantes deben cumplir:
\begin{itemize}
\item {} 
\sphinxAtStartPar
Portar bien el uniforme.

\item {} 
\sphinxAtStartPar
Celulares guardados en bolsos y en silencio.

\item {} 
\sphinxAtStartPar
No tener ni estar consumiendo comida o bebida alguna.

\end{itemize}

\end{itemize}

\end{enumerate}
\begin{enumerate}
\sphinxsetlistlabels{\arabic}{enumi}{enumii}{}{.}%
\setcounter{enumi}{5}
\item {} 
\sphinxAtStartPar
Llegada tarde:
\begin{itemize}
\item {} 
\sphinxAtStartPar
Una vez cumplida la hora de entrada, se cierra la puerta y empieza la clase.
\begin{itemize}
\item {} 
\sphinxAtStartPar
Se hace llamado a lista.

\item {} 
\sphinxAtStartPar
Los estudiantes que lleguen tarde permanecen afuera sin distraer, hacer ruido o desorden y se les coloca la falla por tardanza.

\item {} 
\sphinxAtStartPar
Una vez terminado el llamado a lista, los estudiantes faltantes ingresan en silencio al aula. (Profesores dan 5 min para hacer aseo antes de entrar a la siguiente clase).

\end{itemize}

\end{itemize}

\end{enumerate}
\begin{enumerate}
\sphinxsetlistlabels{\arabic}{enumi}{enumii}{}{.}%
\setcounter{enumi}{6}
\item {} 
\sphinxAtStartPar
El aula de informática es exclusivamente para \sphinxstylestrong{realizar trabajos académicos} (que necesiten de un computador). Los trabajos deben estar relacionados con la Institución Educativa John F. Kennedy. Está prohibido hacer trabajos de cualquier otro tipo o institución ajenos.

\end{enumerate}
\begin{enumerate}
\sphinxsetlistlabels{\arabic}{enumi}{enumii}{}{.}%
\setcounter{enumi}{7}
\item {} 
\sphinxAtStartPar
Si el estudiante es identificado \sphinxstylestrong{realizando cualquier actividad ajena a lo establecido por el docente} o incumpliendo cualquiera de las presentes reglas, su equipo será bloqueado.

\end{enumerate}


\bigskip\hrule\bigskip



\section{PRØHIBICIONES}
\label{\detokenize{reglas:prohibiciones}}\begin{itemize}
\item {} 
\sphinxAtStartPar
Pedir permiso para ir al baño (ORDEN DE COORDINACIÓN) (PROBLEMAS URINARIOS PREVIA EVIDENCIA MÉDICA) \sphinxstylestrong{(Pedir permiso constantemente genera mucha interrupción en el desarrollo de la clase)}

\item {} 
\sphinxAtStartPar
Consumir o ingresar cualquier tipo de \sphinxstylestrong{comida} o \sphinxstylestrong{bebida} en el aula de informática. Los residuos de comida o bebida caen sobre los equipos.

\item {} 
\sphinxAtStartPar
Dejar bolsos o cualquier otro elemento encima de las mesas de los equipos. (\sphinxstyleemphasis{Los cables de alimentación se pueden desconectar y esto daña los equipos}).

\item {} 
\sphinxAtStartPar
\sphinxstylestrong{Acostarse} en el piso del aula.

\item {} 
\sphinxAtStartPar
\sphinxstylestrong{Sentarse de forma incorrecta en las sillas} del aula o \sphinxstylestrong{sentarse en las mesas}.

\item {} 
\sphinxAtStartPar
\sphinxstylestrong{Gritos extemporáneos} o hablar en \sphinxstylestrong{voz alta}.

\item {} 
\sphinxAtStartPar
\sphinxstylestrong{Dañar o escribir} en los muebles del aula. (lápiz, bolígrafos, pinturas, etc.)

\item {} 
\sphinxAtStartPar
\sphinxstylestrong{Maquillarse} en el aula.

\item {} 
\sphinxAtStartPar
Realizar cualquier tipo de bailes o \sphinxstylestrong{expresiones corporales} en el aula.

\item {} 
\sphinxAtStartPar
\sphinxstylestrong{Consumir} cualquier tipo de droga o estupefaciente.

\item {} 
\sphinxAtStartPar
\sphinxstylestrong{Arrojar basura} en el aula.

\item {} 
\sphinxAtStartPar
\sphinxstylestrong{No hacer el aseo} o hacerlo de forma mediocre.

\item {} 
\sphinxAtStartPar
Esconder y robar cualquier objeto del aula o de los compañeros.

\item {} 
\sphinxAtStartPar
Ocupar el puesto del compañero.

\item {} 
\sphinxAtStartPar
Realizar trabajos diferentes a la asignatura en el aula.

\item {} 
\sphinxAtStartPar
Quitarse los zapatos.

\item {} 
\sphinxAtStartPar
Uso inapropiado del aula:
\begin{itemize}
\item {} 
\sphinxAtStartPar
Invadir el espacio del docente.

\item {} 
\sphinxAtStartPar
Bajar los tacos.

\item {} 
\sphinxAtStartPar
Escuchar música en el computador o celular sin autorización del docente.

\item {} 
\sphinxAtStartPar
Ver videos en el computador o celular sin autorización del docente.

\item {} 
\sphinxAtStartPar
Jugar videojuegos en el computador o celular sin autorización del docente.

\item {} 
\sphinxAtStartPar
Manipular, portar o \sphinxstylestrong{dejar visible el celular} sin previa autorización del docente (jugar, redes sociales, etc.)

\item {} 
\sphinxAtStartPar
Tomar fotos, selfis o videos sin autorización del docente.

\item {} 
\sphinxAtStartPar
Manipular aires acondicionados, ventiladores, televisor, tablero.

\item {} 
\sphinxAtStartPar
Acceder al gabinete, armario y escritorio del docente.

\item {} 
\sphinxAtStartPar
Desarmar o dar golpes a los equipos.

\item {} 
\sphinxAtStartPar
Instalar software en los equipos sin previa autorización del profesor.

\item {} 
\sphinxAtStartPar
Borrar archivos de los equipos.

\end{itemize}

\end{itemize}


\bigskip\hrule\bigskip



\section{¿Qué pasa si se incumplen las reglas?}
\label{\detokenize{reglas:que-pasa-si-se-incumplen-las-reglas}}
\sphinxAtStartPar
Si el estudiante incumple alguna de las reglas, se procede con las siguientes \sphinxstylestrong{acciones formativas}:
\begin{itemize}
\item {} 
\sphinxAtStartPar
Disminución de valor de la nota en \sphinxstylestrong{perfil Kennediano}.

\item {} 
\sphinxAtStartPar
Disminución de valor de la nota \sphinxstylestrong{trabajo en equipo}.

\item {} 
\sphinxAtStartPar
Posible cambio de puesto de trabajo.

\end{itemize}

\sphinxAtStartPar
Mientras que el estudiante NO CUMPLA con las \sphinxstylestrong{acciones formativas}, se procede a:
\begin{itemize}
\item {} 
\sphinxAtStartPar
No calificar o evaluar cualquiera de las actividades académicas en proceso.

\item {} 
\sphinxAtStartPar
No firmar paz y salvo.

\end{itemize}


\bigskip\hrule\bigskip



\section{¿Qué pasa si el incumplimiento es reiterativo?}
\label{\detokenize{reglas:que-pasa-si-el-incumplimiento-es-reiterativo}}

\begin{savenotes}\sphinxattablestart
\sphinxthistablewithglobalstyle
\centering
\begin{tabulary}{\linewidth}[t]{T}
\sphinxtoprule
\sphinxtableatstartofbodyhook

\begin{sphinxfigure-in-table}
\centering
\capstart
\noindent\sphinxincludegraphics[width=300\sphinxpxdimen]{{image026}.png}
\sphinxfigcaption{Sin excusas}\label{\detokenize{reglas:id1}}\end{sphinxfigure-in-table}\relax
\\
\sphinxbottomrule
\end{tabulary}
\sphinxtableafterendhook\par
\sphinxattableend\end{savenotes}

\sphinxAtStartPar
En caso que el estudiante incumpla alguna de las reglas continuamente, se iniciará \sphinxstylestrong{proceso disciplinario} citando al acudiente, representante de curso y director de curso.

\sphinxstepscope


\chapter{Lógica computacional 8}
\label{\detokenize{cursos/logica8/index:logica-computacional-8}}\label{\detokenize{cursos/logica8/index::doc}}
\sphinxAtStartPar
Bienvenido al curso lógica computacional de grado octavo.

\sphinxAtStartPar
Acceda a las guías didácticas según el periodo.

\sphinxstepscope


\section{Periodo 01}
\label{\detokenize{cursos/logica8/periodo01/index:periodo-01}}\label{\detokenize{cursos/logica8/periodo01/index::doc}}\begin{quote}\begin{description}
\sphinxlineitem{DESEMPEÑO}
\sphinxAtStartPar
Explica los conceptos de hardware, software, sistemas computacionales, informática y los relaciona con la definición de \sphinxstylestrong{lógica computacional}.

\end{description}\end{quote}


\bigskip\hrule\bigskip



\subsection{Contenidos}
\label{\detokenize{cursos/logica8/periodo01/index:contenidos}}
\sphinxstepscope


\subsubsection{Guía 01: Lógica computacional}
\label{\detokenize{cursos/logica8/periodo01/G01:guia-01-logica-computacional}}\label{\detokenize{cursos/logica8/periodo01/G01::doc}}

\begin{savenotes}\sphinxattablestart
\sphinxthistablewithglobalstyle
\centering
\begin{tabular}[t]{\X{35}{100}\X{65}{100}}
\sphinxtoprule
\sphinxtableatstartofbodyhook

\begin{savenotes}\sphinxattablestart
\sphinxthistablewithglobalstyle
\centering
\begin{tabulary}{\linewidth}[t]{T}
\sphinxtoprule
\sphinxtableatstartofbodyhook
\sphinxAtStartPar
\sphinxincludegraphics{{escudo}.png}
\\
\sphinxbottomrule
\end{tabulary}
\sphinxtableafterendhook\par
\sphinxattableend\end{savenotes}
&

\begin{savenotes}\sphinxattablestart
\sphinxthistablewithglobalstyle
\centering
\begin{tabulary}{\linewidth}[t]{T}
\sphinxtoprule
\sphinxtableatstartofbodyhook
\sphinxAtStartPar
\sphinxstylestrong{INSTITUCIÓN EDUCATIVA}: John F. Kennedy
\\
\sphinxhline
\sphinxAtStartPar
\sphinxstylestrong{DOCENTE}: Julian Camilo Fonseca Romero
\\
\sphinxhline
\sphinxAtStartPar
\sphinxstylestrong{ASIGNATURA}: Lógica computacional
\\
\sphinxhline
\sphinxAtStartPar
\sphinxstylestrong{GRADO}: OCTAVO
\\
\sphinxhline
\sphinxAtStartPar
\sphinxstylestrong{TIEMPO}: 10 semanas (1 hora / semana)
\\
\sphinxhline
\sphinxAtStartPar
\sphinxstylestrong{TIEMPO DE TRABAJO}: 10/10 horas
\\
\sphinxbottomrule
\end{tabulary}
\sphinxtableafterendhook\par
\sphinxattableend\end{savenotes}
\\
\sphinxbottomrule
\end{tabular}
\sphinxtableafterendhook\par
\sphinxattableend\end{savenotes}

\begin{sphinxuseclass}{sd-tab-set}
\begin{sphinxuseclass}{sd-tab-item}\subsubsection*{Objetivo}

\begin{sphinxuseclass}{sd-tab-content}\begin{itemize}
\item {} 
\sphinxAtStartPar
\sphinxstylestrong{Temática}: Conceptos de lógica computacional.

\item {} 
\sphinxAtStartPar
\sphinxstylestrong{Objetivo de aprendizaje}: Comprender la interacción entre hardware, software y sistemas computacionales, aplicando el razonamiento lógico para entender el concepto de lógica computacional.

\end{itemize}

\end{sphinxuseclass}
\end{sphinxuseclass}
\begin{sphinxuseclass}{sd-tab-item}\subsubsection*{Orientaciones curriculares}

\begin{sphinxuseclass}{sd-tab-content}\begin{itemize}
\item {} 
\sphinxAtStartPar
\sphinxstylestrong{Componente}: Naturaleza y Evolución de la T\&I.

\item {} 
\sphinxAtStartPar
\sphinxstylestrong{Competencia}: Relaciono saberes, conocimientos tecnológicos e informáticos con los conocimientos de otras disciplinas.

\item {} 
\sphinxAtStartPar
\sphinxstylestrong{Evidencia de aprendizaje}: Comprendo los principios y conocimientos tecnológicos e informáticos que hacen posible el funcionamiento de productos tecnológicos actuales.

\end{itemize}

\end{sphinxuseclass}
\end{sphinxuseclass}
\begin{sphinxuseclass}{sd-tab-item}\subsubsection*{Criterios de evaluación}

\begin{sphinxuseclass}{sd-tab-content}\begin{itemize}
\item {} 
\sphinxAtStartPar
Actividades desarrolladas en su totalidad.

\item {} 
\sphinxAtStartPar
Participación en clase.

\item {} 
\sphinxAtStartPar
Explicación clara de conceptos.

\item {} 
\sphinxAtStartPar
Argumentación de ideas.

\item {} 
\sphinxAtStartPar
Cumplimiento de las reglas del aula.

\end{itemize}

\end{sphinxuseclass}
\end{sphinxuseclass}
\begin{sphinxuseclass}{sd-tab-item}\subsubsection*{Recursos (1)}

\begin{sphinxuseclass}{sd-tab-content}\begin{itemize}
\item {} 
\sphinxAtStartPar
{\hyperref[\detokenize{cursos/logica8/recursos/T01::doc}]{\sphinxcrossref{\DUrole{std}{\DUrole{std-doc}{TEXTO: Concepto lógica computacional}}}}}

\end{itemize}

\end{sphinxuseclass}
\end{sphinxuseclass}
\end{sphinxuseclass}

\paragraph{Mensaje del autor}
\label{\detokenize{cursos/logica8/periodo01/G01:mensaje-del-autor}}\begin{quote}
\begin{quote}

\sphinxAtStartPar
“En esta guía, el estudiante comprenderá los conceptos de hardware y software. Asimismo, argumentará la importancia del uso de la CPU y la memoria RAM en el procesamiento de la información. Además, analizará el concepto de sistema computacional, el cual forma parte de la definición de informática. Finalmente, integrará los conceptos de razonamiento y lógica para entender la definición de \sphinxstylestrong{lógica computacional}.”
\end{quote}

\begin{flushright}
---Camilo Fonseca
\end{flushright}
\end{quote}


\bigskip\hrule\bigskip



\paragraph{A00: Introducción (+0,1)}
\label{\detokenize{cursos/logica8/periodo01/G01:a00-introduccion-0-1}}
\sphinxAtStartPar
\DUrole{sd-sphinx-override}{\DUrole{sd-badge}{\DUrole{sd-bg-light}{\DUrole{sd-bg-text-light}{Exploración}}}}
\begin{enumerate}
\sphinxsetlistlabels{\arabic}{enumi}{enumii}{}{.}%
\item {} 
\sphinxAtStartPar
TRANSCRIBA en su cuaderno el objetivo, las orientaciones curriculares y los criterios de evaluación de la presente guía.

\item {} 
\sphinxAtStartPar
LEA el \sphinxstylestrong{mensaje del autor} y en su cuaderno responda:
\begin{itemize}
\item {} 
\sphinxAtStartPar
¿Qué tienen en común un smartphone, un smartTV y un computador?

\end{itemize}

\end{enumerate}


\bigskip\hrule\bigskip



\paragraph{A01: Concepto de hardware (+0,1)}
\label{\detokenize{cursos/logica8/periodo01/G01:a01-concepto-de-hardware-0-1}}
\sphinxAtStartPar
\DUrole{sd-sphinx-override}{\DUrole{sd-badge}{\DUrole{sd-bg-info}{\DUrole{sd-bg-text-info}{Estructuración}}}}
\begin{enumerate}
\sphinxsetlistlabels{\arabic}{enumi}{enumii}{}{.}%
\item {} 
\sphinxAtStartPar
En su cuaderno RESPONDA las siguientes preguntas:
\begin{itemize}
\item {} 
\sphinxAtStartPar
¿Qué es el hardware?

\item {} 
\sphinxAtStartPar
¿De que se compone el hardware de un sistema informático?

\item {} 
\sphinxAtStartPar
¿El concepto de hardware solo puede aplicarse a computadores? ¿Por qué?

\end{itemize}

\item {} 
\sphinxAtStartPar
En su cuaderno CREE un \sphinxstylestrong{mapa conceptual} que muestre los diferentes componentes del hardware y dos ejemplos por cada componente.

\end{enumerate}


\bigskip\hrule\bigskip



\paragraph{A02: La CPU (+0,1)}
\label{\detokenize{cursos/logica8/periodo01/G01:a02-la-cpu-0-1}}
\sphinxAtStartPar
\DUrole{sd-sphinx-override}{\DUrole{sd-badge}{\DUrole{sd-bg-info}{\DUrole{sd-bg-text-info}{Estructuración}}}}
\begin{enumerate}
\sphinxsetlistlabels{\arabic}{enumi}{enumii}{}{.}%
\item {} 
\sphinxAtStartPar
En su cuaderno RESPONDA las siguientes preguntas:
\begin{itemize}
\item {} 
\sphinxAtStartPar
¿Qué es la CPU?

\item {} 
\sphinxAtStartPar
¿Por qué la CPU es considerada el cerebro del computador?

\item {} 
\sphinxAtStartPar
¿En que se mide la velocidad de las CPUs?

\item {} 
\sphinxAtStartPar
¿Qué significa que una velocidad en la CPU sea alta o baja?

\end{itemize}

\item {} 
\sphinxAtStartPar
En su cuaderno CREE un \sphinxstylestrong{mapa conceptual} de los tres principales componentes de la CPU y descríbalos.

\item {} 
\sphinxAtStartPar
En INTERNET, CONSULTE 3 marcas de fabricantes de CPU y descríbalas en su cuaderno.

\end{enumerate}


\bigskip\hrule\bigskip



\paragraph{A03: Memoria RAM (+0,1)}
\label{\detokenize{cursos/logica8/periodo01/G01:a03-memoria-ram-0-1}}
\sphinxAtStartPar
\DUrole{sd-sphinx-override}{\DUrole{sd-badge}{\DUrole{sd-bg-info}{\DUrole{sd-bg-text-info}{Estructuración}}}}
\begin{enumerate}
\sphinxsetlistlabels{\arabic}{enumi}{enumii}{}{.}%
\item {} 
\sphinxAtStartPar
Acorde a la definición \sphinxstylestrong{CONCEPTO DE MEMORIA RAM}, en su cuaderno responda las siguientes preguntas:
\begin{itemize}
\item {} 
\sphinxAtStartPar
¿Qué es la memoria RAM?

\item {} 
\sphinxAtStartPar
¿De que se encarga la memoria RAM?

\item {} 
\sphinxAtStartPar
¿Por qué la memoria RAM es tan importante al momento de procesar la información?

\item {} 
\sphinxAtStartPar
¿Cuales son las características de la memoria RAM?

\end{itemize}

\item {} 
\sphinxAtStartPar
En Internet, CONSULTE las diferentes generaciones que han tenido las memorias RAM y con esta información en su cuaderno CREE un mapa conceptual.

\end{enumerate}


\bigskip\hrule\bigskip



\paragraph{A04: CPU \& RAM (+0,5)}
\label{\detokenize{cursos/logica8/periodo01/G01:a04-cpu-ram-0-5}}
\sphinxAtStartPar
\DUrole{sd-sphinx-override}{\DUrole{sd-badge}{\DUrole{sd-bg-danger}{\DUrole{sd-bg-text-danger}{Transferencia}}}}
\begin{enumerate}
\sphinxsetlistlabels{\arabic}{enumi}{enumii}{}{.}%
\item {} 
\sphinxAtStartPar
Con los temas vistos hasta ahora, en \sphinxstylestrong{UNA PÁGINA DE CUADERNO}, con sus palabras (no usar IA) RESPONDA la siguiente pregunta:
\begin{itemize}
\item {} 
\sphinxAtStartPar
¿Por qué la memoria RAM y la CPU son tan importantes al momento de procesar información?

\end{itemize}

\end{enumerate}


\bigskip\hrule\bigskip



\paragraph{A05: Software (+0,2)}
\label{\detokenize{cursos/logica8/periodo01/G01:a05-software-0-2}}
\sphinxAtStartPar
\DUrole{sd-sphinx-override}{\DUrole{sd-badge}{\DUrole{sd-bg-info}{\DUrole{sd-bg-text-info}{Estructuración}}}}
\begin{enumerate}
\sphinxsetlistlabels{\arabic}{enumi}{enumii}{}{.}%
\item {} 
\sphinxAtStartPar
En su cuaderno RESPONDA las siguientes preguntas:
\begin{itemize}
\item {} 
\sphinxAtStartPar
¿Qué es el software?

\item {} 
\sphinxAtStartPar
¿En qué se diferencian el HARDWARE y el SOFTWARE?

\item {} 
\sphinxAtStartPar
¿Por qué los controladores de HARDWARE no necesitan interfaz gráfica?

\end{itemize}

\item {} 
\sphinxAtStartPar
En su cuaderno, CREE un mapa conceptual que presente la \sphinxstylestrong{clasificación del software}.

\item {} 
\sphinxAtStartPar
CONSULTE en Internet y luego en su cuaderno REALICE una descripción de lo siguiente:
\begin{itemize}
\item {} 
\sphinxAtStartPar
10 sistemas operativos de la empresa MICROSOFT.

\item {} 
\sphinxAtStartPar
1 sistema operativo de la empresa GOOGLE.

\item {} 
\sphinxAtStartPar
3 sistemas operativos de la empresa APPLE.

\end{itemize}

\item {} 
\sphinxAtStartPar
En Internet, CONSULTE dos ejemplos para cada una de las siguientes aplicaciones, luego descríbalas en su cuaderno:
\begin{itemize}
\item {} 
\sphinxAtStartPar
Procesadores de texto

\item {} 
\sphinxAtStartPar
Navegadores web

\item {} 
\sphinxAtStartPar
Videojuegos

\item {} 
\sphinxAtStartPar
Editores de imágenes

\item {} 
\sphinxAtStartPar
Modelado 3D

\end{itemize}

\end{enumerate}


\bigskip\hrule\bigskip



\paragraph{A06: Sistema computacional (+0,1)}
\label{\detokenize{cursos/logica8/periodo01/G01:a06-sistema-computacional-0-1}}
\sphinxAtStartPar
\DUrole{sd-sphinx-override}{\DUrole{sd-badge}{\DUrole{sd-bg-info}{\DUrole{sd-bg-text-info}{Estructuración}}}}
\begin{enumerate}
\sphinxsetlistlabels{\arabic}{enumi}{enumii}{}{.}%
\item {} 
\sphinxAtStartPar
En su cuaderno RESPONDA las siguientes preguntas:
\begin{itemize}
\item {} 
\sphinxAtStartPar
¿Qué es un sistema computacional?

\item {} 
\sphinxAtStartPar
¿De qué manera un hardware interactúa con un software? (piense en un smartphone, un smart\sphinxhyphen{}TV, etc.)

\item {} 
\sphinxAtStartPar
¿Cómo un sistema computacional interactúa con la información? (piense en un software como Microsoft EXCEL o WHATSAPP)

\end{itemize}

\end{enumerate}


\bigskip\hrule\bigskip



\paragraph{A07: Ejemplos sistemas computacionales (+0,5)}
\label{\detokenize{cursos/logica8/periodo01/G01:a07-ejemplos-sistemas-computacionales-0-5}}
\sphinxAtStartPar
\DUrole{sd-sphinx-override}{\DUrole{sd-badge}{\DUrole{sd-bg-danger}{\DUrole{sd-bg-text-danger}{Transferencia}}}}
\begin{enumerate}
\sphinxsetlistlabels{\arabic}{enumi}{enumii}{}{.}%
\item {} 
\sphinxAtStartPar
En su cuaderno, con sus palabras ESCRIBA 5 ejemplos de sistemas computacionales (no usar IA). Tenga en cuenta el siguiente ejemplo:
\begin{quote}

\sphinxAtStartPar
\sphinxstylestrong{EJEMPLO 01}: Un \sphinxstyleemphasis{portátil} es un ejemplo de \sphinxstylestrong{sistema computacional} porque su \sphinxstyleemphasis{hardware} (teclado, pantalla, touchpad) y \sphinxstyleemphasis{software} (OS, aplicaciones de escritorio) permite registrar, almacenar y mostrar información para un usuario”.
\end{quote}

\end{enumerate}


\bigskip\hrule\bigskip



\paragraph{A08: Informática (+0,1)}
\label{\detokenize{cursos/logica8/periodo01/G01:a08-informatica-0-1}}
\sphinxAtStartPar
\DUrole{sd-sphinx-override}{\DUrole{sd-badge}{\DUrole{sd-bg-info}{\DUrole{sd-bg-text-info}{Estructuración}}}}
\begin{enumerate}
\sphinxsetlistlabels{\arabic}{enumi}{enumii}{}{.}%
\item {} 
\sphinxAtStartPar
RESPONDA las siguientes preguntas en su cuaderno:
\begin{itemize}
\item {} 
\sphinxAtStartPar
¿En que se relacionan la informática y la información?

\item {} 
\sphinxAtStartPar
Aparte de la \sphinxstylestrong{integridad} y \sphinxstylestrong{disponibilidad}, ¿Qué otras reglas pueden ser aplicadas en la información?

\end{itemize}

\item {} 
\sphinxAtStartPar
En su cuaderno, ESCRIBA los siguientes casos y por cada uno indique si aplica la regla 01 (\sphinxstyleemphasis{Integridad de la información}) o aplica la regla 02 (\sphinxstyleemphasis{disponibilidad de la información}) (escribir el por qué)
\begin{itemize}
\item {} 
\sphinxAtStartPar
\sphinxstylestrong{CASO 01}: \sphinxstyleemphasis{“De nada sirve que un banco entregue el valor de un saldo 10 meses después de ser solicitado. Es por eso que se usan las computadoras, porque permiten que la información este disponible y llegue oportunamente (milésimas de segundo)”}.

\item {} 
\sphinxAtStartPar
\sphinxstylestrong{CASO 02:} \sphinxstyleemphasis{“De nada sirve que una persona vaya al banco a consultar su saldo si este no corresponde al saldo real. Es por eso que se usan las computadoras, porque permiten que la información conserve su estado sin los riesgos que representa tenerla en otros medios como el papel o en la memoria de las personas”}.

\end{itemize}

\end{enumerate}


\begin{savenotes}\sphinxattablestart
\sphinxthistablewithglobalstyle
\centering
\begin{tabulary}{\linewidth}[t]{T}
\sphinxtoprule
\sphinxtableatstartofbodyhook

\begin{sphinxfigure-in-table}
\centering
\capstart
\noindent\sphinxincludegraphics[width=300\sphinxpxdimen]{{image011}.png}
\sphinxfigcaption{\sphinxstylestrong{BANCOS:} Los bancos deben garantizar integridad y disponibilidad}\label{\detokenize{cursos/logica8/periodo01/G01:id1}}\end{sphinxfigure-in-table}\relax
\\
\sphinxbottomrule
\end{tabulary}
\sphinxtableafterendhook\par
\sphinxattableend\end{savenotes}


\bigskip\hrule\bigskip



\paragraph{A09: Razonamiento y lógica (+0,1)}
\label{\detokenize{cursos/logica8/periodo01/G01:a09-razonamiento-y-logica-0-1}}
\sphinxAtStartPar
\DUrole{sd-sphinx-override}{\DUrole{sd-badge}{\DUrole{sd-bg-info}{\DUrole{sd-bg-text-info}{Estructuración}}}}
\begin{enumerate}
\sphinxsetlistlabels{\arabic}{enumi}{enumii}{}{.}%
\item {} 
\sphinxAtStartPar
En su cuaderno RESPONDA las siguientes preguntas:
\begin{itemize}
\item {} 
\sphinxAtStartPar
¿Qué es razonar?

\item {} 
\sphinxAtStartPar
¿Qué es lógica?

\item {} 
\sphinxAtStartPar
¿Cómo se relacionan los conceptos de razonar y lógica?

\item {} 
\sphinxAtStartPar
¿Se podrían comprobar resultados de razonamiento sin aplicar lógica? (explique su respuesta)

\end{itemize}

\item {} 
\sphinxAtStartPar
RAZONE la siguiente situación y en su cuaderno DEDUZCA que pudo haber pasado (3 conclusiones con temáticas diferentes):
\begin{itemize}
\item {} 
\sphinxAtStartPar
Premisa 01: Una persona escuchó ruidos extraños en su casa.

\item {} 
\sphinxAtStartPar
Premisa 02: Los ruidos fueron en horas de la noche.

\item {} 
\sphinxAtStartPar
Premisa 03: Habían varios residuos de comida en el piso.

\item {} 
\sphinxAtStartPar
Premisa 04: El computador estaba desconectado.

\item {} 
\sphinxAtStartPar
Premisa 05: No hay ningún objeto perdido.

\item {} 
\sphinxAtStartPar
\sphinxstylestrong{Conclusión 01:} ?

\item {} 
\sphinxAtStartPar
\sphinxstylestrong{Conclusión 02:} ?

\item {} 
\sphinxAtStartPar
\sphinxstylestrong{Conclusión 03:} ?

\end{itemize}

\end{enumerate}


\bigskip\hrule\bigskip



\paragraph{A10: Lógica computacional (+0,1)}
\label{\detokenize{cursos/logica8/periodo01/G01:a10-logica-computacional-0-1}}
\sphinxAtStartPar
\DUrole{sd-sphinx-override}{\DUrole{sd-badge}{\DUrole{sd-bg-info}{\DUrole{sd-bg-text-info}{Estructuración}}}}
\begin{enumerate}
\sphinxsetlistlabels{\arabic}{enumi}{enumii}{}{.}%
\item {} 
\sphinxAtStartPar
En su cuaderno RESPONDA las siguientes preguntas:
\begin{itemize}
\item {} 
\sphinxAtStartPar
¿Qué es un algoritmo?

\item {} 
\sphinxAtStartPar
¿Qué es lógica computacional?

\item {} 
\sphinxAtStartPar
¿Cómo se relaciona la lógica computacional y la informática?

\item {} 
\sphinxAtStartPar
¿Cómo se relacionan los algoritmos y la lógica computacional?

\end{itemize}

\item {} 
\sphinxAtStartPar
En su cuaderno ESCRIBA las siguientes frases e indique si es FALSO, VERDADERO y explique el por qué:
\begin{itemize}
\item {} 
\sphinxAtStartPar
“\sphinxstyleemphasis{La \sphinxstylestrong{lógica computacional} establece las reglas para programar una persona y modificar su comportamiento}”. \_\_

\item {} 
\sphinxAtStartPar
“\sphinxstyleemphasis{Optimizar \sphinxstylestrong{algoritmos} permite solucionar problemas más rápido pero con los mismos recursos}”. \_\_

\end{itemize}

\end{enumerate}


\bigskip\hrule\bigskip



\paragraph{A11: Examen final (+3,0)}
\label{\detokenize{cursos/logica8/periodo01/G01:a11-examen-final-3-0}}
\sphinxAtStartPar
\DUrole{sd-sphinx-override}{\DUrole{sd-badge}{\DUrole{sd-bg-danger}{\DUrole{sd-bg-text-danger}{Transferencia}}}}
\begin{quote}\begin{description}
\sphinxlineitem{TIEMPO}
\sphinxAtStartPar
30 minutos

\sphinxlineitem{PREGUNTAS}
\sphinxAtStartPar
10

\end{description}\end{quote}

\begin{sphinxadmonition}{note}{Posibles preguntas}
\begin{enumerate}
\sphinxsetlistlabels{\arabic}{enumi}{enumii}{}{.}%
\item {} 
\sphinxAtStartPar
Explique con sus propias palabras qué es el \sphinxstylestrong{hardware} y mencione al menos tres categorías en las que se clasifican sus componentes, dando un ejemplo de cada una.

\item {} 
\sphinxAtStartPar
¿Por qué la \sphinxstylestrong{CPU} es considerada el “cerebro” del computador? Describa sus tres componentes principales y su función.

\item {} 
\sphinxAtStartPar
Escriba las funciones principales de la \sphinxstylestrong{memoria RAM} y explique por qué es fundamental en el procesamiento de información.

\item {} 
\sphinxAtStartPar
Compare y contraste el \sphinxstylestrong{hardware} y el \sphinxstylestrong{software}, resaltando sus diferencias fundamentales.

\item {} 
\sphinxAtStartPar
¿Qué es el \sphinxstylestrong{software} y cómo se clasifica? Describa brevemente cada una de sus categorías principales.

\item {} 
\sphinxAtStartPar
Describa el concepto de \sphinxstylestrong{“interfaz gráfica”} y explique su importancia para los usuarios.

\item {} 
\sphinxAtStartPar
Defina qué es un \sphinxstylestrong{“sistema computacional”} y cómo interactúan sus componentes para procesar información.

\item {} 
\sphinxAtStartPar
¿Qué es la \sphinxstylestrong{informática} y cuáles son sus objetivos principales en relación con los sistemas computacionales?

\item {} 
\sphinxAtStartPar
Explique las dos \sphinxstylestrong{reglas fundamentales de la informática} (Integridad y Disponibilidad de la información) y por qué son importantes.

\item {} 
\sphinxAtStartPar
Describa el concepto de \sphinxstylestrong{“razonar”} y proporcione un ejemplo propio (diferente al ejemplo de sócrates), que incluya premisas y una conclusión.

\item {} 
\sphinxAtStartPar
¿Qué es la \sphinxstylestrong{lógica} y cómo se relaciona con el razonamiento? Explique su papel en la verificación de la validez de los argumentos.

\item {} 
\sphinxAtStartPar
¿Se podría comprobar la \sphinxstylestrong{validez de un razonamiento} sin aplicar principios de lógica? Justifica tu respuesta.

\item {} 
\sphinxAtStartPar
Defina qué es un \sphinxstylestrong{“algoritmo”} y cuáles son sus características principales.

\item {} 
\sphinxAtStartPar
¿Qué es la \sphinxstylestrong{“lógica computacional”} y cómo aplica los principios de la lógica y la informática?

\item {} 
\sphinxAtStartPar
¿Cómo se relacionan los \sphinxstylestrong{algoritmos} y la \sphinxstylestrong{lógica computacional} en la resolución de problemas?

\item {} 
\sphinxAtStartPar
Basándose en la información sobre \sphinxstylestrong{hardware} y \sphinxstylestrong{software}, ¿cómo cree que un smartphone o una tablet funcionan como sistemas computacionales?

\item {} 
\sphinxAtStartPar
¿Por qué es importante que la memoria \sphinxstylestrong{RAM} y la \sphinxstylestrong{CPU} trabajen de manera coordinada para un buen rendimiento del sistema?

\item {} 
\sphinxAtStartPar
Escriba un ejemplo de su vida diaria donde se aplique el concepto de \sphinxstylestrong{“disponibilidad de la información”}.

\item {} 
\sphinxAtStartPar
Explique un escenario donde la \sphinxstylestrong{“integridad de la información”} sea crucial y qué podría suceder si esta se viera comprometida.

\item {} 
\sphinxAtStartPar
¿Cómo la evolución del \sphinxstylestrong{hardware} y el \sphinxstylestrong{software} ha impactado la forma en que utilizamos los \sphinxstylestrong{sistemas computacionales} hoy en día?

\end{enumerate}
\end{sphinxadmonition}

\sphinxstepscope


\subsubsection{Plan de refuerzo 01}
\label{\detokenize{cursos/logica8/periodo01/R01:plan-de-refuerzo-01}}\label{\detokenize{cursos/logica8/periodo01/R01::doc}}
\begin{sphinxadmonition}{note}{EVALUACIÓN: Plan de refuerzo (10 minutos)}

\sphinxAtStartPar
Estudiar los temas a continuación. Se realiza examen de 5 preguntas. Si el estudiante no asiste en la fecha programada, este será evaluado la próxima vez que asista a clase.
\end{sphinxadmonition}


\bigskip\hrule\bigskip



\paragraph{Hardware}
\label{\detokenize{cursos/logica8/periodo01/R01:hardware}}
\sphinxAtStartPar
El hardware es el conjunto de componentes de un sistema informático o computacional que se puede ver o tocar. Por ejemplo:
\begin{itemize}
\item {} 
\sphinxAtStartPar
\sphinxstylestrong{Dispositivos de entrada:} Teclados, mouse, micrófonos.

\item {} 
\sphinxAtStartPar
\sphinxstylestrong{Dispositivos de salida:} Monitores, impresoras, bafles.

\item {} 
\sphinxAtStartPar
\sphinxstylestrong{Procesamiento:} CPU.

\item {} 
\sphinxAtStartPar
\sphinxstylestrong{Memoria:} memoria RAM, memoria CACHE.

\item {} 
\sphinxAtStartPar
\sphinxstylestrong{Almacenamiento:} Discos duros, memorias USB.

\end{itemize}


\bigskip\hrule\bigskip



\paragraph{Software}
\label{\detokenize{cursos/logica8/periodo01/R01:software}}
\sphinxAtStartPar
El SOFTWARE es el conjunto de \sphinxstylestrong{programas}, \sphinxstylestrong{bases de datos}, diferentes tipos de \sphinxstylestrong{archivos} y \sphinxstylestrong{documentación} que le dicen al HARDWARE qué hacer.

\sphinxAtStartPar
A diferencia del hardware que es físico y tangible, el software no se puede tocar. El software solo puede verse su funcionamiento en la pantalla.


\bigskip\hrule\bigskip



\paragraph{Sistema computacional}
\label{\detokenize{cursos/logica8/periodo01/R01:sistema-computacional}}
\sphinxAtStartPar
Es el conjunto de \sphinxstylestrong{hardware} que interactúa con un \sphinxstylestrong{software} para \sphinxstyleemphasis{almacenar}, \sphinxstyleemphasis{procesar} y \sphinxstyleemphasis{mostrar} información. El computador es solo un ejemplo de un sistema computacional.


\bigskip\hrule\bigskip



\paragraph{Informática}
\label{\detokenize{cursos/logica8/periodo01/R01:informatica}}
\sphinxAtStartPar
La informática es la \sphinxstyleemphasis{ciencia} o \sphinxstyleemphasis{disciplina} que estudia el diseño, construcción e implementación de sistemas computacionales (hardware que interactúa con un software para registrar, almacenar, procesar y mostrar la información) para garantizar la integridad y disponibilidad de la información.


\bigskip\hrule\bigskip



\paragraph{Razonar}
\label{\detokenize{cursos/logica8/periodo01/R01:razonar}}
\sphinxAtStartPar
Actividad mental que permite organizar las ideas para llegar a una conclusión (comprobar, deducir).


\bigskip\hrule\bigskip



\paragraph{Lógica}
\label{\detokenize{cursos/logica8/periodo01/R01:logica}}
\sphinxAtStartPar
La lógica es la disciplina que permite analizar y verificar la validez de los razonamientos, determinando si una conclusión se deriva correctamente de unas premisas.


\bigskip\hrule\bigskip



\paragraph{Lógica computacional}
\label{\detokenize{cursos/logica8/periodo01/R01:logica-computacional}}
\sphinxAtStartPar
Aplica los principios de la \sphinxstylestrong{lógica} en los \sphinxstylestrong{sistemas computacionales} para representar la \sphinxstylestrong{información}. La lógica computacional se encarga de:
\begin{itemize}
\item {} 
\sphinxAtStartPar
Establecer las reglas para programar un computador.

\item {} 
\sphinxAtStartPar
Diseñar algoritmos (serie de pasos que se ejecutan de forma ordenada para resolver un problema).

\item {} 
\sphinxAtStartPar
Optimizar los procesos computacionales (hacer más con menos recursos).

\end{itemize}

\sphinxstepscope


\section{Recursos}
\label{\detokenize{cursos/logica8/recursos/index:recursos}}\label{\detokenize{cursos/logica8/recursos/index::doc}}
\sphinxAtStartPar
En esta sección encuentra los recursos educativos necesarios para el desarrollo del curso.


\bigskip\hrule\bigskip



\subsection{Contenidos}
\label{\detokenize{cursos/logica8/recursos/index:contenidos}}
\sphinxstepscope


\subsubsection{T01: Concepto lógica computacional}
\label{\detokenize{cursos/logica8/recursos/T01:t01-concepto-logica-computacional}}\label{\detokenize{cursos/logica8/recursos/T01::doc}}
\sphinxAtStartPar
A continuación se presentan conceptos que ayudarán a entender más adelante el concepto de \sphinxstylestrong{lógica computacional}.

\begin{sphinxadmonition}{note}{RESUMEN}
\begin{equation*}
\begin{split} LogicaComputacional(LOGICA + INFORMATICA) \end{split}
\end{equation*}

\bigskip\hrule\bigskip

\begin{equation*}
\begin{split} LOGICA(Razon) \end{split}
\end{equation*}

\bigskip\hrule\bigskip

\begin{equation*}
\begin{split} INFORMATICA(SistemaComputacional(desarrollo) + INFO(Integridad + Disponibilidad))\end{split}
\end{equation*}\begin{equation*}
\begin{split} SistemaComputacional(HARDWARE + SOFTWARE + INFO(Almacenar + Procesar + Mostrar))\end{split}
\end{equation*}\end{sphinxadmonition}


\bigskip\hrule\bigskip



\paragraph{¿Qué es el HARDWARE?}
\label{\detokenize{cursos/logica8/recursos/T01:que-es-el-hardware}}
\sphinxAtStartPar
%
\begin{footnote}[1]\sphinxAtStartFootnote
Deitel, P. J., \& Deitel, H. (2014). C++: cómo programar (A. V. Romero Elizondo, Trans.). Pearson Educación.
%
\end{footnote} El \sphinxstylestrong{HARDWARE}, es el conjunto de componentes de un \sphinxstylestrong{sistema informático} (por ejemplo un computador) los cuales se pueden ver o tocar.

\sphinxAtStartPar
Estos componentes se clasifican en componentes de entrada, salida, procesamiento, memoria y almacenamiento.

\sphinxAtStartPar
A continuación se muestran algunos ejemplos:
\begin{itemize}
\item {} 
\sphinxAtStartPar
\sphinxstylestrong{Dispositivos de entrada:} Teclados, mouse, micrófonos

\item {} 
\sphinxAtStartPar
\sphinxstylestrong{Dispositivos de salida:} Monitores, impresoras, bafles

\item {} 
\sphinxAtStartPar
\sphinxstylestrong{Procesamiento:} CPU

\item {} 
\sphinxAtStartPar
\sphinxstylestrong{Memoria:} memoria RAM, memoria CACHE

\item {} 
\sphinxAtStartPar
\sphinxstylestrong{Almacenamiento:} Discos duros, memorias USB

\end{itemize}


\bigskip\hrule\bigskip



\subparagraph{¿Qué es una CPU?}
\label{\detokenize{cursos/logica8/recursos/T01:que-es-una-cpu}}

\begin{savenotes}\sphinxattablestart
\sphinxthistablewithglobalstyle
\centering
\begin{tabulary}{\linewidth}[t]{T}
\sphinxtoprule
\sphinxtableatstartofbodyhook

\begin{sphinxfigure-in-table}
\centering
\capstart
\noindent\sphinxincludegraphics[width=300\sphinxpxdimen]{{image012}.png}
\sphinxfigcaption{Modelo de CPU}\label{\detokenize{cursos/logica8/recursos/T01:id13}}\end{sphinxfigure-in-table}\relax
\\
\sphinxbottomrule
\end{tabulary}
\sphinxtableafterendhook\par
\sphinxattableend\end{savenotes}
\begin{quote}

\sphinxAtStartPar
\sphinxstyleemphasis{“Muchas personas confunden la CPU con una torre de PC. Realmente son cosas diferentes”}.
\end{quote}

\sphinxAtStartPar
\sphinxfootnotemark[1] La CPU es considerada el “cerebro” del computador. Es el elemento principal que se encarga de:
\begin{itemize}
\item {} 
\sphinxAtStartPar
Procesar todas las instrucciones y cálculos del sistema

\item {} 
\sphinxAtStartPar
Coordinar el funcionamiento de los demás componentes

\item {} 
\sphinxAtStartPar
Ejecutar los programas y aplicaciones

\end{itemize}

\sphinxAtStartPar
La CPU está compuesta principalmente por:
\begin{itemize}
\item {} 
\sphinxAtStartPar
\sphinxstylestrong{Unidad de control}: Coordina y dirige todas las operaciones.

\item {} 
\sphinxAtStartPar
\sphinxstylestrong{ALU}: Unidad Lógica Aritmética. Realiza cálculos matemáticos y operaciones lógicas

\item {} 
\sphinxAtStartPar
\sphinxstylestrong{Registros}: Son unidades pequeñas de memoria de alta velocidad

\end{itemize}

\sphinxAtStartPar
La velocidad de la CPU se mide en Hertz (Hz) y determina qué tan rápido puede procesar las instrucciones. A mayor velocidad, mejor rendimiento del sistema.


\bigskip\hrule\bigskip



\subparagraph{¿Qué es la MEMORIA RAM?}
\label{\detokenize{cursos/logica8/recursos/T01:que-es-la-memoria-ram}}

\begin{savenotes}\sphinxattablestart
\sphinxthistablewithglobalstyle
\centering
\begin{tabulary}{\linewidth}[t]{T}
\sphinxtoprule
\sphinxtableatstartofbodyhook

\begin{sphinxfigure-in-table}
\centering
\capstart
\noindent\sphinxincludegraphics[width=300\sphinxpxdimen]{{image021}.png}
\sphinxfigcaption{Modelo de memoria RAM}\label{\detokenize{cursos/logica8/recursos/T01:id14}}\end{sphinxfigure-in-table}\relax
\\
\sphinxbottomrule
\end{tabulary}
\sphinxtableafterendhook\par
\sphinxattableend\end{savenotes}

\sphinxAtStartPar
\sphinxfootnotemark[1] La memoria RAM (Memoria de acceso aleatorio) tiene las siguientes funciones:
\begin{itemize}
\item {} 
\sphinxAtStartPar
ALMACENA temporalmente los datos y programas que el computador está usando activamente.

\item {} 
\sphinxAtStartPar
ACCEDE rápidamente a la información (SI y SOLO SI el equipo está encendido).

\end{itemize}

\sphinxAtStartPar
Las características de la memoria RAM son:
\begin{itemize}
\item {} 
\sphinxAtStartPar
Se borra cuando se apaga el computador (memoria volátil).

\item {} 
\sphinxAtStartPar
A mayor cantidad de RAM, más programas pueden ejecutarse simultáneamente.

\item {} 
\sphinxAtStartPar
Su velocidad afecta directamente el rendimiento del sistema.

\item {} 
\sphinxAtStartPar
Es diferente al almacenamiento permanente (disco duro) porque la memoria RAM es temporal.

\end{itemize}


\bigskip\hrule\bigskip



\paragraph{¿Qué es el SOFTWARE?}
\label{\detokenize{cursos/logica8/recursos/T01:que-es-el-software}}

\begin{savenotes}\sphinxattablestart
\sphinxthistablewithglobalstyle
\centering
\begin{tabulary}{\linewidth}[t]{T}
\sphinxtoprule
\sphinxtableatstartofbodyhook

\begin{sphinxfigure-in-table}
\centering
\capstart
\noindent\sphinxincludegraphics[width=300\sphinxpxdimen]{{image031}.png}
\sphinxfigcaption{Ejemplo de interfaz gráfica}\label{\detokenize{cursos/logica8/recursos/T01:id15}}\end{sphinxfigure-in-table}\relax
\\
\sphinxbottomrule
\end{tabulary}
\sphinxtableafterendhook\par
\sphinxattableend\end{savenotes}

\sphinxAtStartPar
\sphinxfootnotemark[1] El SOFTWARE, es el conjunto de: \sphinxstylestrong{programas}, \sphinxstylestrong{bases de datos}, diferentes tipos de \sphinxstylestrong{archivos} y \sphinxstylestrong{documentación} que le dicen al HARDWARE qué hacer. A diferencia del hardware que es físico y tangible, el software no se puede tocar.

\sphinxAtStartPar
El software solo puede verse su funcionamiento en la pantalla. El software es la parte lógica e intangible de un sistema informático. El software se clasifica de la siguiente manera:
\begin{itemize}
\item {} 
\sphinxAtStartPar
\sphinxstylestrong{Sistema:} Software que controla el funcionamiento básico del computador como por ejemplo: Firmware BIOS, Firmware UEFI, OS Windows, OS Linux, MacOS, Controladores o Drivers.

\item {} 
\sphinxAtStartPar
\sphinxstylestrong{Aplicaciones:} Diferente software que funciona sobre un sistema operativo como por ejemplo: Procesadores de texto, navegadores web, juegos, editores de imagen, modelado en 3D. La mayoría suelen tener interfáz gráfica.

\end{itemize}

\sphinxAtStartPar
Una \sphinxstylestrong{interfaz gráfica} es un sistema que permite a los usuarios interactuar con dispositivos electrónicos a través de elementos visuales en la pantalla, como iconos, botones y menús. Esto facilita la realización de acciones sin necesidad de conocimientos de programación, haciendo la experiencia más intuitiva y accesible.


\bigskip\hrule\bigskip



\paragraph{¿Qué es una SISTEMA COMPUTACIONAL?}
\label{\detokenize{cursos/logica8/recursos/T01:que-es-una-sistema-computacional}}
\sphinxAtStartPar
%
\begin{footnote}[2]\sphinxAtStartFootnote
Universidad Nacional del Nordeste. (n.d.). Informática. Conceptos fundamentales. https://exa.unne.edu.ar/ingenieria/computacion/Tema1.pdf
%
\end{footnote} Un \sphinxstylestrong{sistema computacional} es el conjunto de hardware que interactúa con un software para almacenar, procesar y mostrar información.


\begin{savenotes}\sphinxattablestart
\sphinxthistablewithglobalstyle
\centering
\begin{tabulary}{\linewidth}[t]{T}
\sphinxtoprule
\sphinxtableatstartofbodyhook

\begin{sphinxfigure-in-table}
\centering
\capstart
\noindent\sphinxincludegraphics[width=300\sphinxpxdimen]{{image041}.png}
\sphinxfigcaption{Computadora}\label{\detokenize{cursos/logica8/recursos/T01:id16}}\end{sphinxfigure-in-table}\relax
\\
\sphinxbottomrule
\end{tabulary}
\sphinxtableafterendhook\par
\sphinxattableend\end{savenotes}


\bigskip\hrule\bigskip



\paragraph{¿Qué es la INFORMÁTICA?}
\label{\detokenize{cursos/logica8/recursos/T01:que-es-la-informatica}}
\sphinxAtStartPar
\sphinxfootnotemark[2] \sphinxstylestrong{La informática} es la ciencia o disciplina que estudia el diseño, construcción e implementación de \sphinxstylestrong{sistemas computacionales} para garantizar la \sphinxstylestrong{integridad} y \sphinxstylestrong{disponibilidad} de la información.
\begin{itemize}
\item {} 
\sphinxAtStartPar
REGLA 01: (Integridad de la información)
\begin{quote}

\sphinxAtStartPar
\sphinxstyleemphasis{“Toda información debe conservar su integridad para que no pueda ser modifica sin autorización”.}
\end{quote}

\item {} 
\sphinxAtStartPar
REGLA 02: (Disponibilidad de la información)
\begin{quote}

\sphinxAtStartPar
\sphinxstyleemphasis{“Toda información debe estar disponible y llegar oportunamente”.}
\end{quote}

\end{itemize}


\bigskip\hrule\bigskip



\paragraph{¿Qué es RAZONAR?}
\label{\detokenize{cursos/logica8/recursos/T01:que-es-razonar}}
\sphinxAtStartPar
%
\begin{footnote}[3]\sphinxAtStartFootnote
Trejos Buriticá, O. I. (2017). Lógica de programación. Ediciones de la U.
%
\end{footnote} \sphinxstylestrong{Razonar} es la \sphinxstyleemphasis{actividad mental} mediante la cual se relacionan ideas o premisas para inferir una conclusión de manera coherente o justificada.

\sphinxAtStartPar
Por ejemplo:
\begin{itemize}
\item {} 
\sphinxAtStartPar
Premisa 01: Todos los humanos son mortales

\item {} 
\sphinxAtStartPar
Premisa 02: Sócrates es humano

\item {} 
\sphinxAtStartPar
\sphinxstylestrong{Conclusión:} Sócrates es mortal

\end{itemize}


\bigskip\hrule\bigskip



\paragraph{¿Qué es la LÓGICA?}
\label{\detokenize{cursos/logica8/recursos/T01:que-es-la-logica}}
\sphinxAtStartPar
\sphinxfootnotemark[3] La \sphinxstylestrong{lógica} es la \sphinxstyleemphasis{disciplina} que permite analizar y verificar la validez de los razonamientos, determinando si una conclusión se deriva correctamente de unas premisas.

\sphinxAtStartPar
Por ejemplo:

\sphinxAtStartPar
Para el siguiente razonamiento, verifique si es correcto o no:
\begin{itemize}
\item {} 
\sphinxAtStartPar
Premisa 01: Todos los humanos son mortales

\item {} 
\sphinxAtStartPar
Premisa 02: Sócrates es humano

\item {} 
\sphinxAtStartPar
\sphinxstylestrong{Conclusión:} Sócrates es mortal

\end{itemize}
\begin{quote}

\sphinxAtStartPar
\sphinxstylestrong{ANÁLISIS:} Si Sócrates es un ser humano, entonces Sócrates es mortal porque todos los humanos son mortales. El razonamiento es correcto.
\end{quote}


\bigskip\hrule\bigskip



\paragraph{DEFINICIÓN DE LÓGICA COMPUTACIONAL}
\label{\detokenize{cursos/logica8/recursos/T01:definicion-de-logica-computacional}}
\sphinxAtStartPar
\sphinxfootnotemark[3] La \sphinxstylestrong{lógica computacional} es el área de la \sphinxstylestrong{informática} que aplica los principios de la \sphinxstylestrong{lógica} para representar y automatizar el \sphinxstylestrong{razonamiento}, con el objetivo de resolver problemas de forma estructurada y eficiente.


\bigskip\hrule\bigskip


\sphinxstepscope


\chapter{Informática 10}
\label{\detokenize{cursos/informatica10/index:informatica-10}}\label{\detokenize{cursos/informatica10/index::doc}}
\sphinxAtStartPar
Bienvenido al curso Tecnología e Informática de grado décimo.

\sphinxAtStartPar
Acceda a las guías didácticas según el periodo.

\sphinxstepscope


\section{Periodo 01}
\label{\detokenize{cursos/informatica10/periodo01/index:periodo-01}}\label{\detokenize{cursos/informatica10/periodo01/index::doc}}\begin{quote}\begin{description}
\sphinxlineitem{DESEMPEÑO}
\sphinxAtStartPar
Conoce el origen de diferentes tecnologías relacionadas con redes, INTERNET, computación personal y software.

\end{description}\end{quote}


\bigskip\hrule\bigskip



\subsection{Contenidos}
\label{\detokenize{cursos/informatica10/periodo01/index:contenidos}}
\sphinxstepscope


\subsubsection{Guía 01: Historia de la computación personal}
\label{\detokenize{cursos/informatica10/periodo01/G01:guia-01-historia-de-la-computacion-personal}}\label{\detokenize{cursos/informatica10/periodo01/G01::doc}}

\begin{savenotes}\sphinxattablestart
\sphinxthistablewithglobalstyle
\centering
\begin{tabular}[t]{\X{35}{100}\X{65}{100}}
\sphinxtoprule
\sphinxtableatstartofbodyhook

\begin{savenotes}\sphinxattablestart
\sphinxthistablewithglobalstyle
\centering
\begin{tabulary}{\linewidth}[t]{T}
\sphinxtoprule
\sphinxtableatstartofbodyhook
\sphinxAtStartPar
\sphinxincludegraphics{{escudo}.png}
\\
\sphinxbottomrule
\end{tabulary}
\sphinxtableafterendhook\par
\sphinxattableend\end{savenotes}
&

\begin{savenotes}\sphinxattablestart
\sphinxthistablewithglobalstyle
\centering
\begin{tabulary}{\linewidth}[t]{T}
\sphinxtoprule
\sphinxtableatstartofbodyhook
\sphinxAtStartPar
\sphinxstylestrong{INSTITUCIÓN EDUCATIVA}: John F. Kennedy
\\
\sphinxhline
\sphinxAtStartPar
\sphinxstylestrong{DOCENTE}: Julian Camilo Fonseca Romero
\\
\sphinxhline
\sphinxAtStartPar
\sphinxstylestrong{ASIGNATURA}: Tecnología e Informática
\\
\sphinxhline
\sphinxAtStartPar
\sphinxstylestrong{GRADO}: DÉCIMO
\\
\sphinxhline
\sphinxAtStartPar
\sphinxstylestrong{TIEMPO}: 10 semanas (2 horas / semana)
\\
\sphinxhline
\sphinxAtStartPar
\sphinxstylestrong{TIEMPO DE TRABAJO}: 20/20 horas
\\
\sphinxbottomrule
\end{tabulary}
\sphinxtableafterendhook\par
\sphinxattableend\end{savenotes}
\\
\sphinxbottomrule
\end{tabular}
\sphinxtableafterendhook\par
\sphinxattableend\end{savenotes}

\begin{sphinxuseclass}{sd-tab-set}
\begin{sphinxuseclass}{sd-tab-item}\subsubsection*{Objetivo}

\begin{sphinxuseclass}{sd-tab-content}\begin{itemize}
\item {} 
\sphinxAtStartPar
\sphinxstylestrong{Temática}: Historia de la computación personal.

\item {} 
\sphinxAtStartPar
\sphinxstylestrong{Objetivo de aprendizaje}: Conocer los orígenes de la computación personal por medio del estudio de la evolución del internet, el hardware y el software.

\end{itemize}

\end{sphinxuseclass}
\end{sphinxuseclass}
\begin{sphinxuseclass}{sd-tab-item}\subsubsection*{Orientaciones curriculares}

\begin{sphinxuseclass}{sd-tab-content}\begin{itemize}
\item {} 
\sphinxAtStartPar
\sphinxstylestrong{Componente}: Naturaleza y evolución de la TI.

\item {} 
\sphinxAtStartPar
\sphinxstylestrong{Competencia}: Construyo conocimientos y saberes de base tecnológica e informática para la toma de decisiones en el desarrollo de productos tecnológicos.

\item {} 
\sphinxAtStartPar
\sphinxstylestrong{Evidencia de aprendizaje}: Sistematizo la evolución tecnológica e informática y la manera cómo incidieron en los cambios estructurales de la sociedad y la cultura a lo largo de la historia.

\end{itemize}

\end{sphinxuseclass}
\end{sphinxuseclass}
\begin{sphinxuseclass}{sd-tab-item}\subsubsection*{Criterios de evaluación}

\begin{sphinxuseclass}{sd-tab-content}\begin{itemize}
\item {} 
\sphinxAtStartPar
Actividades desarrolladas en su totalidad.

\item {} 
\sphinxAtStartPar
Participación en clase.

\item {} 
\sphinxAtStartPar
Explicación clara de conceptos.

\item {} 
\sphinxAtStartPar
Argumentación de ideas.

\item {} 
\sphinxAtStartPar
Cumplimiento de las reglas del aula.

\end{itemize}

\end{sphinxuseclass}
\end{sphinxuseclass}
\begin{sphinxuseclass}{sd-tab-item}\subsubsection*{Recursos (1)}

\begin{sphinxuseclass}{sd-tab-content}\begin{itemize}
\item {} 
\sphinxAtStartPar
{\hyperref[\detokenize{cursos/informatica10/recursos/T01::doc}]{\sphinxcrossref{\DUrole{std}{\DUrole{std-doc}{\sphinxstylestrong{TEXTO:} Historia de la computación personal}}}}}

\end{itemize}

\end{sphinxuseclass}
\end{sphinxuseclass}
\end{sphinxuseclass}

\paragraph{Mensaje del autor}
\label{\detokenize{cursos/informatica10/periodo01/G01:mensaje-del-autor}}\begin{quote}
\begin{quote}

\sphinxAtStartPar
“En esta guía, el estudiante aprende sobre el origen de la computación personal para entender la evolución de nuestra propia cultura. Se invita a explorar cómo el desarrollo de Internet, junto con la evolución del hardware y el software, transformó la estructura de la sociedad actual. Al finalizar, el estudiante podrá sistematizar estos hitos históricos y reconocer la huella que la informática ha dejado en cada aspecto de nuestra vida cotidiana.”
\end{quote}

\begin{flushright}
---Camilo Fonseca
\end{flushright}
\end{quote}


\bigskip\hrule\bigskip



\paragraph{A00: Introducción (+0,1)}
\label{\detokenize{cursos/informatica10/periodo01/G01:a00-introduccion-0-1}}
\sphinxAtStartPar
\DUrole{sd-sphinx-override}{\DUrole{sd-badge}{\DUrole{sd-bg-light}{\DUrole{sd-bg-text-light}{Exploración}}}}
\begin{enumerate}
\sphinxsetlistlabels{\arabic}{enumi}{enumii}{}{.}%
\item {} 
\sphinxAtStartPar
TRANSCRIBA en su cuaderno el objetivo, las orientaciones curriculares y los criterios de evaluación de la presente guía.

\item {} 
\sphinxAtStartPar
LEA el \sphinxstylestrong{mensaje del autor} y en su cuaderno responda:
\begin{itemize}
\item {} 
\sphinxAtStartPar
¿Por qué es importante estudiar la historia de la computación personal?

\end{itemize}

\end{enumerate}


\bigskip\hrule\bigskip



\paragraph{A01: Orígenes de INTERNET (+0,1)}
\label{\detokenize{cursos/informatica10/periodo01/G01:a01-origenes-de-internet-0-1}}
\sphinxAtStartPar
\DUrole{sd-sphinx-override}{\DUrole{sd-badge}{\DUrole{sd-bg-info}{\DUrole{sd-bg-text-info}{Estructuración}}}}
\begin{enumerate}
\sphinxsetlistlabels{\arabic}{enumi}{enumii}{}{.}%
\item {} 
\sphinxAtStartPar
En su cuaderno, responda las siguientes preguntas:
\begin{itemize}
\item {} 
\sphinxAtStartPar
¿Cómo se llamó la primera red?

\item {} 
\sphinxAtStartPar
¿En que año y en que lugar se creó la primera red?

\item {} 
\sphinxAtStartPar
¿Quienes crearon la tecnología para que funcionara el Internet?

\item {} 
\sphinxAtStartPar
¿En que año se creó la tecnología para que funcionara el Internet y que protocolos fueron usados en dicha tecnología?

\end{itemize}

\end{enumerate}


\bigskip\hrule\bigskip



\paragraph{A02: ARPANET (+0,1)}
\label{\detokenize{cursos/informatica10/periodo01/G01:a02-arpanet-0-1}}
\sphinxAtStartPar
\DUrole{sd-sphinx-override}{\DUrole{sd-badge}{\DUrole{sd-bg-info}{\DUrole{sd-bg-text-info}{Estructuración}}}}
\begin{enumerate}
\sphinxsetlistlabels{\arabic}{enumi}{enumii}{}{.}%
\item {} 
\sphinxAtStartPar
En su cuaderno dibuje el mapa de la primera red y responda las siguientes preguntas:
\begin{itemize}
\item {} 
\sphinxAtStartPar
¿Cuántos nodos en el occidente estaban interconectados?

\item {} 
\sphinxAtStartPar
¿Cuántos nodos en el oriente estaban interconectados?

\end{itemize}

\end{enumerate}


\bigskip\hrule\bigskip



\paragraph{A03: Primera página web (+0,1)}
\label{\detokenize{cursos/informatica10/periodo01/G01:a03-primera-pagina-web-0-1}}
\sphinxAtStartPar
\DUrole{sd-sphinx-override}{\DUrole{sd-badge}{\DUrole{sd-bg-info}{\DUrole{sd-bg-text-info}{Estructuración}}}}
\begin{enumerate}
\sphinxsetlistlabels{\arabic}{enumi}{enumii}{}{.}%
\item {} 
\sphinxAtStartPar
Responda la siguiente pregunta en su cuaderno:
\begin{itemize}
\item {} 
\sphinxAtStartPar
¿En que año se crea el primer navegador y la primera página web?

\end{itemize}

\item {} 
\sphinxAtStartPar
Traduzca los contenidos de la primera página web y haga un breve resumen de dichos contenidos en su cuaderno.

\end{enumerate}


\bigskip\hrule\bigskip



\paragraph{A04: Primer navegador MOSAIC (+0,1)}
\label{\detokenize{cursos/informatica10/periodo01/G01:a04-primer-navegador-mosaic-0-1}}
\sphinxAtStartPar
\DUrole{sd-sphinx-override}{\DUrole{sd-badge}{\DUrole{sd-bg-info}{\DUrole{sd-bg-text-info}{Estructuración}}}}
\begin{enumerate}
\sphinxsetlistlabels{\arabic}{enumi}{enumii}{}{.}%
\item {} 
\sphinxAtStartPar
En su cuaderno DIBUJE un \sphinxstylestrong{boceto} de la \sphinxstyleemphasis{interfaz gráfica} del primer navegador web \sphinxstylestrong{MOSAIC} y realice una breve descripción de las diferencias de este primer navegador con navegadores actuales como GOOGLE CHROME, MOZILLA FIREFOX o EDGE.

\end{enumerate}


\bigskip\hrule\bigskip



\paragraph{A05: Orígenes del hardware (+0,1)}
\label{\detokenize{cursos/informatica10/periodo01/G01:a05-origenes-del-hardware-0-1}}
\sphinxAtStartPar
\DUrole{sd-sphinx-override}{\DUrole{sd-badge}{\DUrole{sd-bg-info}{\DUrole{sd-bg-text-info}{Estructuración}}}}
\begin{enumerate}
\sphinxsetlistlabels{\arabic}{enumi}{enumii}{}{.}%
\item {} 
\sphinxAtStartPar
En su cuaderno responda las siguientes preguntas:
\begin{itemize}
\item {} 
\sphinxAtStartPar
¿Cuál fue la primera computadora personal de la historia y en que año se creo?

\item {} 
\sphinxAtStartPar
¿Qué lenguajes de programación se usaron en la primera computadora personal?

\item {} 
\sphinxAtStartPar
¿Cómo APPLE se relaciona con la historia de la computadora personal?

\item {} 
\sphinxAtStartPar
¿Quién fue Steve Jobs?

\item {} 
\sphinxAtStartPar
¿Cómo se democratizó el uso de las computadoras personales?

\end{itemize}

\end{enumerate}


\bigskip\hrule\bigskip



\paragraph{A06: Apple \sphinxhyphen{} Macintosh (+0,1)}
\label{\detokenize{cursos/informatica10/periodo01/G01:a06-apple-macintosh-0-1}}
\sphinxAtStartPar
\DUrole{sd-sphinx-override}{\DUrole{sd-badge}{\DUrole{sd-bg-info}{\DUrole{sd-bg-text-info}{Estructuración}}}}
\begin{enumerate}
\sphinxsetlistlabels{\arabic}{enumi}{enumii}{}{.}%
\item {} 
\sphinxAtStartPar
En su cuaderno responda:
\begin{itemize}
\item {} 
\sphinxAtStartPar
¿Qué características tenía la primera Macintosh y en que año fue creada?

\end{itemize}

\item {} 
\sphinxAtStartPar
En su cuaderno DIBUJE un bosquejo de la primera Macintosh

\item {} 
\sphinxAtStartPar
Observe el artículo donde se describe la Macintosh y en su cuaderno describa el significado de las siguientes frases:
\begin{itemize}
\item {} 
\sphinxAtStartPar
\sphinxstyleemphasis{“La primera APPLE que se puede llevar en un bolso”}

\item {} 
\sphinxAtStartPar
\sphinxstyleemphasis{“Si puedes señalar, puedes usar Macintosh”}

\item {} 
\sphinxAtStartPar
\sphinxstyleemphasis{“PERSONALIDAD DE MACINTOSH: Lado divertido, Lado serio”}

\item {} 
\sphinxAtStartPar
\sphinxstyleemphasis{“Muy pronto habrán dos tipos de personas: Aquellos que usan computadoras y aquellos que usan APPLES”}

\end{itemize}

\end{enumerate}


\bigskip\hrule\bigskip



\paragraph{A07: Orígenes del software (+0,1)}
\label{\detokenize{cursos/informatica10/periodo01/G01:a07-origenes-del-software-0-1}}
\sphinxAtStartPar
\DUrole{sd-sphinx-override}{\DUrole{sd-badge}{\DUrole{sd-bg-info}{\DUrole{sd-bg-text-info}{Estructuración}}}}
\begin{enumerate}
\sphinxsetlistlabels{\arabic}{enumi}{enumii}{}{.}%
\item {} 
\sphinxAtStartPar
En su cuaderno responda las siguientes preguntas:
\begin{itemize}
\item {} 
\sphinxAtStartPar
¿En que año y como operaban las computadoras en sus orígenes?

\item {} 
\sphinxAtStartPar
¿Qué es y en que año se crea \sphinxstylestrong{FORTRAN}?

\item {} 
\sphinxAtStartPar
¿Por qué la década de los 80 fue revolucionaria en el software?

\end{itemize}

\end{enumerate}


\bigskip\hrule\bigskip



\paragraph{A08: UNIX\sphinxhyphen{}LINUX (+0,1)}
\label{\detokenize{cursos/informatica10/periodo01/G01:a08-unix-linux-0-1}}
\sphinxAtStartPar
\DUrole{sd-sphinx-override}{\DUrole{sd-badge}{\DUrole{sd-bg-info}{\DUrole{sd-bg-text-info}{Estructuración}}}}
\begin{enumerate}
\sphinxsetlistlabels{\arabic}{enumi}{enumii}{}{.}%
\item {} 
\sphinxAtStartPar
Investigue en \sphinxstylestrong{INTERNET} o \sphinxstylestrong{IA} y realice una descripción del sistema operativo \sphinxstylestrong{UNIX} luego en su cuaderno responda la siguiente pregunta:
\begin{itemize}
\item {} 
\sphinxAtStartPar
¿Por qué UNIX es considerado el padre de la mayoría de sistemas operativos en la actualidad?

\end{itemize}

\item {} 
\sphinxAtStartPar
Investigue en \sphinxstylestrong{INTERNET} la diferencia entre un \sphinxstylestrong{equipo cliente} y un \sphinxstylestrong{equipo servidor} y luego realice una descripción en su cuaderno.

\item {} 
\sphinxAtStartPar
Investigue en \sphinxstylestrong{INTERNET} que es el sistema operativo \sphinxstylestrong{LINUX} luego realice una descripción en su cuaderno e indique por que se usa en la mayoría de servidores a nivel mundial.

\item {} 
\sphinxAtStartPar
En su cuaderno describa 10 distribuciones \sphinxstylestrong{LINUX} (con sus fechas de creación).

\end{enumerate}


\bigskip\hrule\bigskip



\paragraph{A09: Windows (+0,1)}
\label{\detokenize{cursos/informatica10/periodo01/G01:a09-windows-0-1}}
\sphinxAtStartPar
\DUrole{sd-sphinx-override}{\DUrole{sd-badge}{\DUrole{sd-bg-info}{\DUrole{sd-bg-text-info}{Estructuración}}}}
\begin{enumerate}
\sphinxsetlistlabels{\arabic}{enumi}{enumii}{}{.}%
\item {} 
\sphinxAtStartPar
En su cuaderno responda:
\begin{itemize}
\item {} 
\sphinxAtStartPar
¿En que año se creo el sistema operativo \sphinxstylestrong{Microsoft DOS (MS\sphinxhyphen{}DOS)} y por que fue tan importante?

\end{itemize}

\item {} 
\sphinxAtStartPar
Con ayuda de INTERNET o IA, en su cuaderno realice una descripción (con fechas de creación) de los siguientes sistemas operativos de Microsoft:
\begin{itemize}
\item {} 
\sphinxAtStartPar
Windows 3.1

\item {} 
\sphinxAtStartPar
Windows 95

\item {} 
\sphinxAtStartPar
Windows 98

\item {} 
\sphinxAtStartPar
Windows 2000

\item {} 
\sphinxAtStartPar
Windows Me

\item {} 
\sphinxAtStartPar
Windows XP

\item {} 
\sphinxAtStartPar
Windows Vista

\item {} 
\sphinxAtStartPar
Windows 7

\item {} 
\sphinxAtStartPar
Windows 8

\item {} 
\sphinxAtStartPar
Windows 8.1

\item {} 
\sphinxAtStartPar
Windows 10

\item {} 
\sphinxAtStartPar
Windows 11

\end{itemize}

\end{enumerate}


\bigskip\hrule\bigskip



\paragraph{A10: Línea de tiempo (+1,0)}
\label{\detokenize{cursos/informatica10/periodo01/G01:a10-linea-de-tiempo-1-0}}
\sphinxAtStartPar
\DUrole{sd-sphinx-override}{\DUrole{sd-badge}{\DUrole{sd-bg-danger}{\DUrole{sd-bg-text-danger}{Transferencia}}}}
\begin{enumerate}
\sphinxsetlistlabels{\arabic}{enumi}{enumii}{}{.}%
\item {} 
\sphinxAtStartPar
En una \sphinxstylestrong{hoja de tamaño oficio}, cree una línea de tiempo de la \sphinxstylestrong{historia de la computación personal}, desde los orígenes de la primera red con \sphinxstylestrong{ARPANET} hasta la creación de los \sphinxstylestrong{sistemas operativos de Microsoft}.

\end{enumerate}


\bigskip\hrule\bigskip



\paragraph{A11: Examen final (+3,0)}
\label{\detokenize{cursos/informatica10/periodo01/G01:a11-examen-final-3-0}}
\sphinxAtStartPar
\DUrole{sd-sphinx-override}{\DUrole{sd-badge}{\DUrole{sd-bg-danger}{\DUrole{sd-bg-text-danger}{Transferencia}}}}
\begin{quote}\begin{description}
\sphinxlineitem{TIEMPO}
\sphinxAtStartPar
30 minutos

\sphinxlineitem{PREGUNTAS}
\sphinxAtStartPar
11

\end{description}\end{quote}

\begin{sphinxadmonition}{note}{posibles preguntas}
\begin{enumerate}
\sphinxsetlistlabels{\arabic}{enumi}{enumii}{}{.}%
\item {} 
\sphinxAtStartPar
¿Qué es una red de datos y cuál fue la primera red a nivel mundial, (en que año y en que lugar se creo)?

\item {} 
\sphinxAtStartPar
¿Quién fue la persona que creo la tecnología para que funcionara el INTERNET?

\item {} 
\sphinxAtStartPar
¿En que año se creo la tecnología para que funcionara el INTERNET y que protocolos fueron usados en dicha tecnología?

\item {} 
\sphinxAtStartPar
¿En que año se crea el primer navegador y la primera página web?

\item {} 
\sphinxAtStartPar
Realice un descripción del navegador MOSAIC e indique sus diferencias con navegadores actuales

\item {} 
\sphinxAtStartPar
¿Qué es una computadora y que características debe tener para que sea considerada personal?

\item {} 
\sphinxAtStartPar
¿Cuál fue la primera computadora personal de la historia?

\item {} 
\sphinxAtStartPar
¿Qué lenguajes de programación se usaron en la primera computadora personal?

\item {} 
\sphinxAtStartPar
¿Cómo APPLE se relaciona con la historia de la computadora personal?

\item {} 
\sphinxAtStartPar
¿Cómo se democratizó el uso de las computadoras personales?

\item {} 
\sphinxAtStartPar
¿Qué características tenía la primera Macintosh y en que año fue creada?

\item {} 
\sphinxAtStartPar
¿En que año se creo el sistema operativo Microsoft DOS y por que fue tan importante?

\item {} 
\sphinxAtStartPar
Realice una descripción de tres sistemas operativos de Microsoft e indique por que se caracterizan

\item {} 
\sphinxAtStartPar
¿En que año y como operaban las computadoras en sus orígenes?

\item {} 
\sphinxAtStartPar
¿Qué es y en que año se crea FORTRAN?

\item {} 
\sphinxAtStartPar
¿Qué es un sistema operativo?

\item {} 
\sphinxAtStartPar
¿Por qué la década de los 80 fue revolucionaria en el software?

\item {} 
\sphinxAtStartPar
¿Qué es el sistema operativo LINUX y en que se diferencia con el sistema operativo UNIX?

\item {} 
\sphinxAtStartPar
Realice la descripción de tres distribuciones de sistema operativo tipo LINUX

\item {} 
\sphinxAtStartPar
¿En que se diferencian los sistemas operativos LINUX y Microsoft WINDOWS?

\end{enumerate}
\end{sphinxadmonition}

\sphinxstepscope


\subsubsection{Plan de refuerzo 01}
\label{\detokenize{cursos/informatica10/periodo01/R01:plan-de-refuerzo-01}}\label{\detokenize{cursos/informatica10/periodo01/R01::doc}}
\begin{sphinxadmonition}{note}{EVALUACIÓN: Plan de refuerzo (10 minutos)}

\sphinxAtStartPar
Estudiar los temas a continuación. Se realiza examen de 5 preguntas. Si el estudiante no asiste en la fecha programada, este será evaluado la próxima vez que asista a clase.
\end{sphinxadmonition}


\bigskip\hrule\bigskip



\paragraph{Historia del Internet}
\label{\detokenize{cursos/informatica10/periodo01/R01:historia-del-internet}}
\sphinxAtStartPar
El Internet tiene sus orígenes en ARPANET, un proyecto militar desarrollado en 1969 por el Departamento de Defensa de Estados Unidos. Esta red inicial conectaba solo cuatro computadoras en diferentes universidades.

\sphinxAtStartPar
En los años 70 y 80 (exactamente en el año 1983), se desarrollaron los protocolos TCP/IP, que permitieron la comunicación entre diferentes redes de computadoras. La década de los 90 marcó un punto de inflexión con la creación de la World Wide Web por Tim Berners\sphinxhyphen{}Lee, que introdujo las páginas web y los navegadores en 1989.

\sphinxAtStartPar
Para facilitar el acceso a diferentes páginas web, en 1993 se creó el primer navegador web con interfaz gráfica de la historia. Este tenía por nombre \sphinxstylestrong{MOSAIC}.

\sphinxAtStartPar
El Internet moderno ha evolucionado, pasando por la era de la Web 2.0 (redes sociales e interactividad con su origen en el año 2004), hasta la actual era del Internet móvil desde el año 2007. Hoy en día la nube y el Internet de las Cosas (IoT) conectan a miles de millones de dispositivos y usuarios en todo el mundo.


\bigskip\hrule\bigskip



\paragraph{Historia del Hardware}
\label{\detokenize{cursos/informatica10/periodo01/R01:historia-del-hardware}}
\sphinxAtStartPar
La historia de la computadora personal comienza en la década de 1970. En 1975, la Altair 8800 se convirtió en la primera computadora personal comercialmente exitosa, aunque requería ser ensamblada por el usuario.

\sphinxAtStartPar
En 1976, Apple Computer lanzó la \sphinxstylestrong{Apple I}, seguida en 1977 por la revolucionaria \sphinxstylestrong{Apple II}, una de las primeras computadoras con gráficos a color.

\sphinxAtStartPar
IBM entró al mercado en 1981 con la \sphinxstylestrong{IBM PC}, estableciendo un estándar que dominaría la industria. Los años 80 vieron la llegada de la \sphinxstylestrong{Commodore 64} (1982) y la \sphinxstylestrong{Apple Macintosh} (1984), que introdujo la interfaz gráfica de usuario (GUI) al mercado masivo.

\sphinxAtStartPar
En los 90, \sphinxstylestrong{Windows 3.0} y posteriores versiones de Microsoft democratizaron el uso de las computadoras personales.

\sphinxAtStartPar
El siglo XXI trajo las laptops ultra\sphinxhyphen{}delgadas, tablets y mini\sphinxhyphen{}computadoras. En 2012, la \sphinxstylestrong{Raspberry\sphinxhyphen{}Pi} revolucionó el concepto de computadora personal al ofrecer un dispositivo del tamaño de una tarjeta de crédito a bajo costo, enfocado en la educación y proyectos de programación.


\bigskip\hrule\bigskip



\paragraph{Historia del Software}
\label{\detokenize{cursos/informatica10/periodo01/R01:historia-del-software}}
\sphinxAtStartPar
El software tiene su origen en los primeros días de la computación. \sphinxstylestrong{En la década de 1940}, las primeras computadoras ejecutaban instrucciones básicas mediante cables e interruptores físicos.

\sphinxAtStartPar
\sphinxstylestrong{Los años 50} marcaron el inicio de los lenguajes de programación modernos, con la creación del lenguaje de alto nivel \sphinxstylestrong{FORTRAN (1957)}. En el \sphinxstylestrong{año 1969} se desarrolló el sistema operativo \sphinxstylestrong{UNIX}. La mayoría de sistemas operativos en la actualidad se basan en UNIX.

\sphinxAtStartPar
Con la llegada de \sphinxstylestrong{ALTAIR en 1975,} Microsoft desarrolla el lenguaje de programación \sphinxstylestrong{BASIC}. Este permitió programar en dicha computadora.

\sphinxAtStartPar
\sphinxstylestrong{La década de 1980} fue revolucionaria con la llegada del software comercial para PC, incluyendo procesadores de texto, hojas de cálculo y los primeros sistemas operativos con interfaz gráfica. Esto inicio el \sphinxstylestrong{año 1981} con la creación del primer sistema operativo de Microsoft: \sphinxstylestrong{MS\sphinxhyphen{}DOS.}

\sphinxAtStartPar
Luego en el \sphinxstylestrong{año 1991} la comunidad de ingenieros y desarrolladores a nivel mundial crean el sistema operativo \sphinxstylestrong{LINUX}. El cual es de gran importancia hasta nuestros días.

\sphinxAtStartPar
En la actualidad, el software ha evolucionado con aplicaciones en la nube, inteligencia artificial y desarrollo ágil, transformando completamente la manera en la que interactuamos con la tecnología.

\sphinxstepscope


\section{Recursos}
\label{\detokenize{cursos/informatica10/recursos/index:recursos}}\label{\detokenize{cursos/informatica10/recursos/index::doc}}
\sphinxAtStartPar
En esta sección encuentra los recursos educativos necesarios para el desarrollo del curso.


\bigskip\hrule\bigskip



\subsection{Contenidos}
\label{\detokenize{cursos/informatica10/recursos/index:contenidos}}
\sphinxstepscope


\subsubsection{T01: Historia de la computación personal}
\label{\detokenize{cursos/informatica10/recursos/T01:t01-historia-de-la-computacion-personal}}\label{\detokenize{cursos/informatica10/recursos/T01::doc}}
\sphinxAtStartPar
En este texto, se muestra la historia de la computación personal. Conocer el origen de esta tecnología es importante porque permite conocer como han evolucionado las herramientas digitales que usamos a diario. De esta manera se muestra la historia del INTERNET, el hardware y el software.


\bigskip\hrule\bigskip



\paragraph{Historia del Internet (ARPANET)}
\label{\detokenize{cursos/informatica10/recursos/T01:historia-del-internet-arpanet}}
\sphinxAtStartPar
%
\begin{footnote}[1]\sphinxAtStartFootnote
ARPANET. (2005, Feb 10). Sistemas.com. Retrieved April 25, 2026, from https://sistemas.com/arpanet.php
%
\end{footnote} El Internet tiene sus orígenes en \sphinxstylestrong{ARPANET}, un proyecto militar desarrollado en 1969 por el \sphinxstyleemphasis{Departamento de defensa de Estados Unidos}. Esta red inicial conectaba un número reducido de computadoras en distintas universidades.


\begin{savenotes}\sphinxattablestart
\sphinxthistablewithglobalstyle
\centering
\begin{tabulary}{\linewidth}[t]{T}
\sphinxtoprule
\sphinxtableatstartofbodyhook

\begin{sphinxfigure-in-table}
\centering
\capstart
\noindent\sphinxincludegraphics[width=450\sphinxpxdimen]{{image01}.png}
\sphinxfigcaption{\sphinxstylestrong{ARPANET:} Primera red del mundo}\label{\detokenize{cursos/informatica10/recursos/T01:id26}}\end{sphinxfigure-in-table}\relax
\\
\sphinxbottomrule
\end{tabulary}
\sphinxtableafterendhook\par
\sphinxattableend\end{savenotes}

\sphinxAtStartPar
En los años 70 y 80 (exactamente en el año 1983), se desarrollaron los protocolos TCP/IP, que permitieron la comunicación entre diferentes redes de computadoras.

\sphinxAtStartPar
%
\begin{footnote}[2]\sphinxAtStartFootnote
Rojas, J. (2022, May 17). El primer sitio web de la historia. Un hito que cambió vidas. Telefónica. https://www.telefonica.com/es/sala\sphinxhyphen{}comunicacion/blog/el\sphinxhyphen{}primer\sphinxhyphen{}sitio\sphinxhyphen{}web\sphinxhyphen{}de\sphinxhyphen{}la\sphinxhyphen{}historia\sphinxhyphen{}un\sphinxhyphen{}hito\sphinxhyphen{}que\sphinxhyphen{}cambio\sphinxhyphen{}vidas/
%
\end{footnote} La década de los 90 marcó un punto de inflexión con la creación de la \sphinxstylestrong{World Wide Web} por \sphinxstylestrong{Tim Berners\sphinxhyphen{}Lee}, que introdujo las páginas web y los navegadores en 1989. A continuación se muestra la primera página web de la historia:


\begin{savenotes}\sphinxattablestart
\sphinxthistablewithglobalstyle
\centering
\begin{tabulary}{\linewidth}[t]{T}
\sphinxtoprule
\sphinxtableatstartofbodyhook

\begin{sphinxfigure-in-table}
\centering
\capstart
\noindent\sphinxincludegraphics[width=700\sphinxpxdimen]{{image02}.png}
\sphinxfigcaption{\sphinxstylestrong{Primera página web:} Word Wide Web}\label{\detokenize{cursos/informatica10/recursos/T01:id27}}\end{sphinxfigure-in-table}\relax
\\
\sphinxbottomrule
\end{tabulary}
\sphinxtableafterendhook\par
\sphinxattableend\end{savenotes}

\sphinxAtStartPar
%
\begin{footnote}[3]\sphinxAtStartFootnote
Pastor, J. (2018, April 28). Cuando Mosaic dominaba el mundo (de los navegadores). Xataka.com; Xataka. https://www.xataka.com/historia\sphinxhyphen{}tecnologica/cuando\sphinxhyphen{}mosaic\sphinxhyphen{}dominaba\sphinxhyphen{}el\sphinxhyphen{}mundo\sphinxhyphen{}de\sphinxhyphen{}los\sphinxhyphen{}navegadores
%
\end{footnote} Para facilitar el acceso a diferentes páginas web, en 1993 se creo el navegador web MOSAIC. Previamente existía otro navegador de nombre Nexus pero MOSAIC fue el primer navegador en usarse de forma masiva.


\begin{savenotes}\sphinxattablestart
\sphinxthistablewithglobalstyle
\centering
\begin{tabulary}{\linewidth}[t]{T}
\sphinxtoprule
\sphinxtableatstartofbodyhook

\begin{sphinxfigure-in-table}
\centering
\capstart
\noindent\sphinxincludegraphics[width=700\sphinxpxdimen]{{image03}.png}
\sphinxfigcaption{\sphinxstylestrong{MOSAIC:} Primer navegador web en usarse de forma masiva.}\label{\detokenize{cursos/informatica10/recursos/T01:id28}}\end{sphinxfigure-in-table}\relax
\\
\sphinxbottomrule
\end{tabulary}
\sphinxtableafterendhook\par
\sphinxattableend\end{savenotes}

\sphinxAtStartPar
%
\begin{footnote}[4]\sphinxAtStartFootnote
Eadicicco, L. (2025, May 11). La tecnologí­a que definió el internet moderno está cambiando y Silicon Valley finalmente lo admite. CNN en Español. https://cnnespanol.cnn.com/2025/05/11/ciencia/tecnologia\sphinxhyphen{}internet\sphinxhyphen{}moderno\sphinxhyphen{}silicon\sphinxhyphen{}valley\sphinxhyphen{}trax
%
\end{footnote} El \sphinxstylestrong{Internet moderno} ha evolucionado considerablemente desde entonces, pasando por la era de la Web 2.0 (redes sociales e interactividad las cuales tienen su origen en el año 2004), hasta la actual era del Internet móvil cuyo origen fue en el año 2007.  Hoy en día la nube y el Internet de las Cosas (IoT) conectan a miles de millones de dispositivos y usuarios en todo el mundo.


\bigskip\hrule\bigskip



\paragraph{Historia del Hardware}
\label{\detokenize{cursos/informatica10/recursos/T01:historia-del-hardware}}
\sphinxAtStartPar
La historia de la computadora personal comienza en la década de 1970.


\begin{savenotes}\sphinxattablestart
\sphinxthistablewithglobalstyle
\centering
\begin{tabulary}{\linewidth}[t]{T}
\sphinxtoprule
\sphinxtableatstartofbodyhook

\begin{sphinxfigure-in-table}
\centering
\capstart
\noindent\sphinxincludegraphics[width=300\sphinxpxdimen]{{image04}.png}
\sphinxfigcaption{\sphinxstylestrong{REVISTA:} Popular Electronics}\label{\detokenize{cursos/informatica10/recursos/T01:id29}}\end{sphinxfigure-in-table}\relax
\\
\sphinxbottomrule
\end{tabulary}
\sphinxtableafterendhook\par
\sphinxattableend\end{savenotes}

\sphinxAtStartPar
%
\begin{footnote}[5]\sphinxAtStartFootnote
Pastor, J. (2015, January 27). Altair 8800: todo comenzó hace ya 40 años. Xataka.com; Xataka. https://www.xataka.com/historia\sphinxhyphen{}tecnologica/altair\sphinxhyphen{}8800\sphinxhyphen{}todo\sphinxhyphen{}comenzo\sphinxhyphen{}hace\sphinxhyphen{}ya\sphinxhyphen{}40\sphinxhyphen{}anos
%
\end{footnote} En 1975, la \sphinxstylestrong{Altair 8800} se convirtió en la \sphinxstylestrong{primera computadora personal} comercialmente exitosa, aunque requería ser ensamblada por el usuario. Esta computadora más adelante sería más fácil de programar por la creación de un lenguaje de programación llamado \sphinxstylestrong{BASIC}. Anteriormente se programaba con el lenguaje \sphinxstylestrong{ENSAMBLADOR}.

\sphinxAtStartPar
%
\begin{footnote}[6]\sphinxAtStartFootnote
Catucci, A. (2020, September 9). La historia de Apple: repaso por este caso de éxito tecnológico. Marketing Insider Review. https://marketinginsiderreview.com/historia\sphinxhyphen{}apple\sphinxhyphen{}exito\sphinxhyphen{}tecnologico/
%
\end{footnote} En 1976, Apple Computer lanzó la \sphinxstylestrong{Apple I}, seguida en 1977 por la revolucionaria \sphinxstylestrong{Apple II}, una de las primeras computadoras con gráficos a color.

\sphinxAtStartPar
%
\begin{footnote}[7]\sphinxAtStartFootnote
The Conversation. (2024, January 22). La Mac cumple 40 años: debutó con este antológico aviso dirigido por Ridley Scott para mostrar cómo había otra forma de usar una computadora. LA NACION. https://www.lanacion.com.ar/tecnologia/la\sphinxhyphen{}mac\sphinxhyphen{}cumple\sphinxhyphen{}40\sphinxhyphen{}anos\sphinxhyphen{}debuto\sphinxhyphen{}con\sphinxhyphen{}este\sphinxhyphen{}antologico\sphinxhyphen{}aviso\sphinxhyphen{}dirigido\sphinxhyphen{}por\sphinxhyphen{}ridley\sphinxhyphen{}scott\sphinxhyphen{}para\sphinxhyphen{}mostrar\sphinxhyphen{}como\sphinxhyphen{}nid22012024/
%
\end{footnote} Los años 80 vieron la llegada de la primera \sphinxstylestrong{Macintosh}, conocida como \sphinxstylestrong{Macintosh 128K}, fue presentada el \sphinxstylestrong{24 de enero de 1984} por \sphinxstylestrong{Steve Jobs}. Este ordenador personal marcó un hito en la historia de la computación por varias razones:
\begin{itemize}
\item {} 
\sphinxAtStartPar
\sphinxstylestrong{Interfaz Gráfica de Usuario (GUI)}: Fue uno de los primeros ordenadores en popularizar el uso de una interfaz gráfica, lo que facilitó su uso para personas no técnicas.

\item {} 
\sphinxAtStartPar
\sphinxstylestrong{Ratón}: Incluía un ratón, lo que era innovador en ese momento y permitía una interacción más intuitiva con el sistema.

\item {} 
\sphinxAtStartPar
\sphinxstylestrong{Diseño Compacto}: Su diseño era más pequeño y accesible en comparación con otros ordenadores de la época, lo que lo hacía ideal para el uso personal.

\item {} 
\sphinxAtStartPar
\sphinxstylestrong{Software}: Venía con aplicaciones básicas como MacPaint y MacWrite, que demostraban las capacidades gráficas y de procesamiento de texto del dispositivo.

\end{itemize}


\begin{savenotes}\sphinxattablestart
\sphinxthistablewithglobalstyle
\centering
\begin{tabulary}{\linewidth}[t]{T}
\sphinxtoprule
\sphinxtableatstartofbodyhook

\begin{sphinxfigure-in-table}
\centering
\capstart
\noindent\sphinxincludegraphics[width=600\sphinxpxdimen]{{image05}.png}
\sphinxfigcaption{\sphinxstylestrong{Artículo:} Descripción de la Macintosh}\label{\detokenize{cursos/informatica10/recursos/T01:id30}}\end{sphinxfigure-in-table}\relax
\\
\sphinxbottomrule
\end{tabulary}
\sphinxtableafterendhook\par
\sphinxattableend\end{savenotes}

\sphinxAtStartPar
La \sphinxstylestrong{Macintosh 128K} no solo fue un éxito comercial, sino que también sentó las bases para el desarrollo de futuras computadoras de Apple.

\sphinxAtStartPar
%
\begin{footnote}[8]\sphinxAtStartFootnote
Historia de Microsoft. (2012, November 10). Blog de tecnología; Sistemas Ibertrónica. https://ibertronica.es/blog/productos/microsoft\sphinxhyphen{}pasado\sphinxhyphen{}presente\sphinxhyphen{}y\sphinxhyphen{}futuro/
%
\end{footnote} En los años 90, Microsoft democratizaron el uso de las computadoras personales con el desarrollo de su sistema operativo Windows. El siglo XXI trajo las \sphinxstyleemphasis{laptops ultradelgadas}, \sphinxstyleemphasis{tablets} y \sphinxstyleemphasis{mini\sphinxhyphen{}computadoras}.


\begin{savenotes}\sphinxattablestart
\sphinxthistablewithglobalstyle
\centering
\begin{tabulary}{\linewidth}[t]{T}
\sphinxtoprule
\sphinxtableatstartofbodyhook

\begin{sphinxfigure-in-table}
\centering
\capstart
\noindent\sphinxincludegraphics[width=500\sphinxpxdimen]{{image06}.png}
\sphinxfigcaption{\sphinxstylestrong{PC:} Raspberry Pi}\label{\detokenize{cursos/informatica10/recursos/T01:id31}}\end{sphinxfigure-in-table}\relax
\\
\sphinxbottomrule
\end{tabulary}
\sphinxtableafterendhook\par
\sphinxattableend\end{savenotes}

\sphinxAtStartPar
%
\begin{footnote}[9]\sphinxAtStartFootnote
Andrés, R. (2019, September 1). Raspberry Pi: la historia de un mini\sphinxhyphen{}PC que se ha convertido en un éxito de masas. Computer Hoy. https://computerhoy.20minutos.es/reportajes/tecnologia/historia\sphinxhyphen{}raspberry\sphinxhyphen{}pi\sphinxhyphen{}482477
%
\end{footnote} En 2012, la \sphinxstylestrong{Raspberry Pi} revolucionó el concepto de computadora personal al ofrecer un dispositivo del tamaño de una tarjeta de crédito a bajo costo, enfocado en la educación y proyectos de programación.


\bigskip\hrule\bigskip



\paragraph{Historia del Software}
\label{\detokenize{cursos/informatica10/recursos/T01:historia-del-software}}
\sphinxAtStartPar
El software tiene su origen en los primeros días de la computación. En la década de 1940, las primeras computadoras ejecutaban instrucciones básicas mediante cables e interruptores físicos.


\begin{savenotes}\sphinxattablestart
\sphinxthistablewithglobalstyle
\centering
\begin{tabulary}{\linewidth}[t]{T}
\sphinxtoprule
\sphinxtableatstartofbodyhook

\begin{sphinxfigure-in-table}
\centering
\capstart
\noindent\sphinxincludegraphics[width=600\sphinxpxdimen]{{image07}.png}
\sphinxfigcaption{Primera computadora programable 1940}\label{\detokenize{cursos/informatica10/recursos/T01:id32}}\end{sphinxfigure-in-table}\relax
\\
\sphinxbottomrule
\end{tabulary}
\sphinxtableafterendhook\par
\sphinxattableend\end{savenotes}

\sphinxAtStartPar
%
\begin{footnote}[10]\sphinxAtStartFootnote
Huertos, A. A. (2019, May 28). La historia de los lenguajes de programación. Computer Hoy. https://computerhoy.20minutos.es/reportajes/tecnologia/historia\sphinxhyphen{}lenguajes\sphinxhyphen{}programacion\sphinxhyphen{}428041
%
\end{footnote} Los años 50 marcaron el inicio de los lenguajes de programación modernos, con la creación del lenguaje de alto nivel \sphinxstylestrong{FORTRAN} (1957).

\sphinxAtStartPar
%
\begin{footnote}[11]\sphinxAtStartFootnote
Velásquez, W. (2022, April 1). 3 historias del software libre que deberías conocer. Wajari. https://wajari.com/blog/historias\sphinxhyphen{}software\sphinxhyphen{}libre\sphinxhyphen{}open\sphinxhyphen{}source/
%
\end{footnote} En el año 1969 se desarrolló el sistema operativo \sphinxstylestrong{UNIX}. La mayoría de sistemas operativos en la actualidad se basan en \sphinxstylestrong{UNIX}.


\begin{savenotes}\sphinxattablestart
\sphinxthistablewithglobalstyle
\centering
\begin{tabulary}{\linewidth}[t]{T}
\sphinxtoprule
\sphinxtableatstartofbodyhook

\begin{sphinxfigure-in-table}
\centering
\capstart
\noindent\sphinxincludegraphics[width=600\sphinxpxdimen]{{image08}.png}
\sphinxfigcaption{Consola terminal UNIX}\label{\detokenize{cursos/informatica10/recursos/T01:id33}}\end{sphinxfigure-in-table}\relax
\\
\sphinxbottomrule
\end{tabulary}
\sphinxtableafterendhook\par
\sphinxattableend\end{savenotes}

\sphinxAtStartPar
Con la llegada de \sphinxstylestrong{ALTAIR} en 1975, Microsoft desarrolla el lenguaje de programación BASIC. Este permitió programar en dicha computadora.

\sphinxAtStartPar
%
\begin{footnote}[12]\sphinxAtStartFootnote
Maturana, J. (2023, June 2). 25 aniversario de Computer Hoy: La historia de Windows, cómo ha sido la evolución del servicio más popular de Microsoft. Computer Hoy. https://computerhoy.20minutos.es/windows/historia\sphinxhyphen{}windows\sphinxhyphen{}como\sphinxhyphen{}ha\sphinxhyphen{}sido\sphinxhyphen{}evolucion\sphinxhyphen{}servicio\sphinxhyphen{}popular\sphinxhyphen{}microsoft\sphinxhyphen{}1237508
%
\end{footnote} La década de 1980 fue revolucionaria con la llegada del software comercial para PC, incluyendo procesadores de texto, hojas de cálculo y los primeros sistemas operativos con interfaz gráfica. Esto inicio el año 1981 con la creación del primer sistema operativo de Microsoft: \sphinxstylestrong{MS\sphinxhyphen{}DOS}.


\begin{savenotes}\sphinxattablestart
\sphinxthistablewithglobalstyle
\centering
\begin{tabulary}{\linewidth}[t]{T}
\sphinxtoprule
\sphinxtableatstartofbodyhook

\begin{sphinxfigure-in-table}
\centering
\capstart
\noindent\sphinxincludegraphics[width=200\sphinxpxdimen]{{image09}.png}
\sphinxfigcaption{Logo MSDOS}\label{\detokenize{cursos/informatica10/recursos/T01:id34}}\end{sphinxfigure-in-table}\relax
\\
\sphinxbottomrule
\end{tabulary}
\sphinxtableafterendhook\par
\sphinxattableend\end{savenotes}

\sphinxAtStartPar
\sphinxfootnotemark[11] En el año 1991 una comunidad de ingenieros y desarrolladores a nivel mundial crean el sistema operativo \sphinxstylestrong{LINUX}. El cual es de gran importancia hasta nuestros días.


\begin{savenotes}\sphinxattablestart
\sphinxthistablewithglobalstyle
\centering
\begin{tabulary}{\linewidth}[t]{T}
\sphinxtoprule
\sphinxtableatstartofbodyhook

\begin{sphinxfigure-in-table}
\centering
\capstart
\noindent\sphinxincludegraphics[width=400\sphinxpxdimen]{{image10}.png}
\sphinxfigcaption{Logo LINUX}\label{\detokenize{cursos/informatica10/recursos/T01:id35}}\end{sphinxfigure-in-table}\relax
\\
\sphinxbottomrule
\end{tabulary}
\sphinxtableafterendhook\par
\sphinxattableend\end{savenotes}

\sphinxAtStartPar
En la actualidad, el software ha evolucionado hacia aplicaciones en la nube, inteligencia artificial y desarrollo ágil, transformando completamente la manera en la que interactuamos con la tecnología.


\bigskip\hrule\bigskip


\sphinxstepscope


\chapter{Ofimática 10}
\label{\detokenize{cursos/ofimatica10/index:ofimatica-10}}\label{\detokenize{cursos/ofimatica10/index::doc}}
\sphinxAtStartPar
Bienvenido al curso de ofimática de grado décimo.

\sphinxAtStartPar
Acceda a las guías didácticas según el periodo.

\sphinxstepscope


\section{Periodo 01}
\label{\detokenize{cursos/ofimatica10/periodo01/index:periodo-01}}\label{\detokenize{cursos/ofimatica10/periodo01/index::doc}}\begin{quote}\begin{description}
\sphinxlineitem{DESEMPEÑO}
\sphinxAtStartPar
Conoce el concepto de ofimática, argumenta las ventajas y desventajas de su uso con software libre en especial con la suite LibreOffice.

\end{description}\end{quote}


\bigskip\hrule\bigskip



\subsection{Contenidos}
\label{\detokenize{cursos/ofimatica10/periodo01/index:contenidos}}
\sphinxstepscope


\subsubsection{Guía 01: Ofimática}
\label{\detokenize{cursos/ofimatica10/periodo01/G01:guia-01-ofimatica}}\label{\detokenize{cursos/ofimatica10/periodo01/G01::doc}}

\begin{savenotes}\sphinxattablestart
\sphinxthistablewithglobalstyle
\centering
\begin{tabular}[t]{\X{35}{100}\X{65}{100}}
\sphinxtoprule
\sphinxtableatstartofbodyhook

\begin{savenotes}\sphinxattablestart
\sphinxthistablewithglobalstyle
\centering
\begin{tabulary}{\linewidth}[t]{T}
\sphinxtoprule
\sphinxtableatstartofbodyhook
\sphinxAtStartPar
\sphinxincludegraphics{{escudo}.png}
\\
\sphinxbottomrule
\end{tabulary}
\sphinxtableafterendhook\par
\sphinxattableend\end{savenotes}
&

\begin{savenotes}\sphinxattablestart
\sphinxthistablewithglobalstyle
\centering
\begin{tabulary}{\linewidth}[t]{T}
\sphinxtoprule
\sphinxtableatstartofbodyhook
\sphinxAtStartPar
\sphinxstylestrong{INSTITUCIÓN EDUCATIVA}: John F. Kennedy
\\
\sphinxhline
\sphinxAtStartPar
\sphinxstylestrong{DOCENTE}: Julian Camilo Fonseca Romero
\\
\sphinxhline
\sphinxAtStartPar
\sphinxstylestrong{ASIGNATURA}: Ofimática
\\
\sphinxhline
\sphinxAtStartPar
\sphinxstylestrong{GRADO}: DÉCIMO
\\
\sphinxhline
\sphinxAtStartPar
\sphinxstylestrong{PERIODO:} 01
\\
\sphinxhline
\sphinxAtStartPar
\sphinxstylestrong{TIEMPO}: 10 semanas (4 horas / semana)
\\
\sphinxhline
\sphinxAtStartPar
\sphinxstylestrong{TIEMPO DE TRABAJO}: 10/40 horas
\\
\sphinxbottomrule
\end{tabulary}
\sphinxtableafterendhook\par
\sphinxattableend\end{savenotes}
\\
\sphinxbottomrule
\end{tabular}
\sphinxtableafterendhook\par
\sphinxattableend\end{savenotes}

\begin{sphinxuseclass}{sd-tab-set}
\begin{sphinxuseclass}{sd-tab-item}\subsubsection*{Objetivo}

\begin{sphinxuseclass}{sd-tab-content}\begin{itemize}
\item {} 
\sphinxAtStartPar
\sphinxstylestrong{Temática}: Concepto de ofimática.

\item {} 
\sphinxAtStartPar
\sphinxstylestrong{Objetivo de aprendizaje}: Entender que es la ofimática y conocer los principales proveedores de servicios ofimáticos, productividad y colaboración en la nube.

\end{itemize}

\end{sphinxuseclass}
\end{sphinxuseclass}
\begin{sphinxuseclass}{sd-tab-item}\subsubsection*{Orientaciones curriculares}

\begin{sphinxuseclass}{sd-tab-content}\begin{itemize}
\item {} 
\sphinxAtStartPar
\sphinxstylestrong{Componente}: Uso y apropiación de la Tecnología y la Informática.

\item {} 
\sphinxAtStartPar
\sphinxstylestrong{Competencia}: Genero propuestas innovadoras para el uso y aprovechamiento de productos tecnológicos.

\item {} 
\sphinxAtStartPar
\sphinxstylestrong{Evidencia de aprendizaje}: Utilizo adecuadamente herramientas informáticas para la búsqueda, organización, procesamiento, sistematización, comunicación y difusión de ideas.

\end{itemize}

\end{sphinxuseclass}
\end{sphinxuseclass}
\begin{sphinxuseclass}{sd-tab-item}\subsubsection*{Criterios de evaluación}

\begin{sphinxuseclass}{sd-tab-content}\begin{itemize}
\item {} 
\sphinxAtStartPar
Actividades desarrolladas en su totalidad.

\item {} 
\sphinxAtStartPar
Participación en clase.

\item {} 
\sphinxAtStartPar
Explicación clara de conceptos.

\item {} 
\sphinxAtStartPar
Argumentación de ideas.

\item {} 
\sphinxAtStartPar
Cumplimiento de las reglas del aula.

\end{itemize}

\end{sphinxuseclass}
\end{sphinxuseclass}
\begin{sphinxuseclass}{sd-tab-item}\subsubsection*{Recursos (1)}

\begin{sphinxuseclass}{sd-tab-content}\begin{itemize}
\item {} 
\sphinxAtStartPar
{\hyperref[\detokenize{cursos/ofimatica10/recursos/T01::doc}]{\sphinxcrossref{\DUrole{std}{\DUrole{std-doc}{TEXTO: Introducción a la Ofimática}}}}}

\end{itemize}

\end{sphinxuseclass}
\end{sphinxuseclass}
\end{sphinxuseclass}

\paragraph{Mensaje del autor}
\label{\detokenize{cursos/ofimatica10/periodo01/G01:mensaje-del-autor}}\begin{quote}
\begin{quote}

\sphinxAtStartPar
“En esta guía, el estudiante conoce las herramientas informáticas para organizar, procesar y difundir información de manera efectiva. Al entender la utilidad de suites como Google Workspace o Microsoft 365, se deduce cómo el trabajo colaborativo en la nube es esencial para las tareas de oficina o empresariales del siglo XXI.”
\end{quote}

\begin{flushright}
---Camilo Fonseca
\end{flushright}
\end{quote}


\bigskip\hrule\bigskip



\paragraph{A00: Introducción (+0,1)}
\label{\detokenize{cursos/ofimatica10/periodo01/G01:a00-introduccion-0-1}}
\sphinxAtStartPar
\DUrole{sd-sphinx-override}{\DUrole{sd-badge}{\DUrole{sd-bg-light}{\DUrole{sd-bg-text-light}{Exploración}}}}
\begin{enumerate}
\sphinxsetlistlabels{\arabic}{enumi}{enumii}{}{.}%
\item {} 
\sphinxAtStartPar
TRANSCRIBA en su cuaderno el objetivo, las orientaciones curriculares y los criterios de evaluación de la presente guía.

\item {} 
\sphinxAtStartPar
LEA el \sphinxstylestrong{mensaje del autor} y en su cuaderno responda:
\begin{itemize}
\item {} 
\sphinxAtStartPar
Imagine que tiene que coordinar un proyecto con un equipo de personas ubicadas en diferentes lugares ¿Qué herramientas, dispositivos o software podría usar para desarrollar este proyecto?

\end{itemize}

\end{enumerate}


\bigskip\hrule\bigskip



\paragraph{A01: Concepto de Ofimática (+0,1)}
\label{\detokenize{cursos/ofimatica10/periodo01/G01:a01-concepto-de-ofimatica-0-1}}
\sphinxAtStartPar
\DUrole{sd-sphinx-override}{\DUrole{sd-badge}{\DUrole{sd-bg-info}{\DUrole{sd-bg-text-info}{Estructuración}}}}
\begin{enumerate}
\sphinxsetlistlabels{\arabic}{enumi}{enumii}{}{.}%
\item {} 
\sphinxAtStartPar
En su cuaderno responda las siguientes preguntas:
\begin{itemize}
\item {} 
\sphinxAtStartPar
¿Qué es la ofimática?

\item {} 
\sphinxAtStartPar
¿Cuáles son las principales \sphinxstylestrong{áreas de trabajo} que cubre la ofimática?

\item {} 
\sphinxAtStartPar
¿Por qué la ofimática es clave para el entorno laboral moderno?

\end{itemize}

\end{enumerate}


\bigskip\hrule\bigskip



\paragraph{A02: Análisis de casos (+0,1)}
\label{\detokenize{cursos/ofimatica10/periodo01/G01:a02-analisis-de-casos-0-1}}
\sphinxAtStartPar
\DUrole{sd-sphinx-override}{\DUrole{sd-badge}{\DUrole{sd-bg-info}{\DUrole{sd-bg-text-info}{Estructuración}}}}
\begin{enumerate}
\sphinxsetlistlabels{\arabic}{enumi}{enumii}{}{.}%
\item {} 
\sphinxAtStartPar
Escriba en su cuaderno los siguientes casos y al final de cada uno, indique que \sphinxstylestrong{área de trabajo} de la ofimática corresponde y explique por que.
\begin{quote}

\sphinxAtStartPar
“Un grupo de estudiantes crea una carpeta compartida en la nube para subir evidencias fotográficas de una salida de campo, permitiendo que todos editen archivos de forma colaborativa y segura desde cualquier dispositivo.”
\end{quote}
\begin{quote}

\sphinxAtStartPar
“Un profesional envía una propuesta comercial a un cliente potencial, adjuntando documentos técnicos y utilizando firmas digitales personalizadas para mantener una comunicación formal y segura.”
\end{quote}
\begin{quote}

\sphinxAtStartPar
“Un líder de proyecto programa una reunión de equipo sincronizada, enviando invitaciones automáticas con recordatorios y verificando la disponibilidad de horario de todos los colaboradores en tiempo real.”
“El encargado de la biblioteca escolar crea un sistema estructurado para registrar el inventario de libros, los datos de los estudiantes y las fechas de devolución, facilitando la búsqueda rápida de ejemplares disponibles.”
\end{quote}
\begin{quote}

\sphinxAtStartPar
“Un estudiante de grado undécimo redacta su proyecto de grado utilizando herramientas de edición para estructurar el índice, insertar citas bibliográficas y exportar el informe final en formato PDF para su entrega formal.”
\end{quote}
\begin{quote}

\sphinxAtStartPar
“Un equipo de trabajo diseña una exposición visual sobre la cuarta revolución industrial, integrando imágenes, animaciones y transiciones para comunicar sus hallazgos de forma atractiva ante un comité evaluador.”
\end{quote}
\begin{quote}

\sphinxAtStartPar
“El tesorero de una pequeña empresa organiza los gastos mensuales mediante el uso de fórmulas automáticas para calcular el IVA, generar gráficos de barras que muestren las tendencias de consumo y filtrar datos por categorías de proveedores.”
\end{quote}

\end{enumerate}


\bigskip\hrule\bigskip



\paragraph{A03: Google Workspace (+0,1)}
\label{\detokenize{cursos/ofimatica10/periodo01/G01:a03-google-workspace-0-1}}
\sphinxAtStartPar
\DUrole{sd-sphinx-override}{\DUrole{sd-badge}{\DUrole{sd-bg-info}{\DUrole{sd-bg-text-info}{Estructuración}}}}
\begin{enumerate}
\sphinxsetlistlabels{\arabic}{enumi}{enumii}{}{.}%
\item {} 
\sphinxAtStartPar
En su cuaderno responda las siguientes preguntas:
\begin{itemize}
\item {} 
\sphinxAtStartPar
¿Qué es Google Workspace?

\item {} 
\sphinxAtStartPar
¿Qué herramientas incluye Google Workspace?

\item {} 
\sphinxAtStartPar
¿En qué consiste la suscripción de Google ONE?

\item {} 
\sphinxAtStartPar
¿Qué servicios recibe el usuario por suscribirse a \sphinxstylestrong{Google ONE \sphinxhyphen{} IA PLUS} y cuanto cuesta?

\end{itemize}

\item {} 
\sphinxAtStartPar
En Internet, consulte que otros planes de \sphinxstylestrong{Google ONE} existen y por cada plan en su cuaderno responda:
\begin{itemize}
\item {} 
\sphinxAtStartPar
¿Cuanto cuesta el plan?

\item {} 
\sphinxAtStartPar
¿Que servicios carece dicho plan respecto a \sphinxstylestrong{Google ONE \sphinxhyphen{} IA PLUS}?

\item {} 
\sphinxAtStartPar
¿Que servicios extra ofrece dicho plan respecto a \sphinxstylestrong{Google ONE \sphinxhyphen{} IA PLUS}?

\end{itemize}

\end{enumerate}


\bigskip\hrule\bigskip



\paragraph{A04: Microsoft Office (+0,1)}
\label{\detokenize{cursos/ofimatica10/periodo01/G01:a04-microsoft-office-0-1}}
\sphinxAtStartPar
\DUrole{sd-sphinx-override}{\DUrole{sd-badge}{\DUrole{sd-bg-info}{\DUrole{sd-bg-text-info}{Estructuración}}}}
\begin{enumerate}
\sphinxsetlistlabels{\arabic}{enumi}{enumii}{}{.}%
\item {} 
\sphinxAtStartPar
En su cuaderno responda las siguientes preguntas:
\begin{itemize}
\item {} 
\sphinxAtStartPar
¿Qué es Microsoft Office?

\item {} 
\sphinxAtStartPar
¿Cúales son las aplicaciones más conocidas de Microsoft Office?

\item {} 
\sphinxAtStartPar
¿Qué sistemas operativos son compatibles con Microsoft Office?

\item {} 
\sphinxAtStartPar
¿En qué consiste la suscripción \sphinxstylestrong{Microsoft 365 Family}?

\item {} 
\sphinxAtStartPar
¿Qué servicios recibe el usuario por suscribirse a \sphinxstylestrong{Microsoft 365 Family} y cuanto cuesta?

\end{itemize}

\item {} 
\sphinxAtStartPar
En Internet, consulte que otros planes de Microsoft 365 existen y por cada plan en su cuaderno responda:
\begin{itemize}
\item {} 
\sphinxAtStartPar
¿Cuanto cuesta dicho plan?

\item {} 
\sphinxAtStartPar
¿Que servicios carece dicho plan respecto a \sphinxstylestrong{Microsoft Family 365}?

\item {} 
\sphinxAtStartPar
¿Que servicios extra ofrece dicho plan respecto a \sphinxstylestrong{Microsoft Family 365}?

\end{itemize}

\end{enumerate}


\bigskip\hrule\bigskip



\paragraph{A05: Servicio indicado (+0,5)}
\label{\detokenize{cursos/ofimatica10/periodo01/G01:a05-servicio-indicado-0-5}}
\sphinxAtStartPar
\DUrole{sd-sphinx-override}{\DUrole{sd-badge}{\DUrole{sd-bg-info}{\DUrole{sd-bg-text-info}{Estructuración}}}}
\begin{enumerate}
\sphinxsetlistlabels{\arabic}{enumi}{enumii}{}{.}%
\item {} 
\sphinxAtStartPar
Para cada uno de los casos, seleccione el servicio que más convenga. No olvide explicar el por qué. Los servicios son:
\begin{itemize}
\item {} 
\sphinxAtStartPar
Microsoft 365 Personal.

\item {} 
\sphinxAtStartPar
Microsoft 365 Familia.

\item {} 
\sphinxAtStartPar
Microsoft 365 Premium.

\item {} 
\sphinxAtStartPar
Office Hogar 2024.

\item {} 
\sphinxAtStartPar
Plan gratuito de Google One.

\item {} 
\sphinxAtStartPar
Plan gratuito de Google.

\item {} 
\sphinxAtStartPar
Planes de suscripción de Google One.

\item {} 
\sphinxAtStartPar
Google One con Google AI.

\item {} 
\sphinxAtStartPar
Uso de la aplicación Google One.

\item {} 
\sphinxAtStartPar
Planes familiares de Google One.

\end{itemize}

\sphinxAtStartPar
Los casos son:
\begin{quote}

\sphinxAtStartPar
\sphinxstylestrong{CASO 01:} Sofía es una diseñadora digital que busca optimizar su tiempo usando herramientas avanzadas de Inteligencia Artificial. Necesita acceso exclusivo a las funciones más recientes de Copilot y un uso extensivo para crear y editar imágenes con IA.
\end{quote}
\begin{quote}

\sphinxAtStartPar
\sphinxstylestrong{CASO 02:} Un usuario prefiere gestionar todas sus suscripciones y copias de seguridad directamente desde su smartphone. Necesita una forma sencilla de monitorear cuánto espacio ocupan sus fotos y archivos en tiempo real. El servicio debe suministrar una herramienta de administración móvil centralizada.
\end{quote}
\begin{quote}

\sphinxAtStartPar
\sphinxstylestrong{CASO 03:} Elena es una estudiante interesada en las últimas tendencias tecnológicas. Quiere integrar funciones avanzadas de IA directamente en sus herramientas de trabajo de Google. El servicio debe permitir acceder a los beneficios exclusivos de Google AI.
\end{quote}
\begin{quote}

\sphinxAtStartPar
\sphinxstylestrong{CASO 04:} Carlos estudia ingeniería y necesita usar Word, Excel y PowerPoint tanto en su PC como en su celular. Vive solo y tiene un presupuesto ajustado, pero requiere al menos 1 TB de espacio para sus proyectos pesados.
\end{quote}
\begin{quote}

\sphinxAtStartPar
\sphinxstylestrong{CASO 05:} Lucía utiliza su cuenta de Google principalmente para enviar correos electrónicos académicos y guardar copias de seguridad de las fotos de su celular. No suele manejar archivos muy pesados. Se necesita un almacenamiento básico sin coste adicional.
\end{quote}
\begin{quote}

\sphinxAtStartPar
\sphinxstylestrong{CASO 06:} Jorge utiliza su correo electrónico para comunicaciones básicas y ocasionalmente guarda documentos PDF y fotos personales. No utiliza aplicaciones de escritorio complejas y sus archivos totales no superan los 12 GB. Un almacenamiento gratuito sin coste adicional.
\end{quote}
\begin{quote}

\sphinxAtStartPar
\sphinxstylestrong{CASO 07:} Un pequeño grupo de estudio quiere una solución donde puedan compartir un fondo común de almacenamiento para sus recursos compartidos, pero manteniendo cada uno su privacidad. El servicio debe permitir compartir beneficios con otros miembros.
\end{quote}
\begin{quote}

\sphinxAtStartPar
\sphinxstylestrong{CASO 08:} Martín ha alcanzado el límite de sus 15 GB gratuitos debido a la gran cantidad de videos y trabajos de investigación guardados en su unidad. Necesita un poco más de espacio para no dejar de recibir correos importantes. Se debe adquirir una ampliación de almacenamiento económica para seguir usando su cuenta personal
\end{quote}
\begin{quote}

\sphinxAtStartPar
\sphinxstylestrong{CASO 09:} Una pequeña oficina de contabilidad prefiere no tener gastos recurrentes mensuales o anuales. Solo necesitan instalar las versiones clásicas de Word y Excel en una computadora de la oficina y no requieren almacenamiento en la nube ni actualizaciones constantes de funciones. No se requiere una suscripción.
\end{quote}
\begin{quote}

\sphinxAtStartPar
\sphinxstylestrong{CASO 10:} La familia Martínez tiene 5 integrantes (padres y 3 hijos en el colegio). Todos necesitan las aplicaciones de Office para sus tareas y quieren que cada uno tenga su propio espacio privado de almacenamiento de 1 TB sin pagar suscripciones por separado.
\end{quote}

\end{enumerate}


\bigskip\hrule\bigskip



\paragraph{A06: Examen escrito (+4,0)}
\label{\detokenize{cursos/ofimatica10/periodo01/G01:a06-examen-escrito-4-0}}
\sphinxAtStartPar
\DUrole{sd-sphinx-override}{\DUrole{sd-badge}{\DUrole{sd-bg-danger}{\DUrole{sd-bg-text-danger}{Transferencia}}}}
\begin{quote}\begin{description}
\sphinxlineitem{TIEMPO}
\sphinxAtStartPar
30 minutos

\sphinxlineitem{PREGUNTAS}
\sphinxAtStartPar
10

\end{description}\end{quote}

\begin{sphinxadmonition}{note}{Posibles preguntas}
\begin{enumerate}
\sphinxsetlistlabels{\arabic}{enumi}{enumii}{}{.}%
\item {} 
\sphinxAtStartPar
Explique con sus propias palabras qué es la \sphinxstylestrong{ofimática} y de qué combinación de términos surge su nombre.

\item {} 
\sphinxAtStartPar
¿Cuál es el propósito fundamental de utilizar \sphinxstylestrong{herramientas ofimáticas} en un entorno de trabajo o estudio?

\item {} 
\sphinxAtStartPar
Mencione al menos cuatro áreas principales que abarca la ofimática.

\item {} 
\sphinxAtStartPar
Describa la función principal de las herramientas de \sphinxstylestrong{procesamiento de textos} e indique un ejemplo de software tanto de Google como de Microsoft.

\item {} 
\sphinxAtStartPar
¿Para qué tipo de tareas es ideal el uso de una \sphinxstylestrong{hoja de cálculo} y cuáles son sus componentes básicos?

\item {} 
\sphinxAtStartPar
Explique la utilidad de aplicaciones como \sphinxstylestrong{PowerPoint} o \sphinxstylestrong{Google Slides} en el ámbito académico o empresarial.

\item {} 
\sphinxAtStartPar
¿En qué consiste la gestión de \sphinxstylestrong{bases de datos} y qué aplicación de \sphinxstylestrong{Microsoft} se menciona para esta tarea?

\item {} 
\sphinxAtStartPar
¿En qué consiste \sphinxstylestrong{Google Workspace} y qué beneficio principal tienen los usuarios cuando necesitan trabajar en equipo?

\item {} 
\sphinxAtStartPar
Describa la función de \sphinxstylestrong{Google Drive} dentro del ecosistema de Google Workspace.

\item {} 
\sphinxAtStartPar
Según el plan \sphinxstylestrong{Google One AI PLUS}, ¿Qué servicios adicionales recibe un usuario al activar esta suscripción específica?

\item {} 
\sphinxAtStartPar
¿Qué es \sphinxstylestrong{Gemini} y bajo qué plan de suscripción de Google se obtienen sus mejores funciones?

\item {} 
\sphinxAtStartPar
¿Qué es \sphinxstylestrong{Microsoft Office} y para qué tipos de dispositivos está disponible?

\item {} 
\sphinxAtStartPar
¿Cuales son las diferencias entre las funciones principales de un \sphinxstylestrong{procesador de texto} y una \sphinxstylestrong{hoja de cálculo}?

\item {} 
\sphinxAtStartPar
¿Qué funciones integra \sphinxstylestrong{Outlook} además de la gestión de \sphinxstylestrong{correos electrónicos}?

\item {} 
\sphinxAtStartPar
Explique tres beneficios clave de contar con la \sphinxstylestrong{suscripción familiar de Microsoft 365}.

\item {} 
\sphinxAtStartPar
Compare la capacidad de almacenamiento y los precios que ofrece el plan \sphinxstylestrong{Google One AI PLUS} frente a lo que recibe cada persona en el plan \sphinxstylestrong{Microsoft 365 Family}.

\item {} 
\sphinxAtStartPar
¿Por qué se considera que las suites en línea permiten “conectar, crear y colaborar de manera segura”?

\item {} 
\sphinxAtStartPar
Según las características de \sphinxstylestrong{Microsoft 365}, ¿qué ventajas ofrece en cuanto a seguridad y acceso a funciones en comparación con versiones estáticas?

\item {} 
\sphinxAtStartPar
¿Por qué la \sphinxstylestrong{ofimática} es considerada fundamental en el \sphinxstylestrong{entorno laboral moderno}?

\item {} 
\sphinxAtStartPar
¿Qué herramientas específicas se mencionan para la gestión de calendarios, agendas y documentos en ambas suites (Google y Microsoft)?

\end{enumerate}
\end{sphinxadmonition}

\sphinxstepscope


\subsubsection{Guía 02: Software Libre}
\label{\detokenize{cursos/ofimatica10/periodo01/G02:guia-02-software-libre}}\label{\detokenize{cursos/ofimatica10/periodo01/G02::doc}}

\begin{savenotes}\sphinxattablestart
\sphinxthistablewithglobalstyle
\centering
\begin{tabular}[t]{\X{35}{100}\X{65}{100}}
\sphinxtoprule
\sphinxtableatstartofbodyhook

\begin{savenotes}\sphinxattablestart
\sphinxthistablewithglobalstyle
\centering
\begin{tabulary}{\linewidth}[t]{T}
\sphinxtoprule
\sphinxtableatstartofbodyhook
\sphinxAtStartPar
\sphinxincludegraphics{{escudo}.png}
\\
\sphinxbottomrule
\end{tabulary}
\sphinxtableafterendhook\par
\sphinxattableend\end{savenotes}
&

\begin{savenotes}\sphinxattablestart
\sphinxthistablewithglobalstyle
\centering
\begin{tabulary}{\linewidth}[t]{T}
\sphinxtoprule
\sphinxtableatstartofbodyhook
\sphinxAtStartPar
\sphinxstylestrong{INSTITUCIÓN EDUCATIVA}: John F. Kennedy
\\
\sphinxhline
\sphinxAtStartPar
\sphinxstylestrong{DOCENTE}: Julian Camilo Fonseca Romero
\\
\sphinxhline
\sphinxAtStartPar
\sphinxstylestrong{ASIGNATURA}: Ofimática
\\
\sphinxhline
\sphinxAtStartPar
\sphinxstylestrong{GRADO}: DÉCIMO
\\
\sphinxhline
\sphinxAtStartPar
\sphinxstylestrong{PERIODO:} 01
\\
\sphinxhline
\sphinxAtStartPar
\sphinxstylestrong{TIEMPO}: 10 semanas (4 horas / semana)
\\
\sphinxhline
\sphinxAtStartPar
\sphinxstylestrong{TIEMPO DE TRABAJO}: 20/40 horas
\\
\sphinxbottomrule
\end{tabulary}
\sphinxtableafterendhook\par
\sphinxattableend\end{savenotes}
\\
\sphinxbottomrule
\end{tabular}
\sphinxtableafterendhook\par
\sphinxattableend\end{savenotes}

\begin{sphinxuseclass}{sd-tab-set}
\begin{sphinxuseclass}{sd-tab-item}\subsubsection*{Objetivo}

\begin{sphinxuseclass}{sd-tab-content}\begin{itemize}
\item {} 
\sphinxAtStartPar
\sphinxstylestrong{Temática}: Software libre.

\item {} 
\sphinxAtStartPar
\sphinxstylestrong{Objetivo de aprendizaje}: Entender la importancia del desarrollo y uso del software libre.

\end{itemize}

\end{sphinxuseclass}
\end{sphinxuseclass}
\begin{sphinxuseclass}{sd-tab-item}\subsubsection*{Orientaciones curriculares}

\begin{sphinxuseclass}{sd-tab-content}\begin{itemize}
\item {} 
\sphinxAtStartPar
\sphinxstylestrong{Componente}: Tecnología, Informática y Sociedad.

\item {} 
\sphinxAtStartPar
\sphinxstylestrong{Competencia}: Actúo críticamente y de forma argumentada frente a las implicaciones éticas, sociales y ambientales del desarrollo, implementación, uso y disposición final de los productos tecnológicos.

\item {} 
\sphinxAtStartPar
\sphinxstylestrong{Evidencia de aprendizaje}: Juzgo las implicaciones que la protección a la propiedad intelectual tiene sobre el desarrollo y uso de diversas manifestaciones tecnológicas en el mundo.

\end{itemize}

\end{sphinxuseclass}
\end{sphinxuseclass}
\begin{sphinxuseclass}{sd-tab-item}\subsubsection*{Criterios de evaluación}

\begin{sphinxuseclass}{sd-tab-content}\begin{itemize}
\item {} 
\sphinxAtStartPar
Actividades desarrolladas en su totalidad.

\item {} 
\sphinxAtStartPar
Participación en clase.

\item {} 
\sphinxAtStartPar
Explicación clara de conceptos.

\item {} 
\sphinxAtStartPar
Argumentación de ideas.

\item {} 
\sphinxAtStartPar
Cumplimiento de las reglas del aula.

\end{itemize}

\end{sphinxuseclass}
\end{sphinxuseclass}
\begin{sphinxuseclass}{sd-tab-item}\subsubsection*{Recursos (1)}

\begin{sphinxuseclass}{sd-tab-content}\begin{itemize}
\item {} 
\sphinxAtStartPar
{\hyperref[\detokenize{cursos/ofimatica10/recursos/T02::doc}]{\sphinxcrossref{\DUrole{std}{\DUrole{std-doc}{TEXTO: Software libre}}}}}

\end{itemize}

\end{sphinxuseclass}
\end{sphinxuseclass}
\end{sphinxuseclass}

\paragraph{Mensaje del autor}
\label{\detokenize{cursos/ofimatica10/periodo01/G02:mensaje-del-autor}}\begin{quote}
\begin{quote}

\sphinxAtStartPar
“El estudio del software libre es esencial ya que, como manifestación tecnológica, surge como una alternativa ética y política a los modelos de propiedad intelectual. Esta filosofía no solo garantiza la equidad en el acceso a recursos digitales, sino que permite trascender del saber tecnológico al conocimiento tecnológico profundo. Al permitir el análisis de su funcionamiento interno (posibilidad negada por el software propietario), se faculta al estudiante para adoptar, adaptar y generar soluciones que respondan a necesidades sociales.”
\end{quote}

\begin{flushright}
---Camilo Fonseca
\end{flushright}
\end{quote}


\bigskip\hrule\bigskip



\paragraph{A00: Introducción (+0,1)}
\label{\detokenize{cursos/ofimatica10/periodo01/G02:a00-introduccion-0-1}}
\sphinxAtStartPar
\DUrole{sd-sphinx-override}{\DUrole{sd-badge}{\DUrole{sd-bg-light}{\DUrole{sd-bg-text-light}{Exploración}}}}
\begin{enumerate}
\sphinxsetlistlabels{\arabic}{enumi}{enumii}{}{.}%
\item {} 
\sphinxAtStartPar
TRANSCRIBA en su cuaderno el objetivo, las orientaciones curriculares y los criterios de evaluación de la presente guía.

\item {} 
\sphinxAtStartPar
LEA el mensaje del autor luego identifique y describa en su cuaderno \sphinxstylestrong{5 ejemplos de software} gratuitos que utilice, reflexionando sobre cómo estos facilitan sus actividades cotidianas.

\end{enumerate}


\bigskip\hrule\bigskip



\paragraph{A01: Definición de software (+0,1)}
\label{\detokenize{cursos/ofimatica10/periodo01/G02:a01-definicion-de-software-0-1}}
\sphinxAtStartPar
\DUrole{sd-sphinx-override}{\DUrole{sd-badge}{\DUrole{sd-bg-info}{\DUrole{sd-bg-text-info}{Estructuración}}}}
\begin{enumerate}
\sphinxsetlistlabels{\arabic}{enumi}{enumii}{}{.}%
\item {} 
\sphinxAtStartPar
En su cuaderno responda:
\begin{itemize}
\item {} 
\sphinxAtStartPar
¿Qué es el software?

\item {} 
\sphinxAtStartPar
¿Cuales son los tipos de software?

\item {} 
\sphinxAtStartPar
¿Cómo se puede clasificar el software?

\end{itemize}

\item {} 
\sphinxAtStartPar
En su cuaderno, escriba una lista de 15 ejemplos de software que usted conozca (de smartphone, de computadores, de internet, etc.) NO OLVIDE HACER UNA BREVE DESCRIPCIÓN.

\item {} 
\sphinxAtStartPar
Los anteriores ejemplos identifique si es \sphinxstylestrong{software de sistema} o \sphinxstylestrong{software de aplicación} y explique el porqué.

\item {} 
\sphinxAtStartPar
Los anteriores ejemplos identifique si es \sphinxstylestrong{software libre}, \sphinxstylestrong{software propietario} o \sphinxstylestrong{software pirata}.

\end{enumerate}


\bigskip\hrule\bigskip



\paragraph{A02: Importancia del software libre (+0,1)}
\label{\detokenize{cursos/ofimatica10/periodo01/G02:a02-importancia-del-software-libre-0-1}}
\sphinxAtStartPar
\DUrole{sd-sphinx-override}{\DUrole{sd-badge}{\DUrole{sd-bg-info}{\DUrole{sd-bg-text-info}{Estructuración}}}}
\begin{enumerate}
\sphinxsetlistlabels{\arabic}{enumi}{enumii}{}{.}%
\item {} 
\sphinxAtStartPar
En su cuaderno responda:
\begin{itemize}
\item {} 
\sphinxAtStartPar
¿Qué es el software libre?

\item {} 
\sphinxAtStartPar
¿Cuáles son las \sphinxstylestrong{4 libertades esenciales} que debe cumplir un software, (para que sea considerado software libre)?

\end{itemize}

\item {} 
\sphinxAtStartPar
Analice los siguientes casos y en su cuaderno explique si el software es libre o no. Identifique si cumple las \sphinxstylestrong{4 libertades esenciales} y explique el por qué:
\begin{quote}

\sphinxAtStartPar
“La institución educativa utiliza el \sphinxstylestrong{software X} para gestionar sus datos. Al adquirirlo, la comunidad recibió no solo el instalador, sino también el código fuente del programa. Debido a una necesidad específica, el docente de tecnología pudo estudiar cómo funciona el software y modificarlo para que acepte nuevos parámetros de entrada que antes no permitía. Posteriormente, el docente distribuyó copias de esta versión mejorada a otros colegios de la región para que ellos también pudieran beneficiarse de los cambios sin restricciones legales.”
\end{quote}
\begin{quote}

\sphinxAtStartPar
“El \sphinxstylestrong{Software Y} se distribuye bajo un modelo de suscripción de pago obligatorio que requiere una cuenta personal para su activación. Unos usuarios en Rusia lo utilizan para capacitación, pero solo pueden usarlo en un máximo de cinco dispositivos simultáneamente según su plan. Cuando usuarios en España intentaron obtener la versión modificada de Rusia, descubrieron que la licencia de uso prohíbe explícitamente cualquier redistribución o copia del programa a terceros. Además, aunque los usuarios son dueños de los datos que generan, no tienen acceso al código fuente ni el control sobre el software, por lo que no pueden estudiar cómo funciona ni realizar mejoras legales para adaptarlo a sus necesidades.”
\end{quote}
\begin{quote}

\sphinxAtStartPar
“Al \sphinxstylestrong{Software Z} se le implementaron diversas adaptaciones gracias a que la comunidad tuvo acceso total a su código fuente, permitiendo que todas las modificaciones funcionaran sin bloqueos legales. Aunque inicialmente los clientes trabajaban gastando más tiempo del habitual, la libertad de examinar el software permitió a los programadores desentrañar su funcionamiento para adecuar y optimizar las funcionalidades críticas. Al final, estas mejoras se compartieron legalmente, permitiendo que usuarios en España lo eligieran para reemplazar un software privativo que, a diferencia del Software Z, prohibía por contrato su libre distribución y estudio.”
\end{quote}

\end{enumerate}


\bigskip\hrule\bigskip



\paragraph{A03: Ejemplos de software libre (+0,1)}
\label{\detokenize{cursos/ofimatica10/periodo01/G02:a03-ejemplos-de-software-libre-0-1}}
\sphinxAtStartPar
\DUrole{sd-sphinx-override}{\DUrole{sd-badge}{\DUrole{sd-bg-info}{\DUrole{sd-bg-text-info}{Estructuración}}}}
\begin{enumerate}
\sphinxsetlistlabels{\arabic}{enumi}{enumii}{}{.}%
\item {} 
\sphinxAtStartPar
A partir de los ejemplos de software libre, en su cuaderno diligencie el siguiente crucigrama:

\end{enumerate}


\begin{savenotes}\sphinxattablestart
\sphinxthistablewithglobalstyle
\centering
\begin{tabulary}{\linewidth}[t]{T}
\sphinxtoprule
\sphinxtableatstartofbodyhook

\begin{sphinxfigure-in-table}
\centering
\capstart
\noindent\sphinxincludegraphics[width=500\sphinxpxdimen]{{image013}.png}
\sphinxfigcaption{\sphinxstylestrong{Crucigrama} ejemplos software libre}\label{\detokenize{cursos/ofimatica10/periodo01/G02:id1}}\end{sphinxfigure-in-table}\relax
\\
\sphinxhline
\sphinxAtStartPar
\sphinxstylestrong{2.} No requiere Codecs para funcionar.
\\
\sphinxhline
\sphinxAtStartPar
\sphinxstylestrong{6.} Se usa para crear logotipos.
\\
\sphinxhline
\sphinxAtStartPar
\sphinxstylestrong{5.} Mismas funcionalidades de Photoshop pero es software libre.
\\
\sphinxhline
\sphinxAtStartPar
\sphinxstylestrong{1.} Mismas funcionalidades de Microsoft Office pero es software libre.
\\
\sphinxhline
\sphinxAtStartPar
\sphinxstylestrong{4.} Es una alternativa libre al navegador Google CHROME.
\\
\sphinxhline
\sphinxAtStartPar
\sphinxstylestrong{3.} Este sistema fue creado inicialmente por Linus Torvalds.
\\
\sphinxbottomrule
\end{tabulary}
\sphinxtableafterendhook\par
\sphinxattableend\end{savenotes}
\begin{enumerate}
\sphinxsetlistlabels{\arabic}{enumi}{enumii}{}{.}%
\setcounter{enumi}{1}
\item {} 
\sphinxAtStartPar
En su cuaderno, REALICE una \sphinxstylestrong{descripción} de cada uno de los ejemplos de software libre, encontrados en el anterior crucigrama.

\end{enumerate}


\bigskip\hrule\bigskip



\paragraph{A04: Relación entre software libre, software propietario y software pirata (+0,1)}
\label{\detokenize{cursos/ofimatica10/periodo01/G02:a04-relacion-entre-software-libre-software-propietario-y-software-pirata-0-1}}
\sphinxAtStartPar
\DUrole{sd-sphinx-override}{\DUrole{sd-badge}{\DUrole{sd-bg-info}{\DUrole{sd-bg-text-info}{Estructuración}}}}
\begin{enumerate}
\sphinxsetlistlabels{\arabic}{enumi}{enumii}{}{.}%
\item {} 
\sphinxAtStartPar
Lea las \sphinxstylestrong{características} del software libre, software propietario y el software pirata. Luego complete el siguiente cuadro en su cuaderno:

\end{enumerate}


\begin{savenotes}\sphinxattablestart
\sphinxthistablewithglobalstyle
\centering
\begin{tabulary}{\linewidth}[t]{TTT}
\sphinxtoprule
\sphinxstyletheadfamily 
\sphinxAtStartPar
SOFTWARE LIBRE
&\sphinxstyletheadfamily 
\sphinxAtStartPar
SOFTWARE PROPIETARIO
&\sphinxstyletheadfamily 
\sphinxAtStartPar
SOFTWARE PIRATA
\\
\sphinxmidrule
\sphinxtableatstartofbodyhook
\sphinxAtStartPar
Impulsado por la \sphinxstylestrong{free Software Foundation}
&
\sphinxAtStartPar
Impulsado por la \sphinxstylestrong{propia empresa} con fines lucrativos
&
\sphinxAtStartPar
Impulsado por \sphinxstylestrong{personas} u \sphinxstylestrong{organizaciones} con fines maliciosos
\\
\sphinxhline
\sphinxAtStartPar
.
&
\sphinxAtStartPar

&
\sphinxAtStartPar

\\
\sphinxhline
\sphinxAtStartPar
.
&
\sphinxAtStartPar

&
\sphinxAtStartPar

\\
\sphinxhline
\sphinxAtStartPar
.
&
\sphinxAtStartPar

&
\sphinxAtStartPar

\\
\sphinxhline
\sphinxAtStartPar
.
&
\sphinxAtStartPar

&
\sphinxAtStartPar

\\
\sphinxbottomrule
\end{tabulary}
\sphinxtableafterendhook\par
\sphinxattableend\end{savenotes}


\bigskip\hrule\bigskip



\paragraph{A05: Ahorro(\$) usando software libre (+0,2)}
\label{\detokenize{cursos/ofimatica10/periodo01/G02:a05-ahorro-usando-software-libre-0-2}}
\sphinxAtStartPar
\DUrole{sd-sphinx-override}{\DUrole{sd-badge}{\DUrole{sd-bg-info}{\DUrole{sd-bg-text-info}{Estructuración}}}}

\sphinxAtStartPar
Desarrolle los siguientes cálculos en su cuaderno:
\begin{enumerate}
\sphinxsetlistlabels{\arabic}{enumi}{enumii}{}{.}%
\item {} 
\sphinxAtStartPar
Una empresa necesita adquirir 11 licencias con \sphinxstylestrong{Windows 11} para poder instalar en 11 equipos. El presupuesto asignado fue \$10,000,000. Al final se dieron cuenta que solo 5 equipos cumplían con los requisitos para poder instalar Windows 11. Después de hacer cuentas, ¿Cuánta plata le sobro a la empresa?
\begin{quote}

\sphinxAtStartPar
\sphinxstylestrong{RESPUESTA:} \$6,150,000
\end{quote}

\item {} 
\sphinxAtStartPar
Una familia quiere hacer un cálculo para identificar cuanto tienen que invertir en 5 años por tener activa una suscripción a \sphinxstylestrong{Microsoft Office}. Cada año la suscripción aumenta un 3\%.
\begin{quote}

\sphinxAtStartPar
\sphinxstylestrong{RESPUESTA:} \$2,427,122.84
\end{quote}

\item {} 
\sphinxAtStartPar
Un diseñador adquirió una suscripción de Adobe Photoshop por un año en diferentes intervalos. Uso Photoshop de ENERO a ABRIL. De MAYO a AGOSTO no uso photoshop porque estaba en vacaciones. De SEPTIEMBRE a NOVIEMBRE Pago nuevamente. En Diciembre la renovó por un año más. ¿Cuánto gasto el diseñador por usar la licencia en el primer año?
\begin{quote}

\sphinxAtStartPar
\sphinxstylestrong{RESPUESTA:} \$391,768
\end{quote}

\item {} 
\sphinxAtStartPar
Un profesional independiente decide contratar la suscripción \sphinxstylestrong{Google ONE \sphinxhyphen{} AI PLUS} para potenciar su trabajo con inteligencia artificial. Si planea utilizar el servicio de manera ininterrumpida durante un año y 3/4, ¿cuál será el valor total de su inversión?
\begin{quote}

\sphinxAtStartPar
\sphinxstylestrong{RESPUESTA:} \$358,200
\end{quote}

\item {} 
\sphinxAtStartPar
Un grupo de 4 amigos desea adquirir una suscripción de \sphinxstylestrong{Microsoft 365 Family} para compartir los gastos de forma equitativa durante un año. Si cada uno aporta la misma cantidad, ¿cuánto dinero debe pagar cada integrante para cubrir el costo total de la licencia anual?
\begin{quote}

\sphinxAtStartPar
\sphinxstylestrong{RESPUESTA:} \$114,999
\end{quote}

\end{enumerate}


\bigskip\hrule\bigskip



\paragraph{A06: El riesgo del software pirata (+0,2)}
\label{\detokenize{cursos/ofimatica10/periodo01/G02:a06-el-riesgo-del-software-pirata-0-2}}
\sphinxAtStartPar
\DUrole{sd-sphinx-override}{\DUrole{sd-badge}{\DUrole{sd-bg-info}{\DUrole{sd-bg-text-info}{Estructuración}}}}
\begin{enumerate}
\sphinxsetlistlabels{\arabic}{enumi}{enumii}{}{.}%
\item {} 
\sphinxAtStartPar
Responda en su cuaderno:
\begin{itemize}
\item {} 
\sphinxAtStartPar
¿Cuáles son las características de un software pirata?

\item {} 
\sphinxAtStartPar
¿Qué consecuencias tiene un sistema que tenga VIRUS y TROYANOS?

\end{itemize}

\item {} 
\sphinxAtStartPar
LEA el siguiente caso:
\begin{quote}

\sphinxAtStartPar
Una empresa está afectada por el uso de \sphinxstylestrong{software pirata}.
\begin{itemize}
\item {} 
\sphinxAtStartPar
Una auditoria del gobierno multó a la empresa por \$100,000,000.00 COP.

\item {} 
\sphinxAtStartPar
La información sensible como fórmulas secretas o datos personales de los empleados está en riesgo de ser expuesta.

\item {} 
\sphinxAtStartPar
Los PCs de la empresa están lentos y se reinician solos. Esto altera el flujo normal de trabajo.

\end{itemize}
\end{quote}

\sphinxAtStartPar
En una hoja de cuaderno responda:
\begin{itemize}
\item {} 
\sphinxAtStartPar
¿Qué medidas tomaría usted para evitar que la empresa quiebre?

\end{itemize}

\end{enumerate}


\bigskip\hrule\bigskip



\paragraph{A07: Identificación de software propietario y software libre (+0,1)}
\label{\detokenize{cursos/ofimatica10/periodo01/G02:a07-identificacion-de-software-propietario-y-software-libre-0-1}}
\sphinxAtStartPar
\DUrole{sd-sphinx-override}{\DUrole{sd-badge}{\DUrole{sd-bg-info}{\DUrole{sd-bg-text-info}{Estructuración}}}}
\begin{enumerate}
\sphinxsetlistlabels{\arabic}{enumi}{enumii}{}{.}%
\item {} 
\sphinxAtStartPar
En una hoja de cuaderno, complete la siguiente tabla:

\end{enumerate}


\begin{savenotes}\sphinxattablestart
\sphinxthistablewithglobalstyle
\centering
\begin{tabulary}{\linewidth}[t]{TTTT}
\sphinxtoprule
\sphinxstyletheadfamily 
\sphinxAtStartPar
.
&\sphinxstyletheadfamily 
\sphinxAtStartPar
SOFTWARE LIBRE
&\sphinxstyletheadfamily 
\sphinxAtStartPar
SOFTWARE PROPIETARIO
&\sphinxstyletheadfamily 
\sphinxAtStartPar
Descripción
\\
\sphinxmidrule
\sphinxtableatstartofbodyhook
\sphinxAtStartPar
\sphinxstylestrong{Edición de fotos}
&
\sphinxAtStartPar
GIMP
&
\sphinxAtStartPar
PHOTOSHOP
&
\sphinxAtStartPar

\\
\sphinxhline
\sphinxAtStartPar
\sphinxstylestrong{Procesamiento de textos}
&
\sphinxAtStartPar

&
\sphinxAtStartPar

&
\sphinxAtStartPar

\\
\sphinxhline
\sphinxAtStartPar
\sphinxstylestrong{Sistema Operativo}
&
\sphinxAtStartPar

&
\sphinxAtStartPar

&
\sphinxAtStartPar

\\
\sphinxhline
\sphinxAtStartPar
\sphinxstylestrong{Hojas de cálculo}
&
\sphinxAtStartPar

&
\sphinxAtStartPar

&
\sphinxAtStartPar

\\
\sphinxhline
\sphinxAtStartPar
\sphinxstylestrong{Navegadores Web}
&
\sphinxAtStartPar

&
\sphinxAtStartPar

&
\sphinxAtStartPar

\\
\sphinxhline
\sphinxAtStartPar
\sphinxstylestrong{Presentaciones}
&
\sphinxAtStartPar

&
\sphinxAtStartPar

&
\sphinxAtStartPar

\\
\sphinxhline
\sphinxAtStartPar
\sphinxstylestrong{Reproductor Multimedia}
&
\sphinxAtStartPar

&
\sphinxAtStartPar

&
\sphinxAtStartPar

\\
\sphinxhline
\sphinxAtStartPar
\sphinxstylestrong{Gestión de calendario y agendas}
&
\sphinxAtStartPar

&
\sphinxAtStartPar

&
\sphinxAtStartPar

\\
\sphinxbottomrule
\end{tabulary}
\sphinxtableafterendhook\par
\sphinxattableend\end{savenotes}


\bigskip\hrule\bigskip



\paragraph{A08: Carta al propietario (+4,0)}
\label{\detokenize{cursos/ofimatica10/periodo01/G02:a08-carta-al-propietario-4-0}}
\sphinxAtStartPar
\DUrole{sd-sphinx-override}{\DUrole{sd-badge}{\DUrole{sd-bg-danger}{\DUrole{sd-bg-text-danger}{Transferencia}}}}
\begin{quote}\begin{description}
\sphinxlineitem{TIEMPO}
\sphinxAtStartPar
59 minutos

\end{description}\end{quote}

\begin{sphinxadmonition}{note}{Descripción de la situación}
\begin{enumerate}
\sphinxsetlistlabels{\arabic}{enumi}{enumii}{}{.}%
\item {} 
\sphinxAtStartPar
Imagine la siguiente situación:

\end{enumerate}
\begin{quote}

\sphinxAtStartPar
“Usted trabaja en un negocio de \sphinxstylestrong{venta de frutas y verduras} con muy \sphinxstylestrong{poco presupuesto}. \sphinxstyleemphasis{Las cuentas de las ganancias no coinciden con los productos vendidos}. Su jefe no tiene el dinero suficiente para adquirir un \sphinxstylestrong{software contable}”. El negocio está ganando mala fama y cada día pierde clientes.
\end{quote}
\begin{enumerate}
\sphinxsetlistlabels{\arabic}{enumi}{enumii}{}{.}%
\setcounter{enumi}{1}
\item {} 
\sphinxAtStartPar
En una \sphinxstylestrong{hoja de cuaderno}, REDACTE una \sphinxstylestrong{carta} aconsejando a su jefe sobre el uso del \sphinxstylestrong{software libre} en el negocio. Tenga en cuenta lo siguiente:
\begin{itemize}
\item {} 
\sphinxAtStartPar
Tenga en cuenta que es un modelo de carta (no solo un escrito).

\item {} 
\sphinxAtStartPar
La carta debe tener:
\begin{itemize}
\item {} 
\sphinxAtStartPar
Buena presentación.

\item {} 
\sphinxAtStartPar
Letra entendible.

\item {} 
\sphinxAtStartPar
Redacción coherente.

\item {} 
\sphinxAtStartPar
Ortografía correcta.

\end{itemize}

\item {} 
\sphinxAtStartPar
Invente un nombre para su jefe.

\item {} 
\sphinxAtStartPar
Cree un nombre para el negocio.

\item {} 
\sphinxAtStartPar
Use una dirección ficticia en Puerto Boyacá. (Justificar la ubicación de la dirección).

\item {} 
\sphinxAtStartPar
Indique las características del software libre.

\item {} 
\sphinxAtStartPar
Indique las ventajas y desventajas de usar el software libre.

\item {} 
\sphinxAtStartPar
Adicione un análisis cuantitativo de ahorro.

\item {} 
\sphinxAtStartPar
Añadir un análisis de seguridad de la información (Malware, Virus, Hacking).

\item {} 
\sphinxAtStartPar
Adicione un análisis frente a la competencia.

\item {} 
\sphinxAtStartPar
Nombrar alguna fruta, verdura, producto o servicio que se ofrezca en el negocio.

\end{itemize}


\begin{savenotes}\sphinxattablestart
\sphinxthistablewithglobalstyle
\centering
\begin{tabulary}{\linewidth}[t]{T}
\sphinxtoprule
\sphinxtableatstartofbodyhook

\begin{sphinxfigure-in-table}
\centering
\capstart
\noindent\sphinxincludegraphics[width=400\sphinxpxdimen]{{image022}.png}
\sphinxfigcaption{\sphinxstylestrong{Situación negocio:} frutas y verduras}\label{\detokenize{cursos/ofimatica10/periodo01/G02:id2}}\end{sphinxfigure-in-table}\relax
\\
\sphinxbottomrule
\end{tabulary}
\sphinxtableafterendhook\par
\sphinxattableend\end{savenotes}

\end{enumerate}
\end{sphinxadmonition}

\sphinxstepscope


\subsubsection{Guía 03: LibreOffice}
\label{\detokenize{cursos/ofimatica10/periodo01/G03:guia-03-libreoffice}}\label{\detokenize{cursos/ofimatica10/periodo01/G03::doc}}

\begin{savenotes}\sphinxattablestart
\sphinxthistablewithglobalstyle
\centering
\begin{tabular}[t]{\X{35}{100}\X{65}{100}}
\sphinxtoprule
\sphinxtableatstartofbodyhook

\begin{savenotes}\sphinxattablestart
\sphinxthistablewithglobalstyle
\centering
\begin{tabulary}{\linewidth}[t]{T}
\sphinxtoprule
\sphinxtableatstartofbodyhook
\sphinxAtStartPar
\sphinxincludegraphics{{escudo}.png}
\\
\sphinxbottomrule
\end{tabulary}
\sphinxtableafterendhook\par
\sphinxattableend\end{savenotes}
&

\begin{savenotes}\sphinxattablestart
\sphinxthistablewithglobalstyle
\centering
\begin{tabulary}{\linewidth}[t]{T}
\sphinxtoprule
\sphinxtableatstartofbodyhook
\sphinxAtStartPar
\sphinxstylestrong{INSTITUCIÓN EDUCATIVA}: John F. Kennedy
\\
\sphinxhline
\sphinxAtStartPar
\sphinxstylestrong{DOCENTE}: Julian Camilo Fonseca Romero
\\
\sphinxhline
\sphinxAtStartPar
\sphinxstylestrong{ASIGNATURA}: Ofimática
\\
\sphinxhline
\sphinxAtStartPar
\sphinxstylestrong{GRADO}: DÉCIMO
\\
\sphinxhline
\sphinxAtStartPar
\sphinxstylestrong{PERIODO:} 01
\\
\sphinxhline
\sphinxAtStartPar
\sphinxstylestrong{TIEMPO}: 10 semanas (4 horas / semana)
\\
\sphinxhline
\sphinxAtStartPar
\sphinxstylestrong{TIEMPO DE TRABAJO}: 10/40 horas
\\
\sphinxbottomrule
\end{tabulary}
\sphinxtableafterendhook\par
\sphinxattableend\end{savenotes}
\\
\sphinxbottomrule
\end{tabular}
\sphinxtableafterendhook\par
\sphinxattableend\end{savenotes}

\begin{sphinxuseclass}{sd-tab-set}
\begin{sphinxuseclass}{sd-tab-item}\subsubsection*{Objetivo}

\begin{sphinxuseclass}{sd-tab-content}\begin{itemize}
\item {} 
\sphinxAtStartPar
\sphinxstylestrong{Temática}: Suite ofimática LibreOffice.

\item {} 
\sphinxAtStartPar
\sphinxstylestrong{Objetivo de aprendizaje}: Identificar las herramientas, funciones y ventajas técnicas de la suite ofimática LibreOffice, reconociendo su naturaleza de código abierto y su capacidad de compatibilidad y portabilidad en diversos entornos operativos.

\end{itemize}

\end{sphinxuseclass}
\end{sphinxuseclass}
\begin{sphinxuseclass}{sd-tab-item}\subsubsection*{Orientaciones curriculares}

\begin{sphinxuseclass}{sd-tab-content}\begin{itemize}
\item {} 
\sphinxAtStartPar
\sphinxstylestrong{Componente}: Uso y apropiación de la Tecnología y la Informática.

\item {} 
\sphinxAtStartPar
\sphinxstylestrong{Competencia}: Genero propuestas innovadoras para el uso y aprovechamiento de productos tecnológicos.

\item {} 
\sphinxAtStartPar
\sphinxstylestrong{Evidencia de aprendizaje}: Reconozco la pertinencia y uso de las licencias de artefactos tecnológicos analógicos y digitales, plataformas y los derechos de autor morales y patrimoniales.

\end{itemize}

\end{sphinxuseclass}
\end{sphinxuseclass}
\begin{sphinxuseclass}{sd-tab-item}\subsubsection*{Criterios de evaluación}

\begin{sphinxuseclass}{sd-tab-content}\begin{itemize}
\item {} 
\sphinxAtStartPar
Actividades desarrolladas en su totalidad.

\item {} 
\sphinxAtStartPar
Participación en clase.

\item {} 
\sphinxAtStartPar
Explicación clara de conceptos.

\item {} 
\sphinxAtStartPar
Argumentación de ideas.

\item {} 
\sphinxAtStartPar
Cumplimiento de las reglas del aula.

\end{itemize}

\end{sphinxuseclass}
\end{sphinxuseclass}
\begin{sphinxuseclass}{sd-tab-item}\subsubsection*{Recursos (1)}

\begin{sphinxuseclass}{sd-tab-content}\begin{itemize}
\item {} 
\sphinxAtStartPar
{\hyperref[\detokenize{cursos/ofimatica10/recursos/T03::doc}]{\sphinxcrossref{\DUrole{std}{\DUrole{std-doc}{TEXTO: LibreOffice}}}}}

\end{itemize}

\end{sphinxuseclass}
\end{sphinxuseclass}
\end{sphinxuseclass}

\paragraph{Mensaje del autor}
\label{\detokenize{cursos/ofimatica10/periodo01/G03:mensaje-del-autor}}\begin{quote}
\begin{quote}

\sphinxAtStartPar
“Esta guía didáctica está diseñada para potenciar su autonomía digital. Se explora la suite LibreOffice, una herramienta de código abierto que permite crear y gestionar documentos de forma gratuita y legal. Se identifican las funciones y ventajas técnicas de LibreOffice, desde el procesador de textos Writer hasta el editor de ecuaciones Math, reconociendo cómo su compatibilidad con formatos como ODF y su naturaleza multiplataforma facilitan el trabajo en diversos entornos operativos como Windows, macOS o Linux.”
\end{quote}

\begin{flushright}
---Camilo Fonseca
\end{flushright}
\end{quote}


\bigskip\hrule\bigskip



\paragraph{A00: Introducción (+0,2)}
\label{\detokenize{cursos/ofimatica10/periodo01/G03:a00-introduccion-0-2}}
\sphinxAtStartPar
\DUrole{sd-sphinx-override}{\DUrole{sd-badge}{\DUrole{sd-bg-light}{\DUrole{sd-bg-text-light}{Exploración}}}}
\begin{enumerate}
\sphinxsetlistlabels{\arabic}{enumi}{enumii}{}{.}%
\item {} 
\sphinxAtStartPar
TRANSCRIBA en su cuaderno el objetivo, las orientaciones curriculares y los criterios de evaluación de la presente guía.

\item {} 
\sphinxAtStartPar
Acorde al mensaje del autor, responda la siguiente pregunta en su cuaderno:
\begin{itemize}
\item {} 
\sphinxAtStartPar
¿Por qué al usar una suite ofimática como \sphinxstylestrong{LibreOffice}, se potencia la autonomía digital?

\end{itemize}

\end{enumerate}


\bigskip\hrule\bigskip



\paragraph{A01: LibreOffice (+0,2)}
\label{\detokenize{cursos/ofimatica10/periodo01/G03:a01-libreoffice-0-2}}
\sphinxAtStartPar
\DUrole{sd-sphinx-override}{\DUrole{sd-badge}{\DUrole{sd-bg-info}{\DUrole{sd-bg-text-info}{Estructuración}}}}
\begin{enumerate}
\sphinxsetlistlabels{\arabic}{enumi}{enumii}{}{.}%
\item {} 
\sphinxAtStartPar
RESPONDA en su cuaderno las siguientes preguntas:
\begin{itemize}
\item {} 
\sphinxAtStartPar
¿Que es LibreOffice?

\item {} 
\sphinxAtStartPar
¿Qué significa \sphinxstylestrong{ODF}?

\end{itemize}

\item {} 
\sphinxAtStartPar
COPIE las siguientes frases en su cuaderno e INDIQUE si es \sphinxstylestrong{falso}, \sphinxstylestrong{verdadero} y por qué:
\begin{itemize}
\item {} 
\sphinxAtStartPar
\sphinxstylestrong{Calc} es el componente de LibreOffice que se utiliza para la creación y edición de hojas de cálculo. \_\_\_

\item {} 
\sphinxAtStartPar
\sphinxstylestrong{LibreOffice} es una suite ofimática exclusiva para usuarios del sistema operativo Windows. \_\_\_

\item {} 
\sphinxAtStartPar
El programa de la suite destinado al diseño y edición de dibujo vectorial recibe el nombre de \sphinxstylestrong{Math}. \_\_\_

\item {} 
\sphinxAtStartPar
Para la gestión y creación de bases de datos, la suite de LibreOffice incorpora la herramienta llamada \sphinxstylestrong{Impress}. \_\_\_

\item {} 
\sphinxAtStartPar
\sphinxstylestrong{Writer} es el procesador de textos que se incluye dentro del conjunto de programas de LibreOffice. \_\_\_

\item {} 
\sphinxAtStartPar
Al ser una suite \sphinxstylestrong{multi\sphinxhyphen{}plataforma}, los documentos generados en LibreOffice no se pueden abrir en sistemas operativos
diferentes al de origen. \_\_\_

\item {} 
\sphinxAtStartPar
El editor que permite trabajar con ecuaciones dentro de esta suite ofimática se llama \sphinxstylestrong{Calc}. \_\_\_

\item {} 
\sphinxAtStartPar
El programa de presentaciones multimedia de LibreOffice se denomina \sphinxstylestrong{Base}. \_\_\_

\item {} 
\sphinxAtStartPar
\sphinxstylestrong{GNU/Linux} y \sphinxstylestrong{MacOS} son dos de los sistemas operativos en los que se puede ejecutar LibreOffice. \_\_\_

\item {} 
\sphinxAtStartPar
El término \sphinxstylestrong{“multi\sphinxhyphen{}plataforma”} significa que el programa solo funciona de manera óptima si se está conectado a una red de servidores local. \_\_\_

\end{itemize}

\end{enumerate}


\bigskip\hrule\bigskip



\paragraph{A02: Paises que usan LibreOffice (+0,2)}
\label{\detokenize{cursos/ofimatica10/periodo01/G03:a02-paises-que-usan-libreoffice-0-2}}
\sphinxAtStartPar
\DUrole{sd-sphinx-override}{\DUrole{sd-badge}{\DUrole{sd-bg-info}{\DUrole{sd-bg-text-info}{Estructuración}}}}
\begin{enumerate}
\sphinxsetlistlabels{\arabic}{enumi}{enumii}{}{.}%
\item {} 
\sphinxAtStartPar
En su cuaderno complete la siguiente tabla:


\begin{savenotes}\sphinxattablestart
\sphinxthistablewithglobalstyle
\centering
\begin{tabulary}{\linewidth}[t]{TTTT}
\sphinxtoprule
\sphinxstyletheadfamily 
\sphinxAtStartPar
PAÍS
&\sphinxstyletheadfamily 
\sphinxAtStartPar
¿En donde se emplea LibreOffice?
&\sphinxstyletheadfamily 
\sphinxAtStartPar
¿Número de Instalaciones de LibreOffice?
&\sphinxstyletheadfamily 
\sphinxAtStartPar
¿Para que se emplea LibreOffice?
\\
\sphinxmidrule
\sphinxtableatstartofbodyhook
\sphinxAtStartPar
Francia
&
\sphinxAtStartPar

&
\sphinxAtStartPar

&
\sphinxAtStartPar

\\
\sphinxhline
\sphinxAtStartPar
España
&
\sphinxAtStartPar

&
\sphinxAtStartPar

&
\sphinxAtStartPar

\\
\sphinxhline
\sphinxAtStartPar
Italia
&
\sphinxAtStartPar

&
\sphinxAtStartPar

&
\sphinxAtStartPar

\\
\sphinxhline
\sphinxAtStartPar
Alemania
&
\sphinxAtStartPar

&
\sphinxAtStartPar

&
\sphinxAtStartPar

\\
\sphinxhline
\sphinxAtStartPar
Taiwan
&
\sphinxAtStartPar

&
\sphinxAtStartPar

&
\sphinxAtStartPar

\\
\sphinxhline
\sphinxAtStartPar
Brasil
&
\sphinxAtStartPar

&
\sphinxAtStartPar

&
\sphinxAtStartPar

\\
\sphinxhline
\sphinxAtStartPar
Colombia
&
\sphinxAtStartPar

&
\sphinxAtStartPar

&
\sphinxAtStartPar

\\
\sphinxbottomrule
\end{tabulary}
\sphinxtableafterendhook\par
\sphinxattableend\end{savenotes}

\end{enumerate}


\bigskip\hrule\bigskip



\paragraph{A03: Herramientas de LibreOffice (+0,2)}
\label{\detokenize{cursos/ofimatica10/periodo01/G03:a03-herramientas-de-libreoffice-0-2}}
\sphinxAtStartPar
\DUrole{sd-sphinx-override}{\DUrole{sd-badge}{\DUrole{sd-bg-info}{\DUrole{sd-bg-text-info}{Estructuración}}}}
\begin{enumerate}
\sphinxsetlistlabels{\arabic}{enumi}{enumii}{}{.}%
\item {} 
\sphinxAtStartPar
En su cuaderno, REALICE un descripción detallada de cada una de las herramientas de LibreOffice.

\item {} 
\sphinxAtStartPar
ANALICE los siguientes casos, eliga la mejor herramienta de LibreOffice para cada caso y explique el por qué:
\begin{itemize}
\item {} \begin{quote}

\sphinxAtStartPar
Un investigador necesita crear fórmulas matemáticas complejas con símbolos avanzados (integrales, matrices) que deben ser interpretadas correctamente por navegadores web. ¿Qué herramienta específica debe usar para generar el código MathML necesario, y por qué es más eficiente usarla como herramienta independiente en lugar de solo dentro de un procesador de textos?
\end{quote}

\item {} \begin{quote}

\sphinxAtStartPar
Una empresa necesita crear diagramas de flujo de procesos industriales donde los nodos deben mantener su posición lógica incluso si se mueven otros elementos, utilizando “Conectores Inteligentes”. Además, requieren exportar estos diagramas en formato vectorial (como SVG o WebP) para una página web. ¿Qué herramienta de LibreOffice usaría y por qué?
\end{quote}

\item {} \begin{quote}

\sphinxAtStartPar
Se requiere crear una presentación que incluya no solo diapositivas, sino también efectos de Fontwork, animaciones personalizadas y gráficos técnicos importados directamente de otros componentes de dibujo vectorial. ¿Qué herramienta utilizaría para orquestar este ecosistema de elementos gráficos y multimedia? ¿Por qué?
\end{quote}

\item {} \begin{quote}

\sphinxAtStartPar
Un administrador de sistemas necesita una solución que permita gestionar relaciones entre tablas, realizar consultas SQL (específicamente soporte para ANSI\sphinxhyphen{}92) y conectarse a motores externos como PostgreSQL o MySQL. ¿Qué herramienta de la suite debe emplear y cómo se diferencia su interfaz de una simple hoja de cálculo para este propósito?
\end{quote}

\item {} \begin{quote}

\sphinxAtStartPar
Un usuario debe redactar un manual que requiere un estilo de formato coherente, una estructura de capítulos compleja, la inserción de gráficos vectoriales y la capacidad de exportar el resultado final a formatos de lectura digital como EPUB y PDF. ¿Qué herramienta de LibreOffice debería elegir y por qué su capacidad de integración con otros componentes es crítica en este caso?
\end{quote}

\item {} \begin{quote}

\sphinxAtStartPar
Un equipo de contabilidad necesita evaluar la viabilidad de un proyecto financiero. Deben procesar grandes volúmenes de datos usando funciones estadísticas y aplicar el “Gestor de escenarios” para analizar diferentes variables de tipo: (“qué pasaría si”). ¿Por qué esta herramienta debe ser superior a una simple hoja de cálculo básica?
\end{quote}

\end{itemize}

\end{enumerate}


\bigskip\hrule\bigskip



\paragraph{A04: Ventajas de LibreOffice (+0,2)}
\label{\detokenize{cursos/ofimatica10/periodo01/G03:a04-ventajas-de-libreoffice-0-2}}
\sphinxAtStartPar
\DUrole{sd-sphinx-override}{\DUrole{sd-badge}{\DUrole{sd-bg-info}{\DUrole{sd-bg-text-info}{Estructuración}}}}
\begin{enumerate}
\sphinxsetlistlabels{\arabic}{enumi}{enumii}{}{.}%
\item {} 
\sphinxAtStartPar
Para cada una de las ventajas de LibreOffice, en su cuaderno responda la siguiente pregunta:
\begin{itemize}
\item {} 
\sphinxAtStartPar
Si se implementara LibreOffice en la Institución Educativa, ¿Realmente podríamos \sphinxstylestrong{beneficiarnos} de dicha ventaja?

\end{itemize}

\end{enumerate}


\bigskip\hrule\bigskip



\paragraph{A05: Examen escrito (+4,0)}
\label{\detokenize{cursos/ofimatica10/periodo01/G03:a05-examen-escrito-4-0}}
\sphinxAtStartPar
\DUrole{sd-sphinx-override}{\DUrole{sd-badge}{\DUrole{sd-bg-danger}{\DUrole{sd-bg-text-danger}{Transferencia}}}}
\begin{quote}\begin{description}
\sphinxlineitem{TIEMPO}
\sphinxAtStartPar
30 minutos

\sphinxlineitem{PREGUNTAS}
\sphinxAtStartPar
10

\end{description}\end{quote}

\begin{sphinxadmonition}{note}{Posibles preguntas}
\begin{enumerate}
\sphinxsetlistlabels{\arabic}{enumi}{enumii}{}{.}%
\item {} 
\sphinxAtStartPar
¿Cómo se define a LibreOffice desde el punto de vista de su licencia y modelo de desarrollo, y qué implica que sea “código abierto” para el usuario final?

\item {} 
\sphinxAtStartPar
Explique qué significa que LibreOffice sea una suite “multiplataforma” y qué ventajas representa esto para la portabilidad de documentos entre diferentes sistemas operativos.

\item {} 
\sphinxAtStartPar
¿Cuál es el formato de archivo nativo de LibreOffice y por qué se considera una ventaja frente al “bloqueo de proveedores”?

\item {} 
\sphinxAtStartPar
Mencione tres entidades gubernamentales o académicas mencionadas en el texto que han migrado a LibreOffice y explique brevemente cuál fue la motivación principal de al menos una de ellas.

\item {} 
\sphinxAtStartPar
Si un usuario necesita redactar un informe que incluya texto y gráficos, ¿qué capacidades le ofrece Writer respecto a la exportación de archivos?

\item {} 
\sphinxAtStartPar
¿Qué función cumple el “Gestor de escenarios” en Calc y por qué es una herramienta clave para la toma de decisiones?

\item {} 
\sphinxAtStartPar
¿Cómo contribuye la integración con Draw y Math al funcionamiento del programa Impress?

\item {} 
\sphinxAtStartPar
Explique la importancia de los “Conectores Inteligentes” en Draw y cómo facilitan la creación de diagramas.

\item {} 
\sphinxAtStartPar
En cuanto a Base, ¿qué tipos de motores de bases de datos puede manejar y qué permite hacer al usuario respecto a la gestión de consultas y relaciones?

\item {} 
\sphinxAtStartPar
¿Para qué situaciones específicas es indispensable utilizar Math en lugar de simplemente escribir símbolos en un procesador de textos convencional?

\item {} 
\sphinxAtStartPar
¿Qué sucede si un usuario intenta abrir un archivo de Draw desde Writer?

\item {} 
\sphinxAtStartPar
¿En qué consiste la ventaja de “Granularidad” en la configuración de las opciones de LibreOffice?

\item {} 
\sphinxAtStartPar
¿Qué implicaciones tiene para una organización el soporte de idiomas en LibreOffice?

\item {} 
\sphinxAtStartPar
A diferencia de las suites de pago, ¿cómo garantiza LibreOffice la disponibilidad de funciones avanzadas como la exportación a PDF?

\item {} 
\sphinxAtStartPar
Si una empresa desea disminuir sus gastos en licencias de software privativo, ¿cómo podría LibreOffice ayudarle a alcanzar ese objetivo, basándose en los ejemplos de España o Alemania?

\item {} 
\sphinxAtStartPar
¿Cómo permite la estructura “abierta y documentada” del formato OpenDocument que los archivos sean leídos incluso por herramientas que no sean parte de la suite LibreOffice?

\item {} 
\sphinxAtStartPar
¿Por qué se dice que Base facilita la administración de bases de datos relacionales siendo similar a otras aplicaciones populares?

\item {} 
\sphinxAtStartPar
Si usted tuviera que convencer a una institución educativa de migrar a LibreOffice, ¿qué tres ventajas principales usaría como argumentos clave?

\item {} 
\sphinxAtStartPar
En relación con la compatibilidad, ¿qué limitaciones o alcances tiene LibreOffice para trabajar con archivos provenientes de Microsoft Office o formatos como CSV y HTML?

\item {} 
\sphinxAtStartPar
¿Es posible abrir y editar en LibreOffice archivos creados originalmente en Microsoft Office sin inconvenientes?

\end{enumerate}
\end{sphinxadmonition}

\sphinxstepscope


\subsubsection{Plan de refuerzo 01}
\label{\detokenize{cursos/ofimatica10/periodo01/R01:plan-de-refuerzo-01}}\label{\detokenize{cursos/ofimatica10/periodo01/R01::doc}}
\begin{sphinxadmonition}{note}{EVALUACIÓN: Plan de refuerzo (10 minutos)}

\sphinxAtStartPar
Estudiar los temas a continuación. Se realiza examen de 5 preguntas. Si el estudiante no asiste en la fecha programada, este será evaluado la próxima vez que asista a clase.
\end{sphinxadmonition}


\bigskip\hrule\bigskip



\paragraph{¿Qué es la ofimática?}
\label{\detokenize{cursos/ofimatica10/periodo01/R01:que-es-la-ofimatica}}
\sphinxAtStartPar
La \sphinxstylestrong{ofimática} es el conjunto de técnicas, aplicaciones y herramientas informáticas que se utilizan para optimizar, automatizar y mejorar los procedimientos o tareas relacionadas con el trabajo de oficina. El término “ofimática” es una combinación de las palabras “oficina” e “informática”.


\bigskip\hrule\bigskip



\paragraph{¿Qué es Google Workspace?}
\label{\detokenize{cursos/ofimatica10/periodo01/R01:que-es-google-workspace}}
\sphinxAtStartPar
\sphinxstylestrong{Google Workspace} es una solución de productividad en línea desarrollada por \sphinxstylestrong{Google} que permite a los usuarios conectarse, crear y colaborar de manera segura. Incluye herramientas esenciales como:
\begin{itemize}
\item {} 
\sphinxAtStartPar
\sphinxstylestrong{Gmail}: Servicio de correo electrónico.

\item {} 
\sphinxAtStartPar
\sphinxstylestrong{Documentos}: Para crear y editar documentos de texto.

\item {} 
\sphinxAtStartPar
\sphinxstylestrong{Hojas de cálculo}: Para trabajar con datos y realizar cálculos.

\item {} 
\sphinxAtStartPar
\sphinxstylestrong{Presentaciones}: Para crear presentaciones visuales.

\item {} 
\sphinxAtStartPar
\sphinxstylestrong{Drive}: Almacenamiento en la nube para guardar y compartir archivos.

\item {} 
\sphinxAtStartPar
\sphinxstylestrong{Calendario}: Para gestionar eventos y citas.

\end{itemize}


\bigskip\hrule\bigskip



\paragraph{¿Qué es Microsoft Office?}
\label{\detokenize{cursos/ofimatica10/periodo01/R01:que-es-microsoft-office}}
\sphinxAtStartPar
\sphinxstylestrong{Microsoft Office} es un conjunto de aplicaciones de software de productividad creado por \sphinxstylestrong{Microsoft}®. Es una herramienta muy popular y ampliamente utilizada en todo el mundo para una variedad de tareas relacionadas con la creación, edición y gestión de documentos, hojas de cálculo, presentaciones y más.

\sphinxAtStartPar
Algunas de las aplicaciones más conocidas que forman parte de Microsoft Office son:
\begin{itemize}
\item {} 
\sphinxAtStartPar
\sphinxstylestrong{Word:} Para crear y editar documentos de texto, como cartas, informes y trabajos académicos.

\item {} 
\sphinxAtStartPar
\sphinxstylestrong{Excel:} Para crear hojas de cálculo, analizar datos, realizar cálculos y crear gráficos.

\item {} 
\sphinxAtStartPar
\sphinxstylestrong{PowerPoint:} Para crear presentaciones visuales con diapositivas, texto, imágenes y animaciones.

\item {} 
\sphinxAtStartPar
\sphinxstylestrong{Outlook:} Para gestionar correos electrónicos, calendarios, contactos y tareas.

\item {} 
\sphinxAtStartPar
\sphinxstylestrong{Access:} Para crear y gestionar bases de datos (menos común en el uso diario).

\end{itemize}


\bigskip\hrule\bigskip



\paragraph{¿Qué es el software libre?}
\label{\detokenize{cursos/ofimatica10/periodo01/R01:que-es-el-software-libre}}
\sphinxAtStartPar
El \sphinxstylestrong{software libre} es un tipo de software que respeta la libertad de los usuarios y la comunidad. Los usuarios tienen la libertad de ejecutar, copiar, distribuir, estudiar, modificar y mejorar el software.

\sphinxAtStartPar
Las cuatro libertades esenciales del software libre son:
\begin{itemize}
\item {} 
\sphinxAtStartPar
\sphinxstylestrong{(Libertad 0):} La libertad de ejecutar el programa como se desee, con cualquier propósito.

\item {} 
\sphinxAtStartPar
\sphinxstylestrong{(Libertad 1):} La libertad de estudiar cómo funciona el programa y cambiarlo para que haga lo que uno quiera. (El acceso al código fuente es una condición necesaria para ello).

\item {} 
\sphinxAtStartPar
\sphinxstylestrong{(Libertad 2):} La libertad de redistribuir copias para ayudar a otros.

\item {} 
\sphinxAtStartPar
\sphinxstylestrong{(Libertad 3):} La libertad de distribuir copias de sus versiones modificadas a terceros.

\end{itemize}


\subparagraph{Ejemplos de software libre}
\label{\detokenize{cursos/ofimatica10/periodo01/R01:ejemplos-de-software-libre}}\begin{itemize}
\item {} 
\sphinxAtStartPar
Sistema Operativo GNU\sphinxhyphen{}LINUX (\$0)

\item {} 
\sphinxAtStartPar
Navegador web Mozilla FIREFOX (\$0)

\item {} 
\sphinxAtStartPar
Editor de imágenes GIMP (\$0)

\item {} 
\sphinxAtStartPar
Editor de gráficos vectoriales INKSCAPE (\$0)

\item {} 
\sphinxAtStartPar
Reproductor multimedia VLC (\$0)

\item {} 
\sphinxAtStartPar
Suite de creación 3D BLENDER (\$0)

\item {} 
\sphinxAtStartPar
Edición de sonido Audacity (\$0)

\item {} 
\sphinxAtStartPar
Suite ofimática Libreoffice (\$0)

\end{itemize}


\bigskip\hrule\bigskip



\paragraph{¿Qué es el software propietario?}
\label{\detokenize{cursos/ofimatica10/periodo01/R01:que-es-el-software-propietario}}
\sphinxAtStartPar
El software propietario, es aquel que no permite a los usuarios acceder, modificar o distribuir su código fuente. Los usuarios tienen restricciones legales y técnicas que les impiden ejercer las libertades que caracterizan al software libre.
\begin{itemize}
\item {} 
\sphinxAtStartPar
El código fuente no está disponible para los usuarios.

\item {} 
\sphinxAtStartPar
Prohibición de copiar, modificar o redistribuir el software.

\item {} 
\sphinxAtStartPar
Generalmente requiere el pago de licencias para su uso.

\item {} 
\sphinxAtStartPar
Las actualizaciones y mejoras dependen exclusivamente del propietario.

\end{itemize}


\subparagraph{Ejemplos de software propietario}
\label{\detokenize{cursos/ofimatica10/periodo01/R01:ejemplos-de-software-propietario}}\begin{itemize}
\item {} 
\sphinxAtStartPar
Sistema operativo Microsoft Windows (\$769,999)

\item {} 
\sphinxAtStartPar
Microsoft Office (suite ofimática) (\$459,999/año)

\item {} 
\sphinxAtStartPar
Adobe Photoshop (editor de imágenes) (\$48,971/mes)

\item {} 
\sphinxAtStartPar
AutoCAD (software de diseño) (\$2,201,999/año)

\end{itemize}


\bigskip\hrule\bigskip



\paragraph{¿Qué es el software pirata?}
\label{\detokenize{cursos/ofimatica10/periodo01/R01:que-es-el-software-pirata}}
\sphinxAtStartPar
El software pirata se refiere a copias ilegales de \sphinxstylestrong{programas informáticos} que se distribuyen o utilizan sin la debida autorización de sus propietarios legales. Esta práctica viola los derechos de autor y la propiedad intelectual


\bigskip\hrule\bigskip



\paragraph{¿Qué es LibreOffice?}
\label{\detokenize{cursos/ofimatica10/periodo01/R01:que-es-libreoffice}}
\sphinxAtStartPar
\sphinxstylestrong{LibreOffice} es una suite de productividad de oficina de código abierto, disponible de forma gratuita, que es compatible con otras suites de oficina y está disponible en una variedad de plataformas. Su formato de archivo nativo es \sphinxstylestrong{Open Document Format (ODF)} y puede abrir y guardar documentos en muchos formatos, incluidos los utilizados por varias versiones de Microsoft Office.

\sphinxAtStartPar
La suite de \sphinxstylestrong{LibreOffice} esta compuesta por las siguientes herramientas:
\begin{itemize}
\item {} 
\sphinxAtStartPar
El procesador de textos \sphinxcode{\sphinxupquote{Writer}}.

\item {} 
\sphinxAtStartPar
La hoja de cálculo \sphinxcode{\sphinxupquote{Calc}}.

\item {} 
\sphinxAtStartPar
El programa de presentaciones \sphinxcode{\sphinxupquote{Impress}}.

\item {} 
\sphinxAtStartPar
El programa de dibujo vectorial \sphinxcode{\sphinxupquote{Draw}}.

\item {} 
\sphinxAtStartPar
La base de datos \sphinxcode{\sphinxupquote{Base}}.

\item {} 
\sphinxAtStartPar
El editor de ecuaciones \sphinxcode{\sphinxupquote{Math}}.

\end{itemize}

\sphinxstepscope


\section{Recursos}
\label{\detokenize{cursos/ofimatica10/recursos/index:recursos}}\label{\detokenize{cursos/ofimatica10/recursos/index::doc}}
\sphinxAtStartPar
En esta sección encuentra los recursos necesarios para el desarrollo del curso.


\bigskip\hrule\bigskip



\subsection{Contenidos}
\label{\detokenize{cursos/ofimatica10/recursos/index:contenidos}}
\sphinxstepscope


\subsubsection{T01: Introducción a la Ofimática}
\label{\detokenize{cursos/ofimatica10/recursos/T01:t01-introduccion-a-la-ofimatica}}\label{\detokenize{cursos/ofimatica10/recursos/T01::doc}}
\sphinxAtStartPar
A continuación se explica el concepto de ofimática.


\bigskip\hrule\bigskip



\paragraph{¿Qué es la ofimática?}
\label{\detokenize{cursos/ofimatica10/recursos/T01:que-es-la-ofimatica}}
\sphinxAtStartPar
%
\begin{footnote}[1]\sphinxAtStartFootnote
Superior, E. F. P. (2023, August 31). ¿Qué es la ofimática y para qué sirve?: herramientas. Esic.edu; ESIC. https://www.esic.edu/business/que\sphinxhyphen{}es\sphinxhyphen{}la\sphinxhyphen{}ofimatica\sphinxhyphen{}para\sphinxhyphen{}que\sphinxhyphen{}sirve\sphinxhyphen{}c
%
\end{footnote} La \sphinxstylestrong{ofimática} es el conjunto de técnicas, aplicaciones y herramientas informáticas que se utilizan para optimizar, automatizar y mejorar los procedimientos o tareas relacionadas con el trabajo de oficina. El término “ofimática” es una combinación de las palabras “oficina” e “informática”.

\sphinxAtStartPar
Las principales áreas que abarca la ofimática incluyen:
\begin{itemize}
\item {} 
\sphinxAtStartPar
Procesamiento de textos (creación y edición de documentos).

\item {} 
\sphinxAtStartPar
Hojas de cálculo (manejo de datos numéricos y fórmulas).

\item {} 
\sphinxAtStartPar
Presentaciones (creación de diapositivas).

\item {} 
\sphinxAtStartPar
Gestión de bases de datos.

\item {} 
\sphinxAtStartPar
Correo electrónico y comunicaciones.

\item {} 
\sphinxAtStartPar
Gestión de calendarios y agendas.

\item {} 
\sphinxAtStartPar
Gestión documental.

\end{itemize}

\sphinxAtStartPar
La ofimática es fundamental en el entorno laboral moderno, ya que permite aumentar la productividad y eficiencia en las tareas administrativas y de gestión de información.


\bigskip\hrule\bigskip



\paragraph{Google Workspace}
\label{\detokenize{cursos/ofimatica10/recursos/T01:google-workspace}}

\begin{savenotes}\sphinxattablestart
\sphinxthistablewithglobalstyle
\centering
\begin{tabulary}{\linewidth}[t]{T}
\sphinxtoprule
\sphinxtableatstartofbodyhook

\begin{sphinxfigure-in-table}
\centering
\capstart
\noindent\sphinxincludegraphics[width=700\sphinxpxdimen]{{image023}.png}
\sphinxfigcaption{Suite de aplicaciones \sphinxstylestrong{Google WorkSpace}}\label{\detokenize{cursos/ofimatica10/recursos/T01:id9}}\end{sphinxfigure-in-table}\relax
\\
\sphinxbottomrule
\end{tabulary}
\sphinxtableafterendhook\par
\sphinxattableend\end{savenotes}

\sphinxAtStartPar
%
\begin{footnote}[2]\sphinxAtStartFootnote
Hill, S. (2024, March 24). Qué es Google One y cómo saber si lo necesitas. WIRED. https://es.wired.com/articulos/que\sphinxhyphen{}es\sphinxhyphen{}google\sphinxhyphen{}one\sphinxhyphen{}y\sphinxhyphen{}como\sphinxhyphen{}saber\sphinxhyphen{}si\sphinxhyphen{}lo\sphinxhyphen{}necesitas
%
\end{footnote} \sphinxstylestrong{Google Workspace} es una solución de productividad en línea desarrollada por \sphinxstylestrong{Google} que permite a los usuarios conectarse, crear y colaborar de manera segura. Incluye herramientas esenciales como:
\begin{itemize}
\item {} 
\sphinxAtStartPar
\sphinxstylestrong{Gmail}: Servicio de correo electrónico.

\item {} 
\sphinxAtStartPar
\sphinxstylestrong{Documentos}: Para crear y editar documentos de texto.

\item {} 
\sphinxAtStartPar
\sphinxstylestrong{Hojas de cálculo}: Para trabajar con datos y realizar cálculos.

\item {} 
\sphinxAtStartPar
\sphinxstylestrong{Presentaciones}: Para crear presentaciones visuales.

\item {} 
\sphinxAtStartPar
\sphinxstylestrong{Drive}: Almacenamiento en la nube para guardar y compartir archivos.

\item {} 
\sphinxAtStartPar
\sphinxstylestrong{Calendario}: Para gestionar eventos y citas.

\end{itemize}


\bigskip\hrule\bigskip



\subparagraph{Google ONE \sphinxhyphen{} AI PLUS (\$19,900/mes) (2026)}
\label{\detokenize{cursos/ofimatica10/recursos/T01:google-one-ai-plus-19-900-mes-2026}}
\sphinxAtStartPar
\sphinxfootnotemark[2] Plan de almacenamiento y la IA de GOOGLE en una sola suscripción.

\sphinxAtStartPar
Teniendo activa esta suscripción, el usuario recibe los siguientes servicios:
\begin{itemize}
\item {} 
\sphinxAtStartPar
Acceso a App de GEMINI con mejores funciones.

\item {} 
\sphinxAtStartPar
200 GB de almacenamiento en la nube.

\item {} 
\sphinxAtStartPar
Copias de seguridad automáticas de dispositivos.

\item {} 
\sphinxAtStartPar
Llamadas de mayor calidad con funciones de video mejoradas con \sphinxstylestrong{Google Meet}.

\item {} 
\sphinxAtStartPar
Herramientas de agenda avanzadas para ahorrar tiempo.

\end{itemize}


\bigskip\hrule\bigskip



\paragraph{Microsoft Office}
\label{\detokenize{cursos/ofimatica10/recursos/T01:microsoft-office}}

\begin{savenotes}\sphinxattablestart
\sphinxthistablewithglobalstyle
\centering
\begin{tabulary}{\linewidth}[t]{T}
\sphinxtoprule
\sphinxtableatstartofbodyhook

\begin{sphinxfigure-in-table}
\centering
\capstart
\noindent\sphinxincludegraphics[width=450\sphinxpxdimen]{{image014}.png}
\sphinxfigcaption{Suite de aplicaciones \sphinxstylestrong{Microsoft Office}}\label{\detokenize{cursos/ofimatica10/recursos/T01:id10}}\end{sphinxfigure-in-table}\relax
\\
\sphinxbottomrule
\end{tabulary}
\sphinxtableafterendhook\par
\sphinxattableend\end{savenotes}

\sphinxAtStartPar
%
\begin{footnote}[3]\sphinxAtStartFootnote
(Visnuh), M. G. (2024, January 24). Con Microsoft Office 365 dominarás la productividad y con el precio de esta licencia en oferta no tendrás que empeñarte para ello. Xataka.com; Xataka. https://www.xataka.com/seleccion/microsoft\sphinxhyphen{}office\sphinxhyphen{}365\sphinxhyphen{}dominaras\sphinxhyphen{}productividad\sphinxhyphen{}precio\sphinxhyphen{}esta\sphinxhyphen{}licencia\sphinxhyphen{}oferta\sphinxhyphen{}no\sphinxhyphen{}tendras\sphinxhyphen{}que\sphinxhyphen{}empenarte\sphinxhyphen{}para\sphinxhyphen{}ello
%
\end{footnote} \sphinxstylestrong{Microsoft Office} es un conjunto de aplicaciones de software de productividad creado por \sphinxstylestrong{Microsoft}®. Es una herramienta muy popular y ampliamente utilizada en todo el mundo para una variedad de tareas relacionadas con la creación, edición y gestión de documentos, hojas de cálculo, presentaciones y más.

\sphinxAtStartPar
Algunas de las aplicaciones más conocidas que forman parte de Microsoft Office son:
\begin{itemize}
\item {} 
\sphinxAtStartPar
\sphinxstylestrong{Word:} Para crear y editar documentos de texto, como cartas, informes y trabajos académicos.

\item {} 
\sphinxAtStartPar
\sphinxstylestrong{Excel:} Para crear hojas de cálculo, analizar datos, realizar cálculos y crear gráficos.

\item {} 
\sphinxAtStartPar
\sphinxstylestrong{PowerPoint:} Para crear presentaciones visuales con diapositivas, texto, imágenes y animaciones.

\item {} 
\sphinxAtStartPar
\sphinxstylestrong{Outlook:} Para gestionar correos electrónicos, calendarios, contactos y tareas.

\item {} 
\sphinxAtStartPar
\sphinxstylestrong{Access:} Para crear y gestionar bases de datos (menos común en el uso diario).

\end{itemize}
\begin{quote}

\sphinxAtStartPar
\sphinxstylestrong{Microsoft Office} está disponible en diferentes versiones y planes de suscripción, incluyendo versiones de escritorio para \sphinxstylestrong{Windows} y \sphinxstylestrong{macOS}, así como versiones en línea accesibles a través de un navegador web. También existen aplicaciones móviles para dispositivos iOS y Android.
\end{quote}


\bigskip\hrule\bigskip



\subparagraph{MICROSOFT 365 Family (\$459,999/año) (2026)}
\label{\detokenize{cursos/ofimatica10/recursos/T01:microsoft-365-family-459-999-ano-2026}}
\sphinxAtStartPar
\sphinxfootnotemark[3] Con una suscripción activa, el usuario recibe los siguientes servicios:
\begin{itemize}
\item {} 
\sphinxAtStartPar
Compartir hasta con 5 personas.

\item {} 
\sphinxAtStartPar
Iniciar sesión hasta en cinco dispositivos diferentes simultáneamente por persona.

\item {} 
\sphinxAtStartPar
Usar el software tanto en PC, Móviles, tabletas y MAC.

\item {} 
\sphinxAtStartPar
Cada persona tiene hasta 6TB de almacenamiento en la nube.

\item {} 
\sphinxAtStartPar
Las aplicaciones tienen características exclusivas y acceso sin conexión.

\item {} 
\sphinxAtStartPar
Seguridad integrada en datos y dispositivos.

\item {} 
\sphinxAtStartPar
Acceso a correo electrónico sin anuncios.

\end{itemize}


\bigskip\hrule\bigskip


\sphinxstepscope


\subsubsection{T02: Software libre}
\label{\detokenize{cursos/ofimatica10/recursos/T02:t02-software-libre}}\label{\detokenize{cursos/ofimatica10/recursos/T02::doc}}
\sphinxAtStartPar
A continuación se explica el concepto de software libre. Sin embargo, primero se debe repasar el concepto de software y sus tipos.


\bigskip\hrule\bigskip



\paragraph{¿Qué es el software?}
\label{\detokenize{cursos/ofimatica10/recursos/T02:que-es-el-software}}
\sphinxAtStartPar
%
\begin{footnote}[1]\sphinxAtStartFootnote
Deitel, P. J., \& Deitel, H. (2014). C++: cómo programar (A. V. Romero Elizondo, Trans.). Pearson Educación.
%
\end{footnote} “El software es un conjunto de programas, instrucciones y reglas informáticas que permiten ejecutar distintas tareas en una computadora. Se trata de las aplicaciones y sistemas que hacen posible que los usuarios puedan trabajar con un \sphinxstylestrong{ordenador} o \sphinxstylestrong{dispositivo electrónico}”.


\subparagraph{Tipos de software}
\label{\detokenize{cursos/ofimatica10/recursos/T02:tipos-de-software}}
\sphinxAtStartPar
\sphinxfootnotemark[1] El software tiene los siguientes tipos:
\begin{itemize}
\item {} 
\sphinxAtStartPar
\sphinxstylestrong{Software de sistema:} Programas básicos que permiten que el hardware funcione (sistemas operativos, controladores).

\item {} 
\sphinxAtStartPar
\sphinxstylestrong{Software de aplicación:} Programas que realizan tareas específicas para los usuarios (procesadores de texto, navegadores web).

\end{itemize}


\subparagraph{Clasificación de software}
\label{\detokenize{cursos/ofimatica10/recursos/T02:clasificacion-de-software}}
\sphinxAtStartPar
%
\begin{footnote}[2]\sphinxAtStartFootnote
Tipos de licencias de software: software libre, propietario y demás. (2022, July 19). Velneo.com. https://velneo.com/blog/licencias\sphinxhyphen{}software\sphinxhyphen{}libre\sphinxhyphen{}propietario\sphinxhyphen{}otros/
%
\end{footnote} El software se puede clasificar de la siguiente manera:
\begin{itemize}
\item {} 
\sphinxAtStartPar
\sphinxstylestrong{Software libre:} Software que puede ser usado con \sphinxstyleemphasis{libertad} y de \sphinxstyleemphasis{forma legal}.

\item {} 
\sphinxAtStartPar
\sphinxstylestrong{Software propietario:} Software que puede ser usado con \sphinxstyleemphasis{restricciones} y de \sphinxstyleemphasis{forma legal}.

\item {} 
\sphinxAtStartPar
\sphinxstylestrong{Software pirata:} Software que es usado de \sphinxstylestrong{forma ilegal} (independiente si es usado con libertad o restricciones).

\end{itemize}


\bigskip\hrule\bigskip



\paragraph{Software libre}
\label{\detokenize{cursos/ofimatica10/recursos/T02:software-libre}}
\sphinxAtStartPar
\sphinxfootnotemark[2] El \sphinxstylestrong{software libre} es un tipo de software que respeta la libertad de los usuarios y la comunidad. Los usuarios tienen la libertad de ejecutar, copiar, distribuir, estudiar, modificar y mejorar el software.

\begin{sphinxadmonition}{note}{COLABORACIÓN}

\sphinxAtStartPar
El software libre promueve la colaboración, la transparencia y el conocimiento compartido, permitiendo que las comunidades de desarrolladores y usuarios trabajen juntos para mejorar el software continuamente.
\end{sphinxadmonition}

\sphinxAtStartPar
El software libre se caracteriza por:
\begin{itemize}
\item {} 
\sphinxAtStartPar
Se centra en la \sphinxstylestrong{libertad} como un principio ético y filosófico.

\item {} 
\sphinxAtStartPar
Enfatiza las cuatro libertades fundamentales mencionadas anteriormente.

\item {} 
\sphinxAtStartPar
Impulsado por la Free Software Foundation (FSF).

\item {} 
\sphinxAtStartPar
Usa licencias que garantizan que las versiones modificadas también sean libres (Copyleft).

\end{itemize}


\subparagraph{4 Libertades del software libre}
\label{\detokenize{cursos/ofimatica10/recursos/T02:libertades-del-software-libre}}
\sphinxAtStartPar
\sphinxfootnotemark[2] Las cuatro libertades esenciales del software libre son:
\begin{itemize}
\item {} 
\sphinxAtStartPar
\sphinxstylestrong{(Libertad 0):} La libertad de ejecutar el programa como se desee, con cualquier propósito.

\item {} 
\sphinxAtStartPar
\sphinxstylestrong{(Libertad 1):} La libertad de estudiar cómo funciona el programa y cambiarlo para que haga lo que uno quiera. (El acceso al código fuente es una condición necesaria para ello).

\item {} 
\sphinxAtStartPar
\sphinxstylestrong{(Libertad 2):} La libertad de redistribuir copias para ayudar a otros.

\item {} 
\sphinxAtStartPar
\sphinxstylestrong{(Libertad 3):} La libertad de distribuir copias de sus versiones modificadas a terceros.

\end{itemize}

\sphinxAtStartPar
Estas \sphinxstylestrong{libertades} son esenciales, no opcionales. Un \sphinxstylestrong{software} es \sphinxstylestrong{libre} si otorga a los usuarios todas estas libertades de manera adecuada.


\bigskip\hrule\bigskip



\paragraph{Ejemplos de software libre}
\label{\detokenize{cursos/ofimatica10/recursos/T02:ejemplos-de-software-libre}}
\sphinxAtStartPar
A continuación se muestran algunos ejemplos del software libre:


\subparagraph{Sistema Operativo GNU\sphinxhyphen{}LINUX (\$0)}
\label{\detokenize{cursos/ofimatica10/recursos/T02:sistema-operativo-gnu-linux-0}}

\begin{savenotes}\sphinxattablestart
\sphinxthistablewithglobalstyle
\centering
\begin{tabulary}{\linewidth}[t]{T}
\sphinxtoprule
\sphinxtableatstartofbodyhook

\begin{sphinxfigure-in-table}
\centering
\capstart
\noindent\sphinxincludegraphics[width=300\sphinxpxdimen]{{image032}.png}
\sphinxfigcaption{Logo GNU/Linux}\label{\detokenize{cursos/ofimatica10/recursos/T02:id35}}\end{sphinxfigure-in-table}\relax
\\
\sphinxbottomrule
\end{tabulary}
\sphinxtableafterendhook\par
\sphinxattableend\end{savenotes}

\sphinxAtStartPar
%
\begin{footnote}[3]\sphinxAtStartFootnote
Rohaut, S. Linux: preparación a la certificación LPIC 1 (exámenes LPI 101 y LPI 102), 3ª Edición. Barcelona: ENI ediciones, 2015.
%
\end{footnote} GNU/Linux es un sistema operativo libre y de código abierto que combina el núcleo Linux con las herramientas del proyecto GNU.

\sphinxAtStartPar
GNU/Linux incluye las siguientes características:
\begin{itemize}
\item {} 
\sphinxAtStartPar
Gratuito y libre para usar, modificar y distribuir.

\item {} 
\sphinxAtStartPar
Alta seguridad y estabilidad.

\item {} 
\sphinxAtStartPar
Gran variedad de distribuciones para diferentes necesidades (Ubuntu, Fedora, Debian, etc.).

\item {} 
\sphinxAtStartPar
Funciona en diversos dispositivos, desde computadores personales hasta servidores.

\end{itemize}
\begin{quote}

\sphinxAtStartPar
El sistema fue creado inicialmente por \sphinxstylestrong{Linus Torvalds} (kernel Linux) en conjunto con el proyecto GNU de \sphinxstylestrong{Richard Stallman}, y hoy es mantenido por una gran comunidad de desarrolladores en todo el mundo.
\end{quote}


\subparagraph{Navegador web Mozilla FIREFOX (\$0)}
\label{\detokenize{cursos/ofimatica10/recursos/T02:navegador-web-mozilla-firefox-0}}

\begin{savenotes}\sphinxattablestart
\sphinxthistablewithglobalstyle
\centering
\begin{tabulary}{\linewidth}[t]{T}
\sphinxtoprule
\sphinxtableatstartofbodyhook

\begin{sphinxfigure-in-table}
\centering
\capstart
\noindent\sphinxincludegraphics[width=360\sphinxpxdimen]{{image042}.png}
\sphinxfigcaption{Logo Navegador Mozilla Firefox}\label{\detokenize{cursos/ofimatica10/recursos/T02:id36}}\end{sphinxfigure-in-table}\relax
\\
\sphinxbottomrule
\end{tabulary}
\sphinxtableafterendhook\par
\sphinxattableend\end{savenotes}

\sphinxAtStartPar
%
\begin{footnote}[4]\sphinxAtStartFootnote
Andreu, A. (2024, September 17). Qué es Mozilla Firefox, diferencias con otros navegadores y extensiones imprescindibles. Computer Hoy. https://computerhoy.20minutos.es/tecnologia/mozilla\sphinxhyphen{}firefox\sphinxhyphen{}diferencias\sphinxhyphen{}otros\sphinxhyphen{}navegadores\sphinxhyphen{}extensiones\sphinxhyphen{}imprescindibles\sphinxhyphen{}1404587
%
\end{footnote} \sphinxstylestrong{Mozilla Firefox} es un navegador web de código abierto y gratuito que se caracteriza por:
\begin{itemize}
\item {} 
\sphinxAtStartPar
Alta velocidad y rendimiento en la navegación.

\item {} 
\sphinxAtStartPar
Fuerte énfasis en la privacidad y seguridad del usuario.

\item {} 
\sphinxAtStartPar
Personalizable a través de extensiones y temas.

\item {} 
\sphinxAtStartPar
Compatible con los estándares web modernos.

\end{itemize}

\sphinxAtStartPar
Firefox es desarrollado por Mozilla Foundation, una organización sin fines de lucro que promueve una internet abierta y accesible para todos. El navegador representa una alternativa libre a opciones privativas como \sphinxstylestrong{Google Chrome}.


\subparagraph{Editor de imágenes GIMP (\$0)}
\label{\detokenize{cursos/ofimatica10/recursos/T02:editor-de-imagenes-gimp-0}}

\begin{savenotes}\sphinxattablestart
\sphinxthistablewithglobalstyle
\centering
\begin{tabulary}{\linewidth}[t]{T}
\sphinxtoprule
\sphinxtableatstartofbodyhook

\begin{sphinxfigure-in-table}
\centering
\capstart
\noindent\sphinxincludegraphics[width=300\sphinxpxdimen]{{image051}.png}
\sphinxfigcaption{Logo \sphinxstylestrong{Gimp}}\label{\detokenize{cursos/ofimatica10/recursos/T02:id37}}\end{sphinxfigure-in-table}\relax
\\
\sphinxbottomrule
\end{tabulary}
\sphinxtableafterendhook\par
\sphinxattableend\end{savenotes}

\sphinxAtStartPar
%
\begin{footnote}[5]\sphinxAtStartFootnote
Delgado, J. M. (2024, September 15). Razones por las que GIMP es la mejor alternativa a Photoshop para editar tus fotos. Computer Hoy. https://computerhoy.20minutos.es/apps/razones\sphinxhyphen{}gimp\sphinxhyphen{}mejor\sphinxhyphen{}alternativa\sphinxhyphen{}photoshop\sphinxhyphen{}editar\sphinxhyphen{}fotos\sphinxhyphen{}1400576
%
\end{footnote} GIMP (\sphinxstylestrong{GNU Image Manipulation Program}) es un editor de imágenes libre y de código abierto que ofrece:
\begin{itemize}
\item {} 
\sphinxAtStartPar
Herramientas profesionales para edición y retoque de imágenes.

\item {} 
\sphinxAtStartPar
Soporte para múltiples formatos de archivo.

\item {} 
\sphinxAtStartPar
Capacidad de trabajar con capas y filtros.

\item {} 
\sphinxAtStartPar
Sistema de plugins para extender funcionalidades.

\end{itemize}

\sphinxAtStartPar
Es una alternativa gratuita a programas comerciales como \sphinxstylestrong{Photoshop}, siendo especialmente popular en la comunidad Linux aunque también está disponible para Windows y MacOS.


\subparagraph{Editor de gráficos vectoriales INKSCAPE (\$0)}
\label{\detokenize{cursos/ofimatica10/recursos/T02:editor-de-graficos-vectoriales-inkscape-0}}

\begin{savenotes}\sphinxattablestart
\sphinxthistablewithglobalstyle
\centering
\begin{tabulary}{\linewidth}[t]{T}
\sphinxtoprule
\sphinxtableatstartofbodyhook

\begin{sphinxfigure-in-table}
\centering
\capstart
\noindent\sphinxincludegraphics[width=360\sphinxpxdimen]{{image061}.png}
\sphinxfigcaption{Logo \sphinxstylestrong{Inkscape}}\label{\detokenize{cursos/ofimatica10/recursos/T02:id38}}\end{sphinxfigure-in-table}\relax
\\
\sphinxbottomrule
\end{tabulary}
\sphinxtableafterendhook\par
\sphinxattableend\end{savenotes}

\sphinxAtStartPar
%
\begin{footnote}[6]\sphinxAtStartFootnote
Inkscape, tu ilustrador de cabecera. (2021, July 19). INTEF. https://intef.es/observatorio\_tecno/inkscape\sphinxhyphen{}tu\sphinxhyphen{}ilustrador\sphinxhyphen{}de\sphinxhyphen{}cabecera/
%
\end{footnote} Inkscape es un editor de gráficos vectoriales libre y de código abierto que ofrece:
\begin{itemize}
\item {} 
\sphinxAtStartPar
Creación y edición de gráficos vectoriales profesionales.

\item {} 
\sphinxAtStartPar
Compatible con estándares como SVG (Scalable Vector Graphics).

\item {} 
\sphinxAtStartPar
Herramientas avanzadas de dibujo y manipulación de formas.

\item {} 
\sphinxAtStartPar
Ideal para crear logotipos, ilustraciones y diseños escalables.

\end{itemize}

\sphinxAtStartPar
Es una alternativa gratuita a programas comerciales como \sphinxstylestrong{Adobe Illustrator}, permitiendo crear gráficos que mantienen su calidad sin importar el tamaño al que se escalen.


\subparagraph{Reproductor multimedia VLC (\$0)}
\label{\detokenize{cursos/ofimatica10/recursos/T02:reproductor-multimedia-vlc-0}}

\begin{savenotes}\sphinxattablestart
\sphinxthistablewithglobalstyle
\centering
\begin{tabulary}{\linewidth}[t]{T}
\sphinxtoprule
\sphinxtableatstartofbodyhook

\begin{sphinxfigure-in-table}
\centering
\capstart
\noindent\sphinxincludegraphics[width=400\sphinxpxdimen]{{image071}.png}
\sphinxfigcaption{Logo \sphinxstylestrong{VLC}}\label{\detokenize{cursos/ofimatica10/recursos/T02:id39}}\end{sphinxfigure-in-table}\relax
\\
\sphinxbottomrule
\end{tabulary}
\sphinxtableafterendhook\par
\sphinxattableend\end{savenotes}

\sphinxAtStartPar
%
\begin{footnote}[7]\sphinxAtStartFootnote
Delgado, J. M. (2025, June 14). VLC sigue muy vivo: cosas sorprendentes que puedes hacer con el mejor reproductor de vídeo para PC. Computer Hoy. https://computerhoy.20minutos.es/tecnologia/vlc\sphinxhyphen{}sigue\sphinxhyphen{}muy\sphinxhyphen{}vivo\sphinxhyphen{}cosas\sphinxhyphen{}sorprendentes\sphinxhyphen{}puedes\sphinxhyphen{}hacer\sphinxhyphen{}mejor\sphinxhyphen{}reproductor\sphinxhyphen{}video\sphinxhyphen{}pc\sphinxhyphen{}1463759
%
\end{footnote} VLC Media Player es un reproductor multimedia libre y de código abierto que se caracteriza por:
\begin{itemize}
\item {} 
\sphinxAtStartPar
Capacidad para reproducir casi cualquier formato de audio y video.

\item {} 
\sphinxAtStartPar
No requiere códecs adicionales para funcionar.

\item {} 
\sphinxAtStartPar
Permite streaming y conversión de archivos multimedia.

\item {} 
\sphinxAtStartPar
Funciona en múltiples sistemas operativos (Windows, Linux, Mac).

\end{itemize}

\sphinxAtStartPar
Desarrollado por VideoLAN Organization, VLC es uno de los reproductores más populares debido a su versatilidad y facilidad de uso. Es completamente gratuito y se mantiene actualizado por una comunidad activa de desarrolladores.


\subparagraph{Suite de creación 3D BLENDER (\$0)}
\label{\detokenize{cursos/ofimatica10/recursos/T02:suite-de-creacion-3d-blender-0}}

\begin{savenotes}\sphinxattablestart
\sphinxthistablewithglobalstyle
\centering
\begin{tabulary}{\linewidth}[t]{T}
\sphinxtoprule
\sphinxtableatstartofbodyhook

\begin{sphinxfigure-in-table}
\centering
\capstart
\noindent\sphinxincludegraphics[width=360\sphinxpxdimen]{{image081}.png}
\sphinxfigcaption{Logo \sphinxstylestrong{BLENDER}}\label{\detokenize{cursos/ofimatica10/recursos/T02:id40}}\end{sphinxfigure-in-table}\relax
\\
\sphinxbottomrule
\end{tabulary}
\sphinxtableafterendhook\par
\sphinxattableend\end{savenotes}

\sphinxAtStartPar
%
\begin{footnote}[8]\sphinxAtStartFootnote
Nava, J. (2026, February 19). ¿Qué es blender? sus usos y beneficios. IMSED Business \& Tech School. https://imsed.com/postgrados/tecnologia/que\sphinxhyphen{}es\sphinxhyphen{}blender\sphinxhyphen{}usos\sphinxhyphen{}beneficios/
%
\end{footnote} Blender es una suite de \sphinxstylestrong{creación 3D} de código abierto y gratuita que abarca todo el ciclo de producción digital. Permite realizar modelado, renderizado, animación, edición de video, composición y creación de efectos visuales (VFX).

\sphinxAtStartPar
Características principales:
\begin{itemize}
\item {} 
\sphinxAtStartPar
\sphinxstylestrong{Modelado y Esculpido:} Herramientas avanzadas para crear formas complejas y personajes.

\item {} 
\sphinxAtStartPar
\sphinxstylestrong{Animación y Rigging:} Sistema profesional para dar movimiento a objetos y modelos orgánicos.

\item {} 
\sphinxAtStartPar
\sphinxstylestrong{Motores de Renderizado:} Incluye Cycles (trazado de rayos de alta fidelidad) y Eevee (renderizado en tiempo real).

\item {} 
\sphinxAtStartPar
\sphinxstylestrong{Versatilidad:} Al ser multiplataforma y ligero, es utilizado tanto por artistas independientes como por estudios de videojuegos y cine.

\end{itemize}


\subparagraph{Edición de sonido Audacity (\$0)}
\label{\detokenize{cursos/ofimatica10/recursos/T02:edicion-de-sonido-audacity-0}}

\begin{savenotes}\sphinxattablestart
\sphinxthistablewithglobalstyle
\centering
\begin{tabulary}{\linewidth}[t]{T}
\sphinxtoprule
\sphinxtableatstartofbodyhook

\begin{sphinxfigure-in-table}
\centering
\capstart
\noindent\sphinxincludegraphics[width=360\sphinxpxdimen]{{image091}.png}
\sphinxfigcaption{Logo \sphinxstylestrong{Audacity}}\label{\detokenize{cursos/ofimatica10/recursos/T02:id41}}\end{sphinxfigure-in-table}\relax
\\
\sphinxbottomrule
\end{tabulary}
\sphinxtableafterendhook\par
\sphinxattableend\end{savenotes}

\sphinxAtStartPar
%
\begin{footnote}[9]\sphinxAtStartFootnote
Audacity: una forma diferente de estar en la onda. (2023, November 20). INTEF. https://intef.es/observatorio\_tecno/audacity\sphinxhyphen{}una\sphinxhyphen{}forma\sphinxhyphen{}diferente\sphinxhyphen{}de\sphinxhyphen{}estar\sphinxhyphen{}en\sphinxhyphen{}la\sphinxhyphen{}onda/
%
\end{footnote} Audacity es un \sphinxstylestrong{editor y grabador de audio digital} de código abierto y gratuito, ampliamente utilizado para la producción de podcasts, edición de música y procesamiento de sonido. Audacity tiene las siguientes características:
\begin{itemize}
\item {} 
\sphinxAtStartPar
\sphinxstylestrong{Grabación multipista:} Permite capturar audio en vivo a través de micrófonos o mezcladores.

\item {} 
\sphinxAtStartPar
\sphinxstylestrong{Edición sencilla:} Ofrece herramientas para cortar, copiar, pegar y mezclar sonidos con una interfaz intuitiva.

\item {} 
\sphinxAtStartPar
\sphinxstylestrong{Procesamiento de efectos:} Incluye funciones para reducción de ruido, ecualización, normalización y cambio de tono o velocidad.

\item {} 
\sphinxAtStartPar
\sphinxstylestrong{Compatibilidad:} Soporta diversos formatos de archivo como MP3, WAV, AIFF y Ogg Vorbis.

\end{itemize}


\subparagraph{Suite ofimática Libreoffice (\$0)}
\label{\detokenize{cursos/ofimatica10/recursos/T02:suite-ofimatica-libreoffice-0}}

\begin{savenotes}\sphinxattablestart
\sphinxthistablewithglobalstyle
\centering
\begin{tabulary}{\linewidth}[t]{T}
\sphinxtoprule
\sphinxtableatstartofbodyhook

\begin{sphinxfigure-in-table}
\centering
\capstart
\noindent\sphinxincludegraphics[width=360\sphinxpxdimen]{{image101}.png}
\sphinxfigcaption{\sphinxstylestrong{LibreOffice}}\label{\detokenize{cursos/ofimatica10/recursos/T02:id42}}\end{sphinxfigure-in-table}\relax
\\
\sphinxbottomrule
\end{tabulary}
\sphinxtableafterendhook\par
\sphinxattableend\end{savenotes}

\sphinxAtStartPar
%
\begin{footnote}[10]\sphinxAtStartFootnote
Onieva, D. (2020, February 17). LibreOffice, la suite gratuita que compite cara a cara con Office. SoftZone. https://www.softzone.es/programas/oficina/libreoffice/
%
\end{footnote} LibreOffice es una suite ofimática libre y de código abierto que incluye:
\begin{itemize}
\item {} 
\sphinxAtStartPar
\sphinxstylestrong{Writer:} Procesador de textos similar a Microsoft Word.

\item {} 
\sphinxAtStartPar
\sphinxstylestrong{Calc:} Hoja de cálculo comparable a Microsoft Excel.

\item {} 
\sphinxAtStartPar
\sphinxstylestrong{Impress:} Programa de presentaciones equivalente a PowerPoint.

\item {} 
\sphinxAtStartPar
\sphinxstylestrong{Base:} Gestor de bases de datos.

\item {} 
\sphinxAtStartPar
\sphinxstylestrong{Draw:} Editor de dibujos y diagramas.

\item {} 
\sphinxAtStartPar
\sphinxstylestrong{Math:} Editor de fórmulas matemáticas.

\end{itemize}

\sphinxAtStartPar
LibreOffice es completamente gratuito, compatible con los formatos de \sphinxstylestrong{Microsoft Office} y está disponible en múltiples idiomas. Es desarrollado por \sphinxstylestrong{The Document Foundation} y cuenta con una gran comunidad de usuarios y desarrolladores que lo mantienen actualizado constantemente.


\bigskip\hrule\bigskip



\paragraph{Software propietario Ø}
\label{\detokenize{cursos/ofimatica10/recursos/T02:software-propietario-o}}
\sphinxAtStartPar
\sphinxfootnotemark[2] El software propietario, es aquel que no permite a los usuarios acceder, modificar o distribuir su código fuente. Los usuarios tienen restricciones legales y técnicas que les impiden ejercer las libertades que caracterizan al software libre.

\begin{sphinxadmonition}{note}{Control total}

\sphinxAtStartPar
El \sphinxstylestrong{software privativo} mantiene el control total sobre el programa en manos de su propietario, lo que puede limitar la libertad de elección y personalización por parte de los usuarios.
\end{sphinxadmonition}

\sphinxAtStartPar
El software privativo se caracteriza por:
\begin{itemize}
\item {} 
\sphinxAtStartPar
El código fuente no está disponible para los usuarios.

\item {} 
\sphinxAtStartPar
Prohibición de copiar, modificar o redistribuir el software.

\item {} 
\sphinxAtStartPar
Generalmente requiere el pago de licencias para su uso.

\item {} 
\sphinxAtStartPar
Las actualizaciones y mejoras dependen exclusivamente del propietario.

\end{itemize}


\bigskip\hrule\bigskip



\paragraph{Ejemplos de software propietario}
\label{\detokenize{cursos/ofimatica10/recursos/T02:ejemplos-de-software-propietario}}
\sphinxAtStartPar
A continuación se muestran algunos ejemplos del software propietario:


\subparagraph{Sistema operativo Microsoft Windows (\$769,999)}
\label{\detokenize{cursos/ofimatica10/recursos/T02:sistema-operativo-microsoft-windows-769-999}}

\begin{savenotes}\sphinxattablestart
\sphinxthistablewithglobalstyle
\centering
\begin{tabulary}{\linewidth}[t]{T}
\sphinxtoprule
\sphinxtableatstartofbodyhook

\begin{sphinxfigure-in-table}
\centering
\capstart
\noindent\sphinxincludegraphics[width=300\sphinxpxdimen]{{image11}.png}
\sphinxfigcaption{Logo Microsoft \sphinxstylestrong{Windows}}\label{\detokenize{cursos/ofimatica10/recursos/T02:id43}}\end{sphinxfigure-in-table}\relax
\\
\sphinxbottomrule
\end{tabulary}
\sphinxtableafterendhook\par
\sphinxattableend\end{savenotes}

\sphinxAtStartPar
%
\begin{footnote}[11]\sphinxAtStartFootnote
Corrales, R. (2024, August 19). Microsoft, el gigante de la informática: historia, Windows, todos sus servicios y la IA de Copilot. Computer Hoy. https://computerhoy.20minutos.es/industria/microsoft\sphinxhyphen{}gigante\sphinxhyphen{}informatica\sphinxhyphen{}historia\sphinxhyphen{}windows\sphinxhyphen{}todos\sphinxhyphen{}servicios\sphinxhyphen{}ia\sphinxhyphen{}copilot\sphinxhyphen{}1399819
%
\end{footnote} Microsoft Windows es el sistema operativo más utilizado en computadoras personales. Sus características principales incluyen:
\begin{itemize}
\item {} 
\sphinxAtStartPar
Interfaz gráfica intuitiva y fácil de usar.

\item {} 
\sphinxAtStartPar
Amplia compatibilidad con hardware y software.

\item {} 
\sphinxAtStartPar
Actualizaciones regulares de seguridad y funcionalidad.

\item {} 
\sphinxAtStartPar
Gran variedad de aplicaciones disponibles.

\end{itemize}

\sphinxAtStartPar
Desarrollado por \sphinxstylestrong{Microsoft Corporation}, Windows es un software privativo que requiere la compra de una licencia para su uso legal. Desde su primera versión en 1985, ha evolucionado constantemente.


\subparagraph{Microsoft Office (suite ofimática) (\$459,999/año)}
\label{\detokenize{cursos/ofimatica10/recursos/T02:microsoft-office-suite-ofimatica-459-999-ano}}

\begin{savenotes}\sphinxattablestart
\sphinxthistablewithglobalstyle
\centering
\begin{tabulary}{\linewidth}[t]{T}
\sphinxtoprule
\sphinxtableatstartofbodyhook

\begin{sphinxfigure-in-table}
\centering
\capstart
\noindent\sphinxincludegraphics[width=300\sphinxpxdimen]{{image12}.png}
\sphinxfigcaption{Microsoft \sphinxstylestrong{Office}}\label{\detokenize{cursos/ofimatica10/recursos/T02:id44}}\end{sphinxfigure-in-table}\relax
\\
\sphinxbottomrule
\end{tabulary}
\sphinxtableafterendhook\par
\sphinxattableend\end{savenotes}

\sphinxAtStartPar
%
\begin{footnote}[12]\sphinxAtStartFootnote
Delgado, J. M. (2024, December 15). Funciones ocultas de Microsoft Office que deberías usar ahora mismo para mejorar tu productividad. Computer Hoy. https://computerhoy.20minutos.es/apps/funciones\sphinxhyphen{}ocultas\sphinxhyphen{}microsoft\sphinxhyphen{}office\sphinxhyphen{}deberias\sphinxhyphen{}usar\sphinxhyphen{}ahora\sphinxhyphen{}mismo\sphinxhyphen{}mejorar\sphinxhyphen{}productividad\sphinxhyphen{}1428291
%
\end{footnote} Microsoft Office es una suite de programas de oficina desarrollada por Microsoft que incluye:
\begin{itemize}
\item {} 
\sphinxAtStartPar
\sphinxstylestrong{Word:} Procesador de textos para crear documentos.

\item {} 
\sphinxAtStartPar
\sphinxstylestrong{Excel:} Hoja de cálculo para análisis de datos y cálculos.

\item {} 
\sphinxAtStartPar
\sphinxstylestrong{PowerPoint:} Programa para crear presentaciones.

\item {} 
\sphinxAtStartPar
\sphinxstylestrong{Outlook:} Cliente de correo electrónico y calendario.

\item {} 
\sphinxAtStartPar
\sphinxstylestrong{Access:} Sistema de gestión de bases de datos.

\end{itemize}

\sphinxAtStartPar
Es una de las suites ofimáticas más utilizadas en el mundo empresarial y educativo. Requiere la compra de una licencia o suscripción para su uso legal, y está disponible tanto en versión de escritorio como en la nube (Microsoft 365).


\subparagraph{Adobe Photoshop (editor de imágenes) (\$48,971/mes)}
\label{\detokenize{cursos/ofimatica10/recursos/T02:adobe-photoshop-editor-de-imagenes-48-971-mes}}

\begin{savenotes}\sphinxattablestart
\sphinxthistablewithglobalstyle
\centering
\begin{tabulary}{\linewidth}[t]{T}
\sphinxtoprule
\sphinxtableatstartofbodyhook

\begin{sphinxfigure-in-table}
\centering
\capstart
\noindent\sphinxincludegraphics[width=360\sphinxpxdimen]{{image13}.png}
\sphinxfigcaption{Logo Adobe \sphinxstylestrong{Photoshop}}\label{\detokenize{cursos/ofimatica10/recursos/T02:id45}}\end{sphinxfigure-in-table}\relax
\\
\sphinxbottomrule
\end{tabulary}
\sphinxtableafterendhook\par
\sphinxattableend\end{savenotes}

\sphinxAtStartPar
%
\begin{footnote}[13]\sphinxAtStartFootnote
de Luis, E. R. (2026, March 15). En 1990 un estudiante tuvo un problema mostrando imágenes en su Mac. La solución fue crear uno de los softwares más exitosos de la historia. Xataka.com; Xataka. https://www.xataka.com/aplicaciones/1987\sphinxhyphen{}tuvo\sphinxhyphen{}problema\sphinxhyphen{}mostrando\sphinxhyphen{}imagenes\sphinxhyphen{}su\sphinxhyphen{}mac\sphinxhyphen{}asi\sphinxhyphen{}que\sphinxhyphen{}creo\sphinxhyphen{}app\sphinxhyphen{}hoy\sphinxhyphen{}editor\sphinxhyphen{}imagenes\sphinxhyphen{}usado\sphinxhyphen{}historia
%
\end{footnote} Adobe Photoshop es el software de edición de imágenes más popular y potente del mercado. Sus características principales incluyen:
\begin{itemize}
\item {} 
\sphinxAtStartPar
Herramientas avanzadas para edición y manipulación de imágenes.

\item {} 
\sphinxAtStartPar
Trabajo con capas, máscaras y efectos profesionales.

\item {} 
\sphinxAtStartPar
Soporte para múltiples formatos de archivo.

\item {} 
\sphinxAtStartPar
Integración con otros productos de Adobe Creative Suite.

\item {} 
\sphinxAtStartPar
Herramientas de retoque fotográfico profesional.

\end{itemize}

\sphinxAtStartPar
Desarrollado por Adobe Inc., Photoshop es el estándar de la industria para el diseño gráfico, fotografía y arte digital. Requiere una suscripción mensual a través de Adobe Creative Cloud para su uso legal.


\subparagraph{AutoCAD (software de diseño) (\$2,201,999/año)}
\label{\detokenize{cursos/ofimatica10/recursos/T02:autocad-software-de-diseno-2-201-999-ano}}

\begin{savenotes}\sphinxattablestart
\sphinxthistablewithglobalstyle
\centering
\begin{tabulary}{\linewidth}[t]{T}
\sphinxtoprule
\sphinxtableatstartofbodyhook

\begin{sphinxfigure-in-table}
\centering
\capstart
\noindent\sphinxincludegraphics[width=360\sphinxpxdimen]{{image14}.png}
\sphinxfigcaption{Logo Autodesk \sphinxstylestrong{AutoCAD}}\label{\detokenize{cursos/ofimatica10/recursos/T02:id46}}\end{sphinxfigure-in-table}\relax
\\
\sphinxbottomrule
\end{tabulary}
\sphinxtableafterendhook\par
\sphinxattableend\end{savenotes}

\sphinxAtStartPar
%
\begin{footnote}[14]\sphinxAtStartFootnote
Lautrec, T. (n.d.). ¿Qué es AutoCAD y para qué sirve? Toulouse Lautrec. Retrieved May 7, 2026, from https://www.toulouselautrec.edu.pe/blogs/que\sphinxhyphen{}es\sphinxhyphen{}autocad
%
\end{footnote} AutoCAD es un software de diseño asistido por computadora (CAD) desarrollado por Autodesk. Sus características principales incluyen:
\begin{itemize}
\item {} 
\sphinxAtStartPar
Diseño y dibujo técnico en 2D y 3D.

\item {} 
\sphinxAtStartPar
Herramientas precisas para arquitectura e ingeniería.

\item {} 
\sphinxAtStartPar
Creación de planos y modelos profesionales.

\item {} 
\sphinxAtStartPar
Amplia biblioteca de componentes y símbolos.

\item {} 
\sphinxAtStartPar
Capacidad de colaboración y compartir proyectos.

\end{itemize}

\sphinxAtStartPar
Es ampliamente utilizado en campos como la arquitectura, ingeniería, diseño industrial y construcción. AutoCAD es un software privativo que requiere una suscripción para su uso legal, siendo uno de los estándares de la industria para el diseño técnico profesional.


\bigskip\hrule\bigskip



\paragraph{Software PIRATA!!!}
\label{\detokenize{cursos/ofimatica10/recursos/T02:software-pirata}}\begin{quote}

\sphinxAtStartPar
\sphinxstyleemphasis{“El software pirata se refiere a copias ilegales de \sphinxstylestrong{programas informáticos} que se distribuyen o utilizan sin la debida autorización de sus propietarios legales. Esta práctica viola los derechos de autor y la propiedad intelectual”.}
\end{quote}

\sphinxAtStartPar
%
\begin{footnote}[15]\sphinxAtStartFootnote
Textos: Jorge Eliécer Lozano / Ilustraciones: Gissell Andreína García Rodríguez. (n.d.). El software pirata representa un riesgo inminente. Edu.co. Retrieved May 7, 2026, from https://notired.univalle.edu.co/el\sphinxhyphen{}software\sphinxhyphen{}pirata\sphinxhyphen{}representa\sphinxhyphen{}un\sphinxhyphen{}riesgo\sphinxhyphen{}inminente/
%
\end{footnote} El \sphinxstylestrong{software pirata} se caracteriza por:
\begin{itemize}
\item {} 
\sphinxAtStartPar
Se obtiene y distribuye sin pagar las licencias correspondientes.

\item {} 
\sphinxAtStartPar
Viola los derechos de autor y la propiedad intelectual.

\item {} 
\sphinxAtStartPar
Puede contener Malware por ejemplo:
\begin{itemize}
\item {} 
\sphinxAtStartPar
\sphinxstylestrong{Virus}: impacta la integridad del sistema. Su objetivo es dañar.

\item {} 
\sphinxAtStartPar
\sphinxstylestrong{Troyanos:} impacta la confidencialidad del sistema. Su objetivo es hacerse pasar por un software confiable pero en realidad esta ejecutando acciones maliciosas.

\end{itemize}

\item {} 
\sphinxAtStartPar
No recibe actualizaciones oficiales ni soporte técnico.

\item {} 
\sphinxAtStartPar
Su uso puede resultar en problemas legales y sanciones.

\end{itemize}

\begin{sphinxadmonition}{note}{ADVERTENCIA}

\sphinxAtStartPar
El uso de \sphinxstylestrong{software pirata} es ilegal y puede tener consecuencias graves, incluyendo:
\begin{itemize}
\item {} 
\sphinxAtStartPar
Multas y sanciones legales (CARCEL Y MULTAS MUY ALTAS).

\item {} 
\sphinxAtStartPar
Riesgos de seguridad informática (VIRUS, TROYANOS).

\item {} 
\sphinxAtStartPar
Pérdida de funcionalidades y actualizaciones (SOFTWARE A MEDIAS, POCO RENDIMIENTO).

\item {} 
\sphinxAtStartPar
Daños a la industria del software y su desarrollo (PERDIDAS ECONÓMICAS, DESPIDOS).

\end{itemize}
\end{sphinxadmonition}


\bigskip\hrule\bigskip


\sphinxstepscope


\subsubsection{T03: LibreOffice}
\label{\detokenize{cursos/ofimatica10/recursos/T03:t03-libreoffice}}\label{\detokenize{cursos/ofimatica10/recursos/T03::doc}}
\sphinxAtStartPar
A continuación se describe en detalle el software LibreOffice.


\bigskip\hrule\bigskip



\paragraph{Suite Ofimática LibreOffice}
\label{\detokenize{cursos/ofimatica10/recursos/T03:suite-ofimatica-libreoffice}}

\begin{savenotes}\sphinxattablestart
\sphinxthistablewithglobalstyle
\centering
\begin{tabulary}{\linewidth}[t]{T}
\sphinxtoprule
\sphinxtableatstartofbodyhook

\begin{sphinxfigure-in-table}
\centering
\capstart
\noindent\sphinxincludegraphics[width=300\sphinxpxdimen]{{image15}.png}
\sphinxfigcaption{Suite LibreOffice}\label{\detokenize{cursos/ofimatica10/recursos/T03:id7}}\end{sphinxfigure-in-table}\relax
\\
\sphinxbottomrule
\end{tabulary}
\sphinxtableafterendhook\par
\sphinxattableend\end{savenotes}

\sphinxAtStartPar
%
\begin{footnote}[1]\sphinxAtStartFootnote
Antonio Fernández, B. (2025, May 9). Guía de iniciación 24.8 \sphinxhyphen{} Prefacio. Libreoffice.org. https://books.libreoffice.org/es/GS248/GS24800\sphinxhyphen{}Prefacio.html
%
\end{footnote} \sphinxstylestrong{LibreOffice} es una suite de productividad de oficina de código abierto, disponible de forma gratuita, que es compatible con otras suites de oficina y está disponible en una variedad de plataformas. Su formato de archivo nativo es \sphinxstylestrong{Open Document Format (ODF)} y puede abrir y guardar documentos en muchos formatos, incluidos los utilizados por varias versiones de Microsoft Office.

\sphinxAtStartPar
La suite de \sphinxstylestrong{LibreOffice} esta compuesta por las siguientes herramientas:
\begin{itemize}
\item {} 
\sphinxAtStartPar
El procesador de textos \sphinxcode{\sphinxupquote{Writer}}.

\item {} 
\sphinxAtStartPar
La hoja de cálculo \sphinxcode{\sphinxupquote{Calc}}.

\item {} 
\sphinxAtStartPar
El programa de presentaciones \sphinxcode{\sphinxupquote{Impress}}.

\item {} 
\sphinxAtStartPar
El programa de dibujo vectorial \sphinxcode{\sphinxupquote{Draw}}.

\item {} 
\sphinxAtStartPar
La base de datos \sphinxcode{\sphinxupquote{Base}}.

\item {} 
\sphinxAtStartPar
El editor de ecuaciones \sphinxcode{\sphinxupquote{Math}}.

\end{itemize}

\sphinxAtStartPar
LibreOffice es una suite ofimática multi\sphinxhyphen{}plataforma, lo que quiere decir que se puede ejecutar en diferentes plataformas y sistemas operativos. Por consiguiente usted puede disfrutar del mismo programa en \sphinxstylestrong{Windows}, \sphinxstylestrong{GNU/Linux} o \sphinxstylestrong{MacOS}, facilitando la portabilidad entre diferentes sistemas de los documentos generados.


\bigskip\hrule\bigskip



\paragraph{Entidades que usan LibreOffice}
\label{\detokenize{cursos/ofimatica10/recursos/T03:entidades-que-usan-libreoffice}}
\sphinxAtStartPar
%
\begin{footnote}[2]\sphinxAtStartFootnote
(N.d.). Libreoffice.org. Retrieved May 17, 2026, from https://es.libreoffice.org/who\sphinxhyphen{}uses\sphinxhyphen{}libreoffice/
%
\end{footnote} Millones de personas utilizan LibreOffice todos los días en sus hogares, negocios, organizaciones benéficas y sectores gubernamentales. De entre estos grupos de usuarios pueden destacarse:


\begin{savenotes}\sphinxattablestart
\sphinxthistablewithglobalstyle
\centering
\begin{tabulary}{\linewidth}[t]{TT}
\sphinxtoprule
\sphinxstyletheadfamily 
\sphinxAtStartPar
País
&\sphinxstyletheadfamily 
\sphinxAtStartPar
Descripción
\\
\sphinxmidrule
\sphinxtableatstartofbodyhook
\sphinxAtStartPar
\sphinxstylestrong{FRANCIA}
&
\sphinxAtStartPar
El equipo de trabajo interministerial para el software libre del Gobierno francés, ejecuta LibreOffice en aproximadamente 500,000 equipos. El software se emplea en varios ministerios, incluidos los de Energía, Defensa, Agricultura y Educación.
\\
\sphinxhline
\sphinxAtStartPar
\sphinxstylestrong{ESPAÑA}
&
\sphinxAtStartPar
La administración de la comunidad autónoma de Valencia, en España, ha instalado LibreOffice en 120,000 equipos. Con esto, la administración ha ganado independencia frente a proveedores de TI y ha disminuido los gastos en licencias de software privativo.
\\
\sphinxhline
\sphinxAtStartPar
\sphinxstylestrong{ITALIA}
&
\sphinxAtStartPar
El Ministerio de Defensa de Italia se encuentra en un proceso de transición hacia LibreOffice y el formato OpenDocument (ODF) en más de 100,000 equipos. El Ministerio, asimismo, ha desarrollado cursos en línea para ayudar con el cambio a LibreOffice.
\\
\sphinxhline
\sphinxAtStartPar
\sphinxstylestrong{ALEMANIA}
&
\sphinxAtStartPar
En Alemania, la ciudad de Múnich ha migrado más de 15,000 equipos a LibreOffice. Esto forma parte de una transición más ambiciosa hacia el software libre y de código abierto (FOSS), que incluye el sistema operativo GNU/Linux.
\\
\sphinxhline
\sphinxAtStartPar
\sphinxstylestrong{TAIWAN}
&
\sphinxAtStartPar
El Ministerio de Finanzas de Taiwán ha instalado LibreOffice en más de 24,000 PC. Con esta iniciativa, el formato estandarizado OpenDocument está convirtiéndose en la norma para el intercambio de datos entre los diversos departamentos.
\\
\sphinxhline
\sphinxAtStartPar
\sphinxstylestrong{BRASIL}
&
\sphinxAtStartPar
En Brasil, la UNESP ―Universidade Estadual Paulista― ha realizado más de 10,000 instalaciones de LibreOffice como parte de una migración general hacia el software libre y de código abierto, la cual incluye el sistema operativo GNU/Linux.
\\
\sphinxhline
\sphinxAtStartPar
\sphinxstylestrong{COLOMBIA}
&
\sphinxAtStartPar
???
\\
\sphinxbottomrule
\end{tabulary}
\sphinxtableafterendhook\par
\sphinxattableend\end{savenotes}


\bigskip\hrule\bigskip



\paragraph{Herramientas de LibreOffice}
\label{\detokenize{cursos/ofimatica10/recursos/T03:herramientas-de-libreoffice}}
\sphinxAtStartPar
\sphinxfootnotemark[1] A continuación se explican las diferentes herramientas de LibreOffice:


\subparagraph{LibreOffice Writer (procesador de textos)}
\label{\detokenize{cursos/ofimatica10/recursos/T03:libreoffice-writer-procesador-de-textos}}

\begin{savenotes}\sphinxattablestart
\sphinxthistablewithglobalstyle
\centering
\begin{tabulary}{\linewidth}[t]{T}
\sphinxtoprule
\sphinxtableatstartofbodyhook

\begin{sphinxfigure-in-table}
\centering
\capstart
\noindent\sphinxincludegraphics[width=270\sphinxpxdimen]{{image16}.png}
\sphinxfigcaption{LibreOffice Writer}\label{\detokenize{cursos/ofimatica10/recursos/T03:id8}}\end{sphinxfigure-in-table}\relax
\\
\sphinxbottomrule
\end{tabulary}
\sphinxtableafterendhook\par
\sphinxattableend\end{savenotes}

\sphinxAtStartPar
\sphinxstylestrong{Writer} es una herramienta rica en funciones para crear cartas, libros, informes, boletines, folletos y otros documentos. Puede insertar gráficos y objetos de otros componentes de LibreOffice en documentos de Writer.

\sphinxAtStartPar
\sphinxstylestrong{Writer} puede exportar archivos a HTML, XHTML, XML, PDF y EPUB, puede guardar archivos en muchos formatos incluidas varias versiones de archivo de Microsoft Word y conectarse con su cliente de correo electrónico.


\subparagraph{LibreOffice Calc (hoja de cálculo)}
\label{\detokenize{cursos/ofimatica10/recursos/T03:libreoffice-calc-hoja-de-calculo}}

\begin{savenotes}\sphinxattablestart
\sphinxthistablewithglobalstyle
\centering
\begin{tabulary}{\linewidth}[t]{T}
\sphinxtoprule
\sphinxtableatstartofbodyhook

\begin{sphinxfigure-in-table}
\centering
\capstart
\noindent\sphinxincludegraphics[width=190\sphinxpxdimen]{{image17}.png}
\sphinxfigcaption{LibreOffice Calc}\label{\detokenize{cursos/ofimatica10/recursos/T03:id9}}\end{sphinxfigure-in-table}\relax
\\
\sphinxbottomrule
\end{tabulary}
\sphinxtableafterendhook\par
\sphinxattableend\end{savenotes}

\sphinxAtStartPar
\sphinxstylestrong{Calc} tiene todas las funciones avanzadas de análisis, gráficos y toma de decisiones que se esperan de una hoja de cálculo de alta gama. Incluye más de 500 funciones para operaciones financieras, estadísticas y matemáticas, entre otras. El Gestor de escenarios proporciona un análisis de «qué pasaría si».

\sphinxAtStartPar
\sphinxstylestrong{Calc} genera gráficos en 2D y 3D, que se pueden integrar en otros documentos de LibreOffice. También puede abrir y trabajar y guardar hojas de cálculo de Microsoft Excel y otros formatos de hoja de cálculo, así como exportar hojas de cálculo a varios formatos, como archivos de valores separados por comas (CSV), PDF y HTML.


\subparagraph{LibreOffice Impress (presentaciones)}
\label{\detokenize{cursos/ofimatica10/recursos/T03:libreoffice-impress-presentaciones}}

\begin{savenotes}\sphinxattablestart
\sphinxthistablewithglobalstyle
\centering
\begin{tabulary}{\linewidth}[t]{T}
\sphinxtoprule
\sphinxtableatstartofbodyhook

\begin{sphinxfigure-in-table}
\centering
\capstart
\noindent\sphinxincludegraphics[width=270\sphinxpxdimen]{{image18}.png}
\sphinxfigcaption{LibreOffice Impress}\label{\detokenize{cursos/ofimatica10/recursos/T03:id10}}\end{sphinxfigure-in-table}\relax
\\
\sphinxbottomrule
\end{tabulary}
\sphinxtableafterendhook\par
\sphinxattableend\end{savenotes}

\sphinxAtStartPar
\sphinxstylestrong{Impress} proporciona todas las herramientas de presentación multimedia comunes, como efectos especiales, animaciones y herramientas de dibujo. Está integrado con las capacidades gráficas avanzadas de los componentes de \sphinxstylestrong{LibreOffice Draw} y \sphinxstylestrong{Math}.

\sphinxAtStartPar
Las presentaciones de diapositivas se pueden mejorar utilizando texto de efectos especiales de Fontwork, así como clips de sonido y vídeo. Impress puede abrir, editar y guardar presentaciones de Microsoft PowerPoint y también puede guardar su trabajo en numerosos formatos gráficos.


\subparagraph{LibreOffice Draw (Dibujo vectorial)}
\label{\detokenize{cursos/ofimatica10/recursos/T03:libreoffice-draw-dibujo-vectorial}}

\begin{savenotes}\sphinxattablestart
\sphinxthistablewithglobalstyle
\centering
\begin{tabulary}{\linewidth}[t]{T}
\sphinxtoprule
\sphinxtableatstartofbodyhook

\begin{sphinxfigure-in-table}
\centering
\capstart
\noindent\sphinxincludegraphics[width=230\sphinxpxdimen]{{image19}.png}
\sphinxfigcaption{LibreOffice Draw}\label{\detokenize{cursos/ofimatica10/recursos/T03:id11}}\end{sphinxfigure-in-table}\relax
\\
\sphinxbottomrule
\end{tabulary}
\sphinxtableafterendhook\par
\sphinxattableend\end{savenotes}

\sphinxAtStartPar
\sphinxstylestrong{Draw} es una herramienta de dibujo vectorial que puede producir desde simples diagramas o diagramas de flujo hasta ilustraciones en 3D. Su característica Conectores Inteligentes le permite definir sus propios puntos de conexión. Puede usar \sphinxstylestrong{Draw} para crear dibujos y emplearlos en cualquiera de los componentes de LibreOffice, y puede crear su propia imagen prediseñada para luego agregarla a la Galería.

\sphinxAtStartPar
\sphinxstylestrong{Draw} puede importar archivos gráficos en muchos formatos y guardar su trabajo en muchos formatos, incluidos PNG, GIF, JPEG, BMP, TIFF, SVG, HTML, PDF y WebP.


\subparagraph{LibreOffice Base (Bases de datos)}
\label{\detokenize{cursos/ofimatica10/recursos/T03:libreoffice-base-bases-de-datos}}

\begin{savenotes}\sphinxattablestart
\sphinxthistablewithglobalstyle
\centering
\begin{tabulary}{\linewidth}[t]{T}
\sphinxtoprule
\sphinxtableatstartofbodyhook

\begin{sphinxfigure-in-table}
\centering
\capstart
\noindent\sphinxincludegraphics[width=270\sphinxpxdimen]{{image20}.png}
\sphinxfigcaption{LibreOffice Base}\label{\detokenize{cursos/ofimatica10/recursos/T03:id12}}\end{sphinxfigure-in-table}\relax
\\
\sphinxbottomrule
\end{tabulary}
\sphinxtableafterendhook\par
\sphinxattableend\end{savenotes}

\sphinxAtStartPar
\sphinxstylestrong{Base} proporciona herramientas para el trabajo diario de bases de datos dentro de una interfaz sencilla. Puede crear y editar formularios, informes, consultas, tablas, vistas y relaciones, por lo que administrar una base de datos relacional es muy similar a otras aplicaciones de bases de datos populares.

\sphinxAtStartPar
\sphinxstylestrong{Base} proporciona muchas características nuevas, como la capacidad de analizar y editar relaciones desde una vista de diagramas. Base incorpora dos motores de bases de datos relacionales, HSQLDB y Firebird. También puede utilizar PostgreSQL, dBASE, Microsoft Access, MySQL, Oracle o cualquier base de datos compatible con ODBC o JDBC. Base también proporciona soporte para un subconjunto de ANSI\sphinxhyphen{}92 SQL.


\subparagraph{LibreOffice Math (Editor de fórmulas)}
\label{\detokenize{cursos/ofimatica10/recursos/T03:libreoffice-math-editor-de-formulas}}

\begin{savenotes}\sphinxattablestart
\sphinxthistablewithglobalstyle
\centering
\begin{tabulary}{\linewidth}[t]{T}
\sphinxtoprule
\sphinxtableatstartofbodyhook

\begin{sphinxfigure-in-table}
\centering
\capstart
\noindent\sphinxincludegraphics[width=240\sphinxpxdimen]{{image21}.png}
\sphinxfigcaption{LibreOffice Math}\label{\detokenize{cursos/ofimatica10/recursos/T03:id13}}\end{sphinxfigure-in-table}\relax
\\
\sphinxbottomrule
\end{tabulary}
\sphinxtableafterendhook\par
\sphinxattableend\end{savenotes}

\sphinxAtStartPar
Math es un editor de fórmulas o ecuaciones. Puede usarlo para crear ecuaciones complejas que incluyen símbolos o caracteres que no están disponibles en conjuntos de fuentes estándar. Si bien se usa más comúnmente para crear fórmulas en otros documentos, como archivos Writer e Impress, Math también puede funcionar como una herramienta independiente. Puede guardar fórmulas en el formato estándar de \sphinxstylestrong{Lenguaje de Marcado Matemático (MathML)} para incluirlas en páginas web u otros documentos creados con otros programas.


\bigskip\hrule\bigskip



\paragraph{Ventajas de LibreOffice}
\label{\detokenize{cursos/ofimatica10/recursos/T03:ventajas-de-libreoffice}}
\sphinxAtStartPar
\sphinxfootnotemark[1] A continuación se presentan las ventajas que tiene el uso de LibreOffice frente a otras suites ofimáticas:
\begin{itemize}
\item {} 
\sphinxAtStartPar
\sphinxstylestrong{Sin tarifas de licencia:} LibreOffice es gratuito para que cualquiera lo use y distribuya sin coste alguno. LibreOffice incluye muchas funciones (como exportación a PDF) de manera gratuita y que en otras suites de oficina son de pago. Con LibreOffice no hay cargos ocultos ni ahora, ni en el futuro.

\item {} 
\sphinxAtStartPar
\sphinxstylestrong{Código abierto:} Puede distribuir, copiar y modificar el software tanto como desee, de acuerdo con las licencias de código abierto de LibreOffice.

\item {} 
\sphinxAtStartPar
\sphinxstylestrong{Multiplataforma:} LibreOffice se ejecuta en varias arquitecturas de hardware y en varios sistemas operativos, como Microsoft Windows, macOS y Linux.

\item {} 
\sphinxAtStartPar
\sphinxstylestrong{Amplio soporte de idiomas:} La interfaz de usuario de LibreOffice, diccionarios ortográficos, separación de sílabas y sinónimos, están disponibles en más de 100 idiomas y dialectos. LibreOffice también ofrece compatibilidad con los idiomas de diseño de texto complejo (CTL) y de escritura de derecha a izquierda (RTL) (como urdu, hebreo y árabe).

\item {} 
\sphinxAtStartPar
\sphinxstylestrong{Interfaz de usuario consistente:} Todos los componentes tienen una apariencia similar, lo que los hace fáciles de usar y aprender.

\item {} 
\sphinxAtStartPar
\sphinxstylestrong{Integración:} Los componentes de LibreOffice están integrados entre sí, todos comparten un corrector ortográfico común y otras herramientas, que se utilizan de forma coherente en toda la suite. Por ejemplo, las herramientas de dibujo disponibles en Writer también se encuentran en Calc, y con versiones similares pero mejoradas en Impress y Draw.
\begin{itemize}
\item {} 
\sphinxAtStartPar
No necesita saber qué aplicación se utilizó para crear un archivo en particular. Por ejemplo, puede abrir un archivo Draw desde Writer; se abrirá automáticamente en Draw.

\end{itemize}

\item {} 
\sphinxAtStartPar
\sphinxstylestrong{Granularidad:} Por lo general, si cambia una opción, afecta a todos los componentes. Sin embargo, las opciones de LibreOffice se pueden configurar en el ámbito de componente o incluso en el ámbito de documento.

\item {} 
\sphinxAtStartPar
\sphinxstylestrong{Compatibilidad de archivos:} Además de sus formatos nativos (OpenDocument), LibreOffice puede y guardar archivos en muchos formatos comunes, entre otros: Microsoft Office, HTML, XML, WordPerfect, Lotus 1\sphinxhyphen{}2\sphinxhyphen{}3 y PDF.

\item {} 
\sphinxAtStartPar
\sphinxstylestrong{Sin bloqueo de proveedores:} LibreOffice utiliza OpenDocument, un formato de archivo XML (eXtensible Markup Language) desarrollado como estándar industrial por OASIS (Organization for the Advancement of Structured Information Standards). Estos archivos se pueden descomprimir y leer fácilmente con cualquier editor de texto, y su estructura es abierta y está documentada.

\item {} 
\sphinxAtStartPar
\sphinxstylestrong{Usted tiene voz:} Las mejoras, las correcciones de software y las fechas de lanzamiento están impulsadas por la comunidad. Puede unirse a la comunidad e influir sobre el curso de LibreOffice.

\end{itemize}


\bigskip\hrule\bigskip


\sphinxstepscope


\chapter{Ofimática 11}
\label{\detokenize{cursos/ofimatica11/index:ofimatica-11}}\label{\detokenize{cursos/ofimatica11/index::doc}}
\sphinxAtStartPar
Bienvenido al curso de ofimática de grado once.

\sphinxAtStartPar
Acceda a las guías didácticas según el periodo.

\sphinxstepscope


\section{Periodo 01}
\label{\detokenize{cursos/ofimatica11/periodo01/index:periodo-01}}\label{\detokenize{cursos/ofimatica11/periodo01/index::doc}}\begin{quote}\begin{description}
\sphinxlineitem{DESEMPEÑO}
\sphinxAtStartPar
Diferencia las partes físicas de la computadora de sus programas y comprende cómo la CPU y la memoria RAM trabajan con los datos para generar información útil. Identifica la organización jerárquica de la información y utiliza correctamente los términos hardware, software y programa en contextos técnicos sencillos.

\end{description}\end{quote}


\bigskip\hrule\bigskip



\subsection{Contenidos}
\label{\detokenize{cursos/ofimatica11/periodo01/index:contenidos}}
\sphinxstepscope


\subsubsection{Guía 01: Conocimientos previos}
\label{\detokenize{cursos/ofimatica11/periodo01/G01:guia-01-conocimientos-previos}}\label{\detokenize{cursos/ofimatica11/periodo01/G01::doc}}

\begin{savenotes}\sphinxattablestart
\sphinxthistablewithglobalstyle
\centering
\begin{tabular}[t]{\X{35}{100}\X{65}{100}}
\sphinxtoprule
\sphinxtableatstartofbodyhook

\begin{savenotes}\sphinxattablestart
\sphinxthistablewithglobalstyle
\centering
\begin{tabulary}{\linewidth}[t]{T}
\sphinxtoprule
\sphinxtableatstartofbodyhook
\sphinxAtStartPar
\sphinxincludegraphics{{escudo}.png}
\\
\sphinxbottomrule
\end{tabulary}
\sphinxtableafterendhook\par
\sphinxattableend\end{savenotes}
&

\begin{savenotes}\sphinxattablestart
\sphinxthistablewithglobalstyle
\centering
\begin{tabulary}{\linewidth}[t]{T}
\sphinxtoprule
\sphinxtableatstartofbodyhook
\sphinxAtStartPar
\sphinxstylestrong{INSTITUCIÓN EDUCATIVA}: John F. Kennedy
\\
\sphinxhline
\sphinxAtStartPar
\sphinxstylestrong{DOCENTE}: Julian Camilo Fonseca Romero
\\
\sphinxhline
\sphinxAtStartPar
\sphinxstylestrong{ASIGNATURA}: Ofimática
\\
\sphinxhline
\sphinxAtStartPar
\sphinxstylestrong{GRADO}: ONCE
\\
\sphinxhline
\sphinxAtStartPar
\sphinxstylestrong{TIEMPO}: 10 semanas (4 horas / semana)
\\
\sphinxhline
\sphinxAtStartPar
\sphinxstylestrong{TIEMPO DE TRABAJO}: 10/40 horas
\\
\sphinxbottomrule
\end{tabulary}
\sphinxtableafterendhook\par
\sphinxattableend\end{savenotes}
\\
\sphinxbottomrule
\end{tabular}
\sphinxtableafterendhook\par
\sphinxattableend\end{savenotes}

\begin{sphinxuseclass}{sd-tab-set}
\begin{sphinxuseclass}{sd-tab-item}\subsubsection*{Objetivo}

\begin{sphinxuseclass}{sd-tab-content}\begin{itemize}
\item {} 
\sphinxAtStartPar
\sphinxstylestrong{Temática}: Concepto de hardware, software, diferencias entre software y programa, información, datos y jerarquía de datos.

\item {} 
\sphinxAtStartPar
\sphinxstylestrong{Objetivo de aprendizaje}: Entender los fundamentos de arquitectura y lógica computacional.

\end{itemize}

\end{sphinxuseclass}
\end{sphinxuseclass}
\begin{sphinxuseclass}{sd-tab-item}\subsubsection*{Orientaciones curriculares}

\begin{sphinxuseclass}{sd-tab-content}\begin{itemize}
\item {} 
\sphinxAtStartPar
\sphinxstylestrong{Componente}: Naturaleza y Evolución de la T\&I.

\item {} 
\sphinxAtStartPar
\sphinxstylestrong{Competencia}: Relaciono saberes, conocimientos tecnológicos e informáticos con los conocimientos de otras disciplinas.

\item {} 
\sphinxAtStartPar
\sphinxstylestrong{Evidencia de aprendizaje}: Comprendo los principios y conocimientos tecnológicos e informáticos que hacen posible el funcionamiento de productos tecnológicos actuales.

\end{itemize}

\end{sphinxuseclass}
\end{sphinxuseclass}
\begin{sphinxuseclass}{sd-tab-item}\subsubsection*{Criterios de evaluación}

\begin{sphinxuseclass}{sd-tab-content}\begin{itemize}
\item {} 
\sphinxAtStartPar
Actividades desarrolladas en su totalidad.

\item {} 
\sphinxAtStartPar
Participación en clase.

\item {} 
\sphinxAtStartPar
Explicación clara de conceptos.

\item {} 
\sphinxAtStartPar
Argumentación de ideas.

\item {} 
\sphinxAtStartPar
Cumplimiento de las reglas del aula.

\end{itemize}

\end{sphinxuseclass}
\end{sphinxuseclass}
\begin{sphinxuseclass}{sd-tab-item}\subsubsection*{Recursos (1)}

\begin{sphinxuseclass}{sd-tab-content}\begin{itemize}
\item {} 
\sphinxAtStartPar
{\hyperref[\detokenize{cursos/ofimatica11/recursos/T01::doc}]{\sphinxcrossref{\DUrole{std}{\DUrole{std-doc}{TEXTO: Conocimientos previos}}}}}

\end{itemize}

\end{sphinxuseclass}
\end{sphinxuseclass}
\end{sphinxuseclass}

\paragraph{Mensaje del autor}
\label{\detokenize{cursos/ofimatica11/periodo01/G01:mensaje-del-autor}}\begin{quote}
\begin{quote}

\sphinxAtStartPar
“En esta guía, el estudiante recuerda conceptos básicos de informática para aplicar más adelante en el aprendizaje de la programación.”
\end{quote}

\begin{flushright}
---Camilo Fonseca
\end{flushright}
\end{quote}


\bigskip\hrule\bigskip



\paragraph{A00: Introducción (+0,1)}
\label{\detokenize{cursos/ofimatica11/periodo01/G01:a00-introduccion-0-1}}
\sphinxAtStartPar
\DUrole{sd-sphinx-override}{\DUrole{sd-badge}{\DUrole{sd-bg-light}{\DUrole{sd-bg-text-light}{Exploración}}}}
\begin{enumerate}
\sphinxsetlistlabels{\arabic}{enumi}{enumii}{}{.}%
\item {} 
\sphinxAtStartPar
TRANSCRIBA en su cuaderno el objetivo, las orientaciones curriculares y los criterios de evaluación de la presente guía.

\item {} 
\sphinxAtStartPar
LEA el \sphinxstylestrong{mensaje del autor} y en su cuaderno responda:
\begin{itemize}
\item {} 
\sphinxAtStartPar
¿Es posible entender la tecnología actual sin conocer como funciona el hardware y como se organizan los datos?

\item {} 
\sphinxAtStartPar
¿Qué es para usted un hardware?

\item {} 
\sphinxAtStartPar
¿Qué es para usted un software?

\item {} 
\sphinxAtStartPar
¿Qué es para usted información y en que se diferencia con los datos?

\item {} 
\sphinxAtStartPar
¿Cómo se relacionan los datos con el hardware y el software?

\end{itemize}

\end{enumerate}


\bigskip\hrule\bigskip



\paragraph{A01: Datos, información y jerarquía de datos (+0,1)}
\label{\detokenize{cursos/ofimatica11/periodo01/G01:a01-datos-informacion-y-jerarquia-de-datos-0-1}}
\sphinxAtStartPar
\DUrole{sd-sphinx-override}{\DUrole{sd-badge}{\DUrole{sd-bg-info}{\DUrole{sd-bg-text-info}{Estructuración}}}}
\begin{enumerate}
\sphinxsetlistlabels{\arabic}{enumi}{enumii}{}{.}%
\item {} 
\sphinxAtStartPar
En su cuaderno responda las siguientes preguntas:
\begin{itemize}
\item {} 
\sphinxAtStartPar
¿Qué son los datos?

\item {} 
\sphinxAtStartPar
¿Qué es información?

\item {} 
\sphinxAtStartPar
¿Cuál es la diferencia entre un valor cualitativo y un valor cuantitativo?

\item {} 
\sphinxAtStartPar
¿Qué es una jerarquía de datos?

\end{itemize}

\item {} 
\sphinxAtStartPar
En su cuaderno describa cada uno de los componentes de la jerarquía de datos y de un ejemplo de cada uno.

\end{enumerate}


\bigskip\hrule\bigskip



\paragraph{A02: Relación entre datos e información (+0,2)}
\label{\detokenize{cursos/ofimatica11/periodo01/G01:a02-relacion-entre-datos-e-informacion-0-2}}
\sphinxAtStartPar
\DUrole{sd-sphinx-override}{\DUrole{sd-badge}{\DUrole{sd-bg-info}{\DUrole{sd-bg-text-info}{Estructuración}}}}
\begin{enumerate}
\sphinxsetlistlabels{\arabic}{enumi}{enumii}{}{.}%
\item {} 
\sphinxAtStartPar
Responda las siguientes preguntas en su cuaderno:
\begin{itemize}
\item {} 
\sphinxAtStartPar
¿Qué es tener sentido?

\item {} 
\sphinxAtStartPar
¿Por qué es importante tener “sentido” para diferenciar entre datos e información?

\end{itemize}

\item {} 
\sphinxAtStartPar
Dibuje una tabla en su cuaderno y en la primera columna escriba 15 ejemplos de datos. En la segunda columna convierta estos datos en información. En la tercera columna indique si la información es un valor cualitativo o cuantitativo.


\begin{savenotes}\sphinxattablestart
\sphinxthistablewithglobalstyle
\centering
\begin{tabulary}{\linewidth}[t]{TTT}
\sphinxtoprule
\sphinxstyletheadfamily 
\sphinxAtStartPar
Datos
&\sphinxstyletheadfamily 
\sphinxAtStartPar
Información
&\sphinxstyletheadfamily 
\sphinxAtStartPar
Valor
\\
\sphinxmidrule
\sphinxtableatstartofbodyhook
\sphinxAtStartPar
Rojo
&
\sphinxAtStartPar
“Semáforo en estado rojo, carros no pueden pasar”
&
\sphinxAtStartPar
cualitativo
\\
\sphinxhline
\sphinxAtStartPar
16
&
\sphinxAtStartPar
“La edad del estudiante es 16 años, por tanto es menor de edad”
&
\sphinxAtStartPar
cuantitativo
\\
\sphinxbottomrule
\end{tabulary}
\sphinxtableafterendhook\par
\sphinxattableend\end{savenotes}

\end{enumerate}


\bigskip\hrule\bigskip



\paragraph{A03: Jerarquía de datos (+0,2)}
\label{\detokenize{cursos/ofimatica11/periodo01/G01:a03-jerarquia-de-datos-0-2}}
\sphinxAtStartPar
\DUrole{sd-sphinx-override}{\DUrole{sd-badge}{\DUrole{sd-bg-info}{\DUrole{sd-bg-text-info}{Estructuración}}}}
\begin{enumerate}
\sphinxsetlistlabels{\arabic}{enumi}{enumii}{}{.}%
\item {} 
\sphinxAtStartPar
Dibuje los siguientes círculos, ingrese ejemplos de datos según su jerarquía (ver la sección de \sphinxstylestrong{jerarquía de datos}).

\end{enumerate}


\begin{savenotes}\sphinxattablestart
\sphinxthistablewithglobalstyle
\centering
\begin{tabulary}{\linewidth}[t]{T}
\sphinxtoprule
\sphinxtableatstartofbodyhook

\begin{sphinxfigure-in-table}
\centering
\capstart
\noindent\sphinxincludegraphics[width=400\sphinxpxdimen]{{image015}.png}
\sphinxfigcaption{\sphinxstylestrong{Diagrama:} Jerarquía de datos}\label{\detokenize{cursos/ofimatica11/periodo01/G01:id1}}\end{sphinxfigure-in-table}\relax
\\
\sphinxbottomrule
\end{tabulary}
\sphinxtableafterendhook\par
\sphinxattableend\end{savenotes}


\bigskip\hrule\bigskip



\paragraph{A04: Concepto de hardware (RAM y CPU) (+0,1)}
\label{\detokenize{cursos/ofimatica11/periodo01/G01:a04-concepto-de-hardware-ram-y-cpu-0-1}}
\sphinxAtStartPar
\DUrole{sd-sphinx-override}{\DUrole{sd-badge}{\DUrole{sd-bg-info}{\DUrole{sd-bg-text-info}{Estructuración}}}}
\begin{enumerate}
\sphinxsetlistlabels{\arabic}{enumi}{enumii}{}{.}%
\item {} 
\sphinxAtStartPar
En su cuaderno responda:
\begin{itemize}
\item {} 
\sphinxAtStartPar
¿Qué es hardware?

\item {} 
\sphinxAtStartPar
¿Cómo se puede clasificar el hardware de un sistema computacional?

\end{itemize}

\item {} 
\sphinxAtStartPar
En su cuaderno mencione diferentes ejemplos de hardware y clasifíquelos.

\item {} 
\sphinxAtStartPar
En su cuaderno responda:
\begin{itemize}
\item {} 
\sphinxAtStartPar
¿Qué es la memoria RAM?

\item {} 
\sphinxAtStartPar
¿Cuales son las características de la memoria RAM?

\end{itemize}

\item {} 
\sphinxAtStartPar
En INTERNET consulte 3 fabricantes de memorias RAM y escríbalos en su cuaderno. Tenga en cuenta lo siguiente:
\begin{itemize}
\item {} 
\sphinxAtStartPar
Mencionar la nacionalidad del fabricante.

\item {} 
\sphinxAtStartPar
Mencionar los clientes del fabricante.

\end{itemize}

\item {} 
\sphinxAtStartPar
En su cuaderno responda:
\begin{itemize}
\item {} 
\sphinxAtStartPar
¿Qué es la CPU?

\item {} 
\sphinxAtStartPar
¿Qué son los Hertz(Hz)?

\item {} 
\sphinxAtStartPar
La CPU tiene 3 componentes. En su cuaderno complete el siguiente gráfico:


\begin{savenotes}\sphinxattablestart
\sphinxthistablewithglobalstyle
\centering
\begin{tabulary}{\linewidth}[t]{T}
\sphinxtoprule
\sphinxtableatstartofbodyhook

\begin{sphinxfigure-in-table}
\centering
\capstart
\noindent\sphinxincludegraphics[width=400\sphinxpxdimen]{{image024}.png}
\sphinxfigcaption{\sphinxstylestrong{Diagrama:} Componentes CPU}\label{\detokenize{cursos/ofimatica11/periodo01/G01:id2}}\end{sphinxfigure-in-table}\relax
\\
\sphinxbottomrule
\end{tabulary}
\sphinxtableafterendhook\par
\sphinxattableend\end{savenotes}

\end{itemize}

\end{enumerate}


\bigskip\hrule\bigskip



\paragraph{A05: Utilidad del hardware (+0,2)}
\label{\detokenize{cursos/ofimatica11/periodo01/G01:a05-utilidad-del-hardware-0-2}}
\sphinxAtStartPar
\DUrole{sd-sphinx-override}{\DUrole{sd-badge}{\DUrole{sd-bg-info}{\DUrole{sd-bg-text-info}{Estructuración}}}}
\begin{enumerate}
\sphinxsetlistlabels{\arabic}{enumi}{enumii}{}{.}%
\item {} 
\sphinxAtStartPar
Teniendo en cuenta la siguiente imagen, responda en su cuaderno la siguiente pregunta:
\begin{itemize}
\item {} 
\sphinxAtStartPar
¿En que se diferencian los dispositivos de entrada con los dispositivos de salida?

\end{itemize}


\begin{savenotes}\sphinxattablestart
\sphinxthistablewithglobalstyle
\centering
\begin{tabulary}{\linewidth}[t]{T}
\sphinxtoprule
\sphinxtableatstartofbodyhook

\begin{sphinxfigure-in-table}
\centering
\capstart
\noindent\sphinxincludegraphics[width=360\sphinxpxdimen]{{image033}.png}
\sphinxfigcaption{\sphinxstylestrong{Diagrama:} CPU}\label{\detokenize{cursos/ofimatica11/periodo01/G01:id3}}\end{sphinxfigure-in-table}\relax
\\
\sphinxbottomrule
\end{tabulary}
\sphinxtableafterendhook\par
\sphinxattableend\end{savenotes}

\item {} 
\sphinxAtStartPar
En media hoja de cuaderno, explique por que la CPU y la memoria RAM son los dispositivos más importantes para dar instrucciones.

\item {} 
\sphinxAtStartPar
Respecto a la memoria \sphinxstylestrong{RAM}, en las siguientes frases diga si es falso o verdadero y explique el por qué:
\begin{itemize}
\item {} 
\sphinxAtStartPar
La memoria \sphinxstylestrong{RAM} almacena infinitamente los datos y programas \_\_\_

\item {} 
\sphinxAtStartPar
Para acceder velozmente a la información en el equipo usamos la memoria \sphinxstylestrong{RAM} \_\_\_

\item {} 
\sphinxAtStartPar
La memoriam \sphinxstylestrong{RAM} no es volátil \_\_\_

\item {} 
\sphinxAtStartPar
La velocidad de trabajo de un computador, depende exclusivamente de la velocidad de la memoria \sphinxstylestrong{RAM} \_\_\_

\end{itemize}

\end{enumerate}


\bigskip\hrule\bigskip



\paragraph{A06: Concepto de software (+0,1)}
\label{\detokenize{cursos/ofimatica11/periodo01/G01:a06-concepto-de-software-0-1}}
\sphinxAtStartPar
\DUrole{sd-sphinx-override}{\DUrole{sd-badge}{\DUrole{sd-bg-info}{\DUrole{sd-bg-text-info}{Estructuración}}}}
\begin{enumerate}
\sphinxsetlistlabels{\arabic}{enumi}{enumii}{}{.}%
\item {} 
\sphinxAtStartPar
En su cuaderno responda la siguientes preguntas:
\begin{itemize}
\item {} 
\sphinxAtStartPar
¿Qué es el software?

\item {} 
\sphinxAtStartPar
¿Cómo se clasifica el software?

\item {} 
\sphinxAtStartPar
¿Cuál es la diferencia entre un software y un programa?

\end{itemize}

\end{enumerate}


\bigskip\hrule\bigskip



\paragraph{A07: Clasificación del software (+0,2)}
\label{\detokenize{cursos/ofimatica11/periodo01/G01:a07-clasificacion-del-software-0-2}}
\sphinxAtStartPar
\DUrole{sd-sphinx-override}{\DUrole{sd-badge}{\DUrole{sd-bg-info}{\DUrole{sd-bg-text-info}{Estructuración}}}}
\begin{enumerate}
\sphinxsetlistlabels{\arabic}{enumi}{enumii}{}{.}%
\item {} 
\sphinxAtStartPar
Clasifique el siguiente software en la siguiente tabla:
\begin{itemize}
\item {} 
\sphinxAtStartPar
Windows 11, WORD, cpa000.exe, Linux UBUNTU, LibreOffice Calc, UNIX, hii.bin, ANDROID, FIREFOX, LibreOffice Writer, windowSys.exe, IOS, CHROME, FIFA, Photoshop, BIOS

\end{itemize}


\begin{savenotes}\sphinxattablestart
\sphinxthistablewithglobalstyle
\centering
\begin{tabulary}{\linewidth}[t]{TT}
\sphinxtoprule
\sphinxstyletheadfamily 
\sphinxAtStartPar
Software de sistema
&\sphinxstyletheadfamily 
\sphinxAtStartPar
Software de aplicación
\\
\sphinxmidrule
\sphinxtableatstartofbodyhook
\sphinxAtStartPar
\_
&
\sphinxAtStartPar

\\
\sphinxhline
\sphinxAtStartPar
\_
&
\sphinxAtStartPar

\\
\sphinxhline
\sphinxAtStartPar
\_
&
\sphinxAtStartPar

\\
\sphinxhline
\sphinxAtStartPar
\_
&
\sphinxAtStartPar

\\
\sphinxhline
\sphinxAtStartPar
\_
&
\sphinxAtStartPar

\\
\sphinxbottomrule
\end{tabulary}
\sphinxtableafterendhook\par
\sphinxattableend\end{savenotes}

\end{enumerate}


\bigskip\hrule\bigskip



\paragraph{A08: Examen final (+4,0)}
\label{\detokenize{cursos/ofimatica11/periodo01/G01:a08-examen-final-4-0}}
\sphinxAtStartPar
\DUrole{sd-sphinx-override}{\DUrole{sd-badge}{\DUrole{sd-bg-danger}{\DUrole{sd-bg-text-danger}{Transferencia}}}}
\begin{quote}\begin{description}
\sphinxlineitem{TIEMPO DEL EXAMEN}
\sphinxAtStartPar
30 minutos

\sphinxlineitem{NÚMERO DE PREGUNTAS}
\sphinxAtStartPar
10

\end{description}\end{quote}

\begin{sphinxadmonition}{note}{Posibles preguntas}
\begin{enumerate}
\sphinxsetlistlabels{\arabic}{enumi}{enumii}{}{.}%
\item {} 
\sphinxAtStartPar
Explique la diferencia entre datos e información y de un ejemplo de cada uno.

\item {} 
\sphinxAtStartPar
¿Por qué es importante que los datos tengan un contexto para convertirse en información útil?

\item {} 
\sphinxAtStartPar
Describa qué son los valores cualitativos y cuantitativos y proporcione un ejemplo de cada uno.

\item {} 
\sphinxAtStartPar
Explique la jerarquía de datos desde el bit hasta la base de datos, indicando qué representa cada nivel.

\item {} 
\sphinxAtStartPar
¿Por qué un byte está compuesto por 8 bits y qué importancia tiene esta unidad en la informática?

\item {} 
\sphinxAtStartPar
Define qué es un registro y cómo se relaciona con los campos y archivos en la organización de datos.

\item {} 
\sphinxAtStartPar
Explique qué es el hardware y mencione al menos cinco componentes físicos que lo conforman.

\item {} 
\sphinxAtStartPar
Describa la función principal de la memoria RAM y por qué se considera una memoria volátil.

\item {} 
\sphinxAtStartPar
¿Cuál es el papel de la CPU en un sistema computacional y cuáles son sus componentes principales?

\item {} 
\sphinxAtStartPar
Explique la diferencia entre software de sistema y software de aplicación con ejemplos.

\item {} 
\sphinxAtStartPar
¿Qué es un programa y cómo se relaciona con el software en general?

\item {} 
\sphinxAtStartPar
Describa cómo la información procesada por un computador puede influir en la toma de decisiones.

\item {} 
\sphinxAtStartPar
¿Por qué es fundamental entender la arquitectura de un computador para programar eficientemente?

\item {} 
\sphinxAtStartPar
Explique cómo los datos se almacenan y organizan en un sistema informático para facilitar su acceso.

\item {} 
\sphinxAtStartPar
Describa el proceso que sigue un computador para transformar datos en información.

\item {} 
\sphinxAtStartPar
¿Qué papel juegan los Hertz dentro de la CPU y cómo afecta el rendimiento del sistema?

\item {} 
\sphinxAtStartPar
Explique la importancia de la memoria RAM en la ejecución de programas y qué sucede cuando se apaga el computador.

\item {} 
\sphinxAtStartPar
Describa cómo el hardware y el software trabajan juntos para realizar tareas computacionales.

\item {} 
\sphinxAtStartPar
¿Qué diferencias existen entre los distintos tipos de memoria en un computador y para qué se usa cada una?

\item {} 
\sphinxAtStartPar
Reflexione sobre cómo la evolución de la arquitectura computacional ha impactado en el desarrollo de tecnologías actuales.

\end{enumerate}
\end{sphinxadmonition}

\sphinxstepscope


\subsubsection{Guía 02: Diagramas de flujo}
\label{\detokenize{cursos/ofimatica11/periodo01/G02:guia-02-diagramas-de-flujo}}\label{\detokenize{cursos/ofimatica11/periodo01/G02::doc}}

\begin{savenotes}\sphinxattablestart
\sphinxthistablewithglobalstyle
\centering
\begin{tabular}[t]{\X{35}{100}\X{65}{100}}
\sphinxtoprule
\sphinxtableatstartofbodyhook

\begin{savenotes}\sphinxattablestart
\sphinxthistablewithglobalstyle
\centering
\begin{tabulary}{\linewidth}[t]{T}
\sphinxtoprule
\sphinxtableatstartofbodyhook
\sphinxAtStartPar
\sphinxincludegraphics{{escudo}.png}
\\
\sphinxbottomrule
\end{tabulary}
\sphinxtableafterendhook\par
\sphinxattableend\end{savenotes}
&

\begin{savenotes}\sphinxattablestart
\sphinxthistablewithglobalstyle
\centering
\begin{tabulary}{\linewidth}[t]{T}
\sphinxtoprule
\sphinxtableatstartofbodyhook
\sphinxAtStartPar
\sphinxstylestrong{INSTITUCIÓN EDUCATIVA}: John F. Kennedy
\\
\sphinxhline
\sphinxAtStartPar
\sphinxstylestrong{DOCENTE}: Julian Camilo Fonseca Romero
\\
\sphinxhline
\sphinxAtStartPar
\sphinxstylestrong{ASIGNATURA}: Ofimática
\\
\sphinxhline
\sphinxAtStartPar
\sphinxstylestrong{GRADO}: ONCE
\\
\sphinxhline
\sphinxAtStartPar
\sphinxstylestrong{TIEMPO}: 10 semanas (4 horas / semana)
\\
\sphinxhline
\sphinxAtStartPar
\sphinxstylestrong{TIEMPO DE TRABAJO}: 10/40 horas
\\
\sphinxbottomrule
\end{tabulary}
\sphinxtableafterendhook\par
\sphinxattableend\end{savenotes}
\\
\sphinxbottomrule
\end{tabular}
\sphinxtableafterendhook\par
\sphinxattableend\end{savenotes}

\begin{sphinxuseclass}{sd-tab-set}
\begin{sphinxuseclass}{sd-tab-item}\subsubsection*{Objetivo}

\begin{sphinxuseclass}{sd-tab-content}\begin{itemize}
\item {} 
\sphinxAtStartPar
\sphinxstylestrong{Temática}: Diagramas de flujo, instrucciones.

\item {} 
\sphinxAtStartPar
\sphinxstylestrong{Objetivo de aprendizaje}: Reconocer y diferenciar la simbología estándar (ANSI/ISO) utilizada en los diagramas de flujo (inicio, proceso, decisión, entrada/salida).

\end{itemize}

\end{sphinxuseclass}
\end{sphinxuseclass}
\begin{sphinxuseclass}{sd-tab-item}\subsubsection*{Orientaciones curriculares}

\begin{sphinxuseclass}{sd-tab-content}\begin{itemize}
\item {} 
\sphinxAtStartPar
\sphinxstylestrong{Componente}: Solución de problemas con T\&I.

\item {} 
\sphinxAtStartPar
\sphinxstylestrong{Competencia}: Propongo innovaciones tecnológicas e informáticas para la solución de problemas dando cumplimiento a restricciones, condiciones y especificaciones técnicas y contextuales.

\item {} 
\sphinxAtStartPar
\sphinxstylestrong{Evidencia de aprendizaje}: Propongo, analizo y comparo diferentes soluciones a un mismo problema de la tecnología o la informática, explicando su origen, ventajas y dificultades.

\end{itemize}

\end{sphinxuseclass}
\end{sphinxuseclass}
\begin{sphinxuseclass}{sd-tab-item}\subsubsection*{Criterios de evaluación}

\begin{sphinxuseclass}{sd-tab-content}\begin{itemize}
\item {} 
\sphinxAtStartPar
Actividades desarrolladas en su totalidad.

\item {} 
\sphinxAtStartPar
Participación en clase.

\item {} 
\sphinxAtStartPar
Explicación clara de conceptos.

\item {} 
\sphinxAtStartPar
Argumentación de ideas.

\item {} 
\sphinxAtStartPar
Cumplimiento de las reglas del aula.

\end{itemize}

\end{sphinxuseclass}
\end{sphinxuseclass}
\begin{sphinxuseclass}{sd-tab-item}\subsubsection*{Recursos (1)}

\begin{sphinxuseclass}{sd-tab-content}\begin{itemize}
\item {} 
\sphinxAtStartPar
{\hyperref[\detokenize{cursos/ofimatica11/recursos/T02::doc}]{\sphinxcrossref{\DUrole{std}{\DUrole{std-doc}{TEXTO: Diagramas de flujo}}}}}

\end{itemize}

\end{sphinxuseclass}
\end{sphinxuseclass}
\end{sphinxuseclass}

\paragraph{Mensaje del autor}
\label{\detokenize{cursos/ofimatica11/periodo01/G02:mensaje-del-autor}}\begin{quote}
\begin{quote}

\sphinxAtStartPar
“En esta guía, el estudiante aprenderá que la simbología estándar (ANSI/ISO) es el lenguaje universal para modelar soluciones. Al dominar el uso de inicios, procesos y decisiones, será posible proponer y comparar diferentes caminos para resolver un mismo problema, y así analizar la ruta más eficiente.”
\end{quote}

\begin{flushright}
---Camilo Fonseca
\end{flushright}
\end{quote}


\bigskip\hrule\bigskip



\paragraph{A00: Introducción (+0,3)}
\label{\detokenize{cursos/ofimatica11/periodo01/G02:a00-introduccion-0-3}}
\sphinxAtStartPar
\DUrole{sd-sphinx-override}{\DUrole{sd-badge}{\DUrole{sd-bg-light}{\DUrole{sd-bg-text-light}{Exploración}}}}
\begin{enumerate}
\sphinxsetlistlabels{\arabic}{enumi}{enumii}{}{.}%
\item {} 
\sphinxAtStartPar
TRANSCRIBA en su cuaderno el objetivo, las orientaciones curriculares y los criterios de evaluación de la presente guía.

\item {} 
\sphinxAtStartPar
LEA el \sphinxstylestrong{mensaje del autor} y en su cuaderno responda:
\begin{itemize}
\item {} 
\sphinxAtStartPar
¿Qué es el estandar (ANSI/ISO)? (consulte en INTERNET)

\item {} 
\sphinxAtStartPar
¿Qué entiende por proceso?

\item {} 
\sphinxAtStartPar
¿Qué entiende por decisión?

\item {} 
\sphinxAtStartPar
Para usted, ¿Qué es ser eficiente?

\end{itemize}

\end{enumerate}


\bigskip\hrule\bigskip



\paragraph{A01: Diagramas de flujo (+0,7)}
\label{\detokenize{cursos/ofimatica11/periodo01/G02:a01-diagramas-de-flujo-0-7}}
\sphinxAtStartPar
\DUrole{sd-sphinx-override}{\DUrole{sd-badge}{\DUrole{sd-bg-info}{\DUrole{sd-bg-text-info}{Estructuración}}}}
\begin{enumerate}
\sphinxsetlistlabels{\arabic}{enumi}{enumii}{}{.}%
\item {} 
\sphinxAtStartPar
En su cuaderno responda las siguientes preguntas:
\begin{itemize}
\item {} 
\sphinxAtStartPar
¿Qué es una \sphinxstylestrong{instrucción}?

\item {} 
\sphinxAtStartPar
¿Por que la \sphinxstylestrong{claridad}, la \sphinxstylestrong{eficiencia} y la \sphinxstylestrong{seguridad} son importantes al momento de dar instrucciones?

\end{itemize}

\item {} 
\sphinxAtStartPar
En su cuaderno escriba una serie de \sphinxstylestrong{instrucciones} (seguras, claras y eficientes) para que un estudiante vaya al colegio desde que se despierta hasta que llega al aula de clase.

\item {} 
\sphinxAtStartPar
En su cuaderno describa cada una de las figuras de un \sphinxstylestrong{diagrama de flujo} y dibuje sus símbolos.

\end{enumerate}


\bigskip\hrule\bigskip



\paragraph{A02: Orientaciones al extranjero (+1,0)}
\label{\detokenize{cursos/ofimatica11/periodo01/G02:a02-orientaciones-al-extranjero-1-0}}
\sphinxAtStartPar
\DUrole{sd-sphinx-override}{\DUrole{sd-badge}{\DUrole{sd-bg-danger}{\DUrole{sd-bg-text-danger}{Transferencia}}}}
\begin{enumerate}
\sphinxsetlistlabels{\arabic}{enumi}{enumii}{}{.}%
\item {} 
\sphinxAtStartPar
Lea el siguiente caso:
\begin{quote}

\sphinxAtStartPar
“Suponga que un extranjero visita \sphinxstylestrong{Puerto Boyacá} y quiere hacer turismo en el pueblo.
El extranjero se encuentra en las afueras del hospital \sphinxstylestrong{JOSE CAYETANO VASQUEZ} y quiere conocer cada una de las \sphinxstylestrong{iglesias católicas} en Puerto Boyacá.”
\begin{itemize}
\item {} 
\sphinxAtStartPar
Iglesia San José Obrero (Barrio la Paz).

\item {} 
\sphinxAtStartPar
Iglesia San Pedro Claver (Centro).

\item {} 
\sphinxAtStartPar
Iglesia Cristo Rey (Cra 4a, Calle 5).

\end{itemize}
\end{quote}


\begin{savenotes}\sphinxattablestart
\sphinxthistablewithglobalstyle
\centering
\begin{tabulary}{\linewidth}[t]{T}
\sphinxtoprule
\sphinxtableatstartofbodyhook

\begin{sphinxfigure-in-table}
\centering
\capstart
\noindent\sphinxincludegraphics[width=650\sphinxpxdimen]{{image043}.png}
\sphinxfigcaption{Mapa Puerto Boyacá}\label{\detokenize{cursos/ofimatica11/periodo01/G02:id1}}\end{sphinxfigure-in-table}\relax
\\
\sphinxbottomrule
\end{tabulary}
\sphinxtableafterendhook\par
\sphinxattableend\end{savenotes}

\item {} 
\sphinxAtStartPar
En su \sphinxstylestrong{cuaderno}, cree un \sphinxstylestrong{diagrama de flujo} con 15 instrucciones que le permitan al extranjero llegar a cada una de las iglesias. Tenga en cuenta lo siguiente:
\begin{quote}
\begin{itemize}
\item {} 
\sphinxAtStartPar
5 instrucciones por iglesia (15 instrucciones).

\item {} 
\sphinxAtStartPar
2 direcciones por instrucción (30 direcciones).

\item {} 
\sphinxAtStartPar
Cree un inventario de las 30 direcciones por medio de variables así:

\begin{sphinxVerbatim}[commandchars=\\\{\}]
\PYG{n+nv}{d01}\PYG{+w}{ }\PYG{o}{=}\PYG{+w}{ }\PYG{l+s+s2}{\PYGZdq{}Hospital José Cayetano Vasquez, (Cra 5 con calle 29)\PYGZdq{}}
\PYG{n+nv}{d02}\PYG{+w}{ }\PYG{o}{=}\PYG{+w}{ }\PYG{l+s+s2}{\PYGZdq{}UPTC, (Cra 5 \PYGZsh{} 25 \PYGZhy{} 2)\PYGZdq{}}

\PYG{n+nv}{d10}\PYG{+w}{ }\PYG{o}{=}\PYG{+w}{ }\PYG{l+s+s2}{\PYGZdq{}Iglesia San José Obrero, (Calle 28 \PYGZsh{} 3A \PYGZhy{} 54)\PYGZdq{}}

\PYG{n+nv}{d20}\PYG{+w}{ }\PYG{o}{=}\PYG{+w}{ }\PYG{l+s+s2}{\PYGZdq{}Iglesia San Pedro Claver, (Calle 14 \PYGZsh{} 3 \PYGZhy{} 21)\PYGZdq{}}

\PYG{n+nv}{d30}\PYG{+w}{ }\PYG{o}{=}\PYG{+w}{ }\PYG{l+s+s2}{\PYGZdq{}Iglesia Cristo Rey, (Cra 4a con calle 5)\PYGZdq{}}
\end{sphinxVerbatim}

\item {} 
\sphinxAtStartPar
Cada instrucción(operación) debe describir como ir de una \sphinxstylestrong{dirección A} a una \sphinxstylestrong{dirección B}. Por ejemplo:

\end{itemize}


\begin{savenotes}\sphinxattablestart
\sphinxthistablewithglobalstyle
\centering
\begin{tabulary}{\linewidth}[t]{T}
\sphinxtoprule
\sphinxtableatstartofbodyhook

\begin{sphinxfigure-in-table}
\centering
\capstart
\noindent\sphinxincludegraphics[width=700\sphinxpxdimen]{{image052}.png}
\sphinxfigcaption{\sphinxstylestrong{Diagrama de flujo:} Orientaciones al extranjero}\label{\detokenize{cursos/ofimatica11/periodo01/G02:id2}}\end{sphinxfigure-in-table}\relax
\\
\sphinxbottomrule
\end{tabulary}
\sphinxtableafterendhook\par
\sphinxattableend\end{savenotes}
\end{quote}
\begin{itemize}
\item {} 
\sphinxAtStartPar
Usar una \sphinxstylestrong{página de cuaderno} por cada diagrama.

\item {} 
\sphinxAtStartPar
Para identificar las direcciones, se recomienda usar \sphinxhref{https://www.google.com/maps/@5.9736064,-74.5865216,7242m}{Google Maps}.

\item {} 
\sphinxAtStartPar
No olvide que el extranjero prefiere caminar.

\item {} 
\sphinxAtStartPar
Antes de orientar al extranjero para ir a una nueva dirección, siempre verificar la ubicación actual para evitar que el extranjero se \sphinxstylestrong{teletransporte por error}.

\end{itemize}

\end{enumerate}


\bigskip\hrule\bigskip



\paragraph{A03: Orientaciones seguras (+3,0)}
\label{\detokenize{cursos/ofimatica11/periodo01/G02:a03-orientaciones-seguras-3-0}}
\sphinxAtStartPar
\DUrole{sd-sphinx-override}{\DUrole{sd-badge}{\DUrole{sd-bg-danger}{\DUrole{sd-bg-text-danger}{Transferencia}}}}
\begin{enumerate}
\sphinxsetlistlabels{\arabic}{enumi}{enumii}{}{.}%
\item {} 
\sphinxAtStartPar
En su cuaderno, adapte el algoritmo de la actividad anterior, de tal forma que contenga 3 condicionales.
\begin{quote}

\sphinxAtStartPar
Los condicionales permitiran al extranjero tomar caminos alternos en caso de la ocurrencia de algún evento inesperado (vías cerradas, semáforos en rojo, perros bravos, negocios cerrados, etc.).
\end{quote}

\sphinxAtStartPar
Por ejemplo:


\begin{savenotes}\sphinxattablestart
\sphinxthistablewithglobalstyle
\centering
\begin{tabulary}{\linewidth}[t]{T}
\sphinxtoprule
\sphinxtableatstartofbodyhook

\begin{sphinxfigure-in-table}
\centering
\capstart
\noindent\sphinxincludegraphics[width=700\sphinxpxdimen]{{image062}.png}
\sphinxfigcaption{\sphinxstylestrong{Diagrama de flujo:} Instrucciones seguras}\label{\detokenize{cursos/ofimatica11/periodo01/G02:id3}}\end{sphinxfigure-in-table}\relax
\\
\sphinxbottomrule
\end{tabulary}
\sphinxtableafterendhook\par
\sphinxattableend\end{savenotes}
\begin{itemize}
\item {} 
\sphinxAtStartPar
Usar una \sphinxstylestrong{página de cuaderno} por cada diagrama.

\end{itemize}

\end{enumerate}

\sphinxstepscope


\subsubsection{Guía 03: Desarrollo de algoritmos}
\label{\detokenize{cursos/ofimatica11/periodo01/G03:guia-03-desarrollo-de-algoritmos}}\label{\detokenize{cursos/ofimatica11/periodo01/G03::doc}}

\begin{savenotes}\sphinxattablestart
\sphinxthistablewithglobalstyle
\centering
\begin{tabular}[t]{\X{35}{100}\X{65}{100}}
\sphinxtoprule
\sphinxtableatstartofbodyhook

\begin{savenotes}\sphinxattablestart
\sphinxthistablewithglobalstyle
\centering
\begin{tabulary}{\linewidth}[t]{T}
\sphinxtoprule
\sphinxtableatstartofbodyhook
\sphinxAtStartPar
\sphinxincludegraphics{{escudo}.png}
\\
\sphinxbottomrule
\end{tabulary}
\sphinxtableafterendhook\par
\sphinxattableend\end{savenotes}
&

\begin{savenotes}\sphinxattablestart
\sphinxthistablewithglobalstyle
\centering
\begin{tabulary}{\linewidth}[t]{T}
\sphinxtoprule
\sphinxtableatstartofbodyhook
\sphinxAtStartPar
\sphinxstylestrong{INSTITUCIÓN EDUCATIVA}: John F. Kennedy
\\
\sphinxhline
\sphinxAtStartPar
\sphinxstylestrong{DOCENTE}: Julian Camilo Fonseca Romero
\\
\sphinxhline
\sphinxAtStartPar
\sphinxstylestrong{ASIGNATURA}: Ofimática
\\
\sphinxhline
\sphinxAtStartPar
\sphinxstylestrong{GRADO}: ONCE
\\
\sphinxhline
\sphinxAtStartPar
\sphinxstylestrong{TIEMPO}: 10 semanas (4 horas / semana)
\\
\sphinxhline
\sphinxAtStartPar
\sphinxstylestrong{TIEMPO DE TRABAJO}: 20/40 horas
\\
\sphinxbottomrule
\end{tabulary}
\sphinxtableafterendhook\par
\sphinxattableend\end{savenotes}
\\
\sphinxbottomrule
\end{tabular}
\sphinxtableafterendhook\par
\sphinxattableend\end{savenotes}

\begin{sphinxuseclass}{sd-tab-set}
\begin{sphinxuseclass}{sd-tab-item}\subsubsection*{Objetivo}

\begin{sphinxuseclass}{sd-tab-content}\begin{itemize}
\item {} 
\sphinxAtStartPar
\sphinxstylestrong{Temática}: Desarrollo de algoritmos.

\item {} 
\sphinxAtStartPar
\sphinxstylestrong{Objetivo de aprendizaje}: Desarrollar la capacidad de abstracción y lógica algorítmica mediante la representación gráfica de procesos, permitiéndole modelar soluciones secuenciales y ordenadas a problemas de la vida cotidiana.

\end{itemize}

\end{sphinxuseclass}
\end{sphinxuseclass}
\begin{sphinxuseclass}{sd-tab-item}\subsubsection*{Orientaciones curriculares}

\begin{sphinxuseclass}{sd-tab-content}\begin{itemize}
\item {} 
\sphinxAtStartPar
\sphinxstylestrong{Componente}: Solución de problemas con T\&I.

\item {} 
\sphinxAtStartPar
\sphinxstylestrong{Competencia}: Propongo innovaciones tecnológicas e informáticas para la solución de problemas dando cumplimiento a restricciones, condiciones y especificaciones técnicas y contextuales.

\item {} 
\sphinxAtStartPar
\sphinxstylestrong{Evidencia de aprendizaje}: Propongo, analizo y comparo diferentes soluciones a un mismo problema de la tecnología o la informática, explicando su origen, ventajas y dificultades.

\end{itemize}

\end{sphinxuseclass}
\end{sphinxuseclass}
\begin{sphinxuseclass}{sd-tab-item}\subsubsection*{Criterios de evaluación}

\begin{sphinxuseclass}{sd-tab-content}\begin{itemize}
\item {} 
\sphinxAtStartPar
Actividades desarrolladas en su totalidad.

\item {} 
\sphinxAtStartPar
Participación en clase.

\item {} 
\sphinxAtStartPar
Explicación clara de conceptos.

\item {} 
\sphinxAtStartPar
Argumentación de ideas.

\item {} 
\sphinxAtStartPar
Cumplimiento de las reglas del aula.

\end{itemize}

\end{sphinxuseclass}
\end{sphinxuseclass}
\begin{sphinxuseclass}{sd-tab-item}\subsubsection*{Recursos (1)}

\begin{sphinxuseclass}{sd-tab-content}\begin{itemize}
\item {} 
\sphinxAtStartPar
{\hyperref[\detokenize{cursos/ofimatica11/recursos/T03::doc}]{\sphinxcrossref{\DUrole{std}{\DUrole{std-doc}{TEXTO: Algoritmos}}}}}

\end{itemize}

\end{sphinxuseclass}
\end{sphinxuseclass}
\end{sphinxuseclass}

\paragraph{Mensaje del autor}
\label{\detokenize{cursos/ofimatica11/periodo01/G03:mensaje-del-autor}}\begin{quote}
\begin{quote}

\sphinxAtStartPar
“Entender cómo funciona el mundo actual requiere aprender a pensar de forma organizada. Esta guía invita a usar el desarrollo de algoritmos como una herramienta para abstraer y modelar soluciones. Al visualizar sus ideas, usted podrá comparar distintas rutas lógicas para un mismo problema, evaluando sus beneficios y desafíos.”
\end{quote}

\begin{flushright}
---Camilo Fonseca
\end{flushright}
\end{quote}


\bigskip\hrule\bigskip



\paragraph{A00: Introducción (+0,2)}
\label{\detokenize{cursos/ofimatica11/periodo01/G03:a00-introduccion-0-2}}
\sphinxAtStartPar
\DUrole{sd-sphinx-override}{\DUrole{sd-badge}{\DUrole{sd-bg-light}{\DUrole{sd-bg-text-light}{Exploración}}}}
\begin{enumerate}
\sphinxsetlistlabels{\arabic}{enumi}{enumii}{}{.}%
\item {} 
\sphinxAtStartPar
TRANSCRIBA en su cuaderno el objetivo, las orientaciones curriculares y los criterios de evaluación de la presente guía.

\item {} 
\sphinxAtStartPar
LEA el \sphinxstylestrong{mensaje del autor} y en su cuaderno responda:
\begin{itemize}
\item {} 
\sphinxAtStartPar
¿Qué es abstracción?

\item {} 
\sphinxAtStartPar
¿Qué es modelar una solución?

\end{itemize}

\end{enumerate}


\bigskip\hrule\bigskip



\paragraph{A01: ¿Qué son los algoritmos? (+0,2)}
\label{\detokenize{cursos/ofimatica11/periodo01/G03:a01-que-son-los-algoritmos-0-2}}
\sphinxAtStartPar
\DUrole{sd-sphinx-override}{\DUrole{sd-badge}{\DUrole{sd-bg-info}{\DUrole{sd-bg-text-info}{Estructuración}}}}
\begin{enumerate}
\sphinxsetlistlabels{\arabic}{enumi}{enumii}{}{.}%
\item {} 
\sphinxAtStartPar
Escriba las siguientes frases en su cuaderno e indique si es \sphinxstylestrong{falso} o \sphinxstylestrong{verdadero} e indique el por qué.
\begin{itemize}
\item {} 
\sphinxAtStartPar
Un algoritmo está compuesto por elementos informáticos que componen una secuencia organizacional únicamente \_\_\_\_\_\_\_.

\item {} 
\sphinxAtStartPar
Un algoritmo son un conjunto de instrucciones que le dicen a una computadora o un ser humano como realizar una tarea específica, dichas instrucciones pueden estar escritas en código o lenguaje natural. \_\_\_\_\_\_\_.

\item {} 
\sphinxAtStartPar
Un algoritmo se define como cuadros de funcionamiento infinitos que permiten y tienen como propósito reducir la información. \_\_\_\_\_\_\_.

\end{itemize}

\end{enumerate}


\bigskip\hrule\bigskip



\paragraph{A02: Tipos del algoritmos (+0,2)}
\label{\detokenize{cursos/ofimatica11/periodo01/G03:a02-tipos-del-algoritmos-0-2}}
\sphinxAtStartPar
\DUrole{sd-sphinx-override}{\DUrole{sd-badge}{\DUrole{sd-bg-info}{\DUrole{sd-bg-text-info}{Estructuración}}}}
\begin{enumerate}
\sphinxsetlistlabels{\arabic}{enumi}{enumii}{}{.}%
\item {} 
\sphinxAtStartPar
En su cuaderno relacione con líneas los siguientes cuadros según corresponda su significado:

\end{enumerate}


\begin{savenotes}\sphinxattablestart
\sphinxthistablewithglobalstyle
\centering
\begin{tabulary}{\linewidth}[t]{T}
\sphinxtoprule
\sphinxtableatstartofbodyhook

\begin{sphinxfigure-in-table}
\centering
\capstart
\noindent\sphinxincludegraphics[width=500\sphinxpxdimen]{{image072}.png}
\sphinxfigcaption{Cuadro comparativo (Tipos de algoritmos)}\label{\detokenize{cursos/ofimatica11/periodo01/G03:id1}}\end{sphinxfigure-in-table}\relax
\\
\sphinxbottomrule
\end{tabulary}
\sphinxtableafterendhook\par
\sphinxattableend\end{savenotes}


\bigskip\hrule\bigskip



\paragraph{A03: Clasificación de algoritmos (+0,2)}
\label{\detokenize{cursos/ofimatica11/periodo01/G03:a03-clasificacion-de-algoritmos-0-2}}
\sphinxAtStartPar
\DUrole{sd-sphinx-override}{\DUrole{sd-badge}{\DUrole{sd-bg-info}{\DUrole{sd-bg-text-info}{Estructuración}}}}
\begin{enumerate}
\sphinxsetlistlabels{\arabic}{enumi}{enumii}{}{.}%
\item {} 
\sphinxAtStartPar
Escriba la siguiente tabla en su cuaderno y clasifique los algoritmos según corresponda:

\end{enumerate}


\begin{savenotes}\sphinxattablestart
\sphinxthistablewithglobalstyle
\raggedright
\begin{tabulary}{\linewidth}[t]{TTT}
\sphinxtoprule
\sphinxstyletheadfamily 
\sphinxAtStartPar
EJEMPLO
&\sphinxstyletheadfamily 
\sphinxAtStartPar
Clasificación
&\sphinxstyletheadfamily 
\sphinxAtStartPar
¿Por qué?
\\
\sphinxmidrule
\sphinxtableatstartofbodyhook
\sphinxAtStartPar
Instrucciones para llegar a un lugar
&
\sphinxAtStartPar

&
\sphinxAtStartPar

\\
\sphinxhline
\sphinxAtStartPar
Instrucciones con medición de tiempo para llegar a un lugar
&
\sphinxAtStartPar

&
\sphinxAtStartPar

\\
\sphinxhline
\sphinxAtStartPar
A=10,B=12C=A+B
&
\sphinxAtStartPar

&
\sphinxAtStartPar

\\
\sphinxhline
\sphinxAtStartPar
Microsoft Word
&
\sphinxAtStartPar

&
\sphinxAtStartPar

\\
\sphinxhline
\sphinxAtStartPar
LINUX
&
\sphinxAtStartPar

&
\sphinxAtStartPar

\\
\sphinxhline
\sphinxAtStartPar
\sphinxcode{\sphinxupquote{print("Digite su nombre")}}
&
\sphinxAtStartPar

&
\sphinxAtStartPar

\\
\sphinxhline
\sphinxAtStartPar
PASO 01: Abra la bandeja de entradaPASO 02: Inserte hojas tamaño cartaPASO 03: Cierre la bandeja de entrada
&
\sphinxAtStartPar

&
\sphinxAtStartPar

\\
\sphinxhline
\sphinxAtStartPar
BIOS
&
\sphinxAtStartPar

&
\sphinxAtStartPar

\\
\sphinxhline
\sphinxAtStartPar
\sphinxstylestrong{CANVA}
&
\sphinxAtStartPar

&
\sphinxAtStartPar

\\
\sphinxhline
\sphinxAtStartPar
LibreOffice Writer
&
\sphinxAtStartPar

&
\sphinxAtStartPar

\\
\sphinxhline
\sphinxAtStartPar
Control de vuelo con API de GOOGLE
&
\sphinxAtStartPar

&
\sphinxAtStartPar

\\
\sphinxhline
\sphinxAtStartPar
Sensor Temperatura con ARDUINO
&
\sphinxAtStartPar

&
\sphinxAtStartPar

\\
\sphinxhline
\sphinxAtStartPar
Instrucciones corte cabello
&
\sphinxAtStartPar

&
\sphinxAtStartPar

\\
\sphinxbottomrule
\end{tabulary}
\sphinxtableafterendhook\par
\sphinxattableend\end{savenotes}


\bigskip\hrule\bigskip



\paragraph{A04: Metodología para desarrollar un algoritmo (+0,2)}
\label{\detokenize{cursos/ofimatica11/periodo01/G03:a04-metodologia-para-desarrollar-un-algoritmo-0-2}}
\sphinxAtStartPar
\DUrole{sd-sphinx-override}{\DUrole{sd-badge}{\DUrole{sd-bg-info}{\DUrole{sd-bg-text-info}{Estructuración}}}}
\begin{enumerate}
\sphinxsetlistlabels{\arabic}{enumi}{enumii}{}{.}%
\item {} 
\sphinxAtStartPar
En su cuaderno escriba las 3 fases para desarrollar un algoritmo.

\item {} 
\sphinxAtStartPar
Responda la siguiente pregunta en su cuaderno:
\begin{itemize}
\item {} 
\sphinxAtStartPar
¿Que diferencia hay entre desarrollar un algoritmo, diseñar un algoritmo y construir un algoritmo?

\end{itemize}

\item {} 
\sphinxAtStartPar
Consulte en INTERNET que hace un arquitecto de software y luego en su cuaderno responda:
\begin{itemize}
\item {} 
\sphinxAtStartPar
¿Cual es la diferencia entre un arquitecto de software y un programador?

\end{itemize}

\end{enumerate}


\bigskip\hrule\bigskip



\paragraph{A05: Algoritmos (+4,0)}
\label{\detokenize{cursos/ofimatica11/periodo01/G03:a05-algoritmos-4-0}}
\sphinxAtStartPar
\DUrole{sd-sphinx-override}{\DUrole{sd-badge}{\DUrole{sd-bg-danger}{\DUrole{sd-bg-text-danger}{Transferencia}}}}
\begin{enumerate}
\sphinxsetlistlabels{\arabic}{enumi}{enumii}{}{.}%
\item {} 
\sphinxAtStartPar
En su cuaderno desarrolle cada uno de los siguientes ejercicios, siguiendo la metodología para desarrollar un algoritmo.

\end{enumerate}
\subsubsection*{Adquirir un revista física de arte}

\sphinxAtStartPar
Desarrolle un algoritmo para que una persona adquiera una revista física de arte. Tenga en cuenta lo siguiente:
\begin{quote}

\sphinxAtStartPar
La persona inicialmente busca en INTERNET una revista física de arte.

\sphinxAtStartPar
Una vez encontrada la revista, revisar cuanto cuesta y ver si tiene el dinero suficiente para comprarla. Si no tiene el dinero suficiente volver a buscar la revista en INTERNET cuantas veces sea necesario. Si tiene el dinero suficiente alistar el dinero.

\sphinxAtStartPar
Luego la persona debe revisar en GOOGLE MAPS para localizar la dirección de la librería. Activar la función de GOOGLE MAPS y dirigirse caminando hacia la librería. Intentar entrar a la librería. Si la librería no está abierta: volver a buscar la revista en INTERNET e iniciar de nuevo. Si la librería está abierta, entrar a la librería.

\sphinxAtStartPar
Revisar si la revista está en STOCK. Si la revista no está en STOCK, volver a buscar la revista en INTERNET y empezar de nuevo. SI la revista si está en STOCK, comprar la revista física de arte.
\end{quote}
\subsubsection*{Entrar a una casa con llave}

\sphinxAtStartPar
Desarrolle un algoritmo para que una persona entre a una casa cuya cerradura está asegurada con llave. Tenga en cuenta lo siguiente:
\begin{quote}

\sphinxAtStartPar
La persona inicialmente buscará la llave en su bolsillo. La persona intentará abrir la cerradura con la llave de su bolsillo.

\sphinxAtStartPar
Si la cerradura no abre ahora buscará otra llave en un tapete exterior y volverá a intentar abrir la cerradura con dicha llave.

\sphinxAtStartPar
Si la cerradura no abre ahora buscará otra llave en la guantera de un carro y volverá a intentar abrir la cerradura con dicha llave.

\sphinxAtStartPar
Si la cerradura no abre, llamará a un cerrajero para que abra la puerta. Si el cerrajero no contesta, entonces seguir intentando hasta que conteste.

\sphinxAtStartPar
Si la cerradura si abre, intentar abrir la puerta. Si la puerta abre entrar a la casa y confirmar que se logró entrar. Si la cerradura no abre, confirmar que no se logró entrar.
\end{quote}
\subsubsection*{Saludar y despedirse de una persona}

\sphinxAtStartPar
Desarrolle un algoritmo para que el usuario salude y se despida de una persona. Tenga en cuenta lo siguiente:
\begin{quote}

\sphinxAtStartPar
La persona debe verificar inicialmente si ya saludo a la otra persona

\sphinxAtStartPar
Si la persona no ha saludado, saludar a la persona. Esperar a que la persona salude. Si la persona no saluda, seguir intentándolo hasta que salude.

\sphinxAtStartPar
Una vez que la persona salude, preguntar como esta. Si la persona responde NO esta bien, preguntar que le pasa, escuchar a la persona y decirle que el tiempo lo arregla todo. Volver a preguntar si la persona esta bien (cuantas veces sea necesario) hasta que la persona responda que esta bien.  Si la persona responder que SI esta bien entonces manifestar tranquilidad.

\sphinxAtStartPar
Volver a verificar si la persona ya saludo.

\sphinxAtStartPar
Si la persona ya saludo, entonces determinar un tiempo para poder irse.

\sphinxAtStartPar
El usuario debe establecer el tiempo para irse.

\sphinxAtStartPar
Verificar que el tiempo se cumplió. Mientras el tiempo no se cumpla, entonces hablar con la persona de cualquier tema mientras se cumple el tiempo. Ir disminuyendo el tiempo.

\sphinxAtStartPar
Una vez el tiempo esté cumplido, despedirse. Esperar a que la persona responda la despedida. Si la persona no contesta seguir despidiéndose hasta que la persona conteste.

\sphinxAtStartPar
Una vez se despida de la persona generar una mensaje de IRSE.
\end{quote}
\subsubsection*{Calmar a una persona}

\sphinxAtStartPar
Desarrolle un algoritmo para calmar una persona. Tenga en cuenta lo siguiente:
\begin{quote}

\sphinxAtStartPar
Identificar el estado de la persona.

\sphinxAtStartPar
Si la persona esta bien, entonces irse del lugar. Si la persona esta en mal estado realizar lo siguiente:

\sphinxAtStartPar
Identificar la edad de la persona.

\sphinxAtStartPar
Si la persona es un bebe entonces intentar calmarla con un tetero. Si la persona es un niño entonces calmarla dandole un dulce. Si la persona es un adolescente entonces intentar calmarla dandole una gaseosa. Si la persona es un adulto entonces intentar calmarla dandole una cerveza. Si la persona es un anciano entonces intentar calmarla dandole un café.

\sphinxAtStartPar
Luego volver a identificar el estado de la persona y si todavía no se ha calmado, seguir intentando hasta que se calme.
\end{quote}

\sphinxstepscope


\subsubsection{Plan de refuerzo 01}
\label{\detokenize{cursos/ofimatica11/periodo01/R01:plan-de-refuerzo-01}}\label{\detokenize{cursos/ofimatica11/periodo01/R01::doc}}
\begin{sphinxadmonition}{note}{EVALUACIÓN: Plan de refuerzo (10 minutos)}

\sphinxAtStartPar
Estudiar los temas a continuación. Se realiza examen de 5 preguntas. Si el estudiante no asiste en la fecha programada, este será evaluado la próxima vez que asista a clase.
\end{sphinxadmonition}


\bigskip\hrule\bigskip



\paragraph{Instrucción}
\label{\detokenize{cursos/ofimatica11/periodo01/R01:instruccion}}
\sphinxAtStartPar
Una instrucción es una orden o indicación que se da para realizar una acción específica.


\bigskip\hrule\bigskip



\paragraph{Hardware}
\label{\detokenize{cursos/ofimatica11/periodo01/R01:hardware}}
\sphinxAtStartPar
El hardware es el conjunto de componentes de un sistema informático o computacional que se puede ver o tocar. Por ejemplo:
\begin{itemize}
\item {} 
\sphinxAtStartPar
\sphinxstylestrong{Dispositivos de entrada:} Teclados, mouse, micrófonos.

\item {} 
\sphinxAtStartPar
\sphinxstylestrong{Dispositivos de salida:} Monitores, impresoras, bafles.

\item {} 
\sphinxAtStartPar
\sphinxstylestrong{Procesamiento:} CPU.

\item {} 
\sphinxAtStartPar
\sphinxstylestrong{Memoria:} memoria RAM, memoria CACHE.

\item {} 
\sphinxAtStartPar
\sphinxstylestrong{Almacenamiento:} Discos duros, memorias USB.

\end{itemize}


\bigskip\hrule\bigskip



\paragraph{Software}
\label{\detokenize{cursos/ofimatica11/periodo01/R01:software}}
\sphinxAtStartPar
El software es el conjunto de programas, bases de datos, diferentes tipos de archivos y documentación que le dicen al hardware qué hacer. Es la parte lógica e intangible de un sistema informático. Ejemplos de software:
\begin{itemize}
\item {} 
\sphinxAtStartPar
Sistemas operativos

\item {} 
\sphinxAtStartPar
Aplicaciones

\item {} 
\sphinxAtStartPar
Drivers

\end{itemize}


\bigskip\hrule\bigskip



\paragraph{Información y datos}
\label{\detokenize{cursos/ofimatica11/periodo01/R01:informacion-y-datos}}
\sphinxAtStartPar
Los datos son valores en bruto, hechos o símbolos que por sí solos no tienen significado. Por ejemplo:
\begin{itemize}
\item {} 
\sphinxAtStartPar
Números sueltos.

\item {} 
\sphinxAtStartPar
Caracteres sin contexto.

\item {} 
\sphinxAtStartPar
Valores sin procesar.

\end{itemize}

\sphinxAtStartPar
La información es el resultado de procesar, organizar e interpretar los datos de manera que tengan sentido y sean útiles. La información:
\begin{itemize}
\item {} 
\sphinxAtStartPar
Tiene significado y contexto.

\item {} 
\sphinxAtStartPar
Es comprensible para el usuario.

\item {} 
\sphinxAtStartPar
Ayuda en la toma de decisiones.

\end{itemize}


\bigskip\hrule\bigskip



\paragraph{Software Vs Programa}
\label{\detokenize{cursos/ofimatica11/periodo01/R01:software-vs-programa}}

\subparagraph{Software}
\label{\detokenize{cursos/ofimatica11/periodo01/R01:id1}}
\sphinxAtStartPar
Un \sphinxstylestrong{software} es el conjunto de \sphinxstylestrong{programas}, bases de datos, archivos y documentación que permiten realizar diferentes tareas en un sistema computacional.


\subparagraph{Programa}
\label{\detokenize{cursos/ofimatica11/periodo01/R01:programa}}
\sphinxAtStartPar
Un \sphinxstylestrong{programa} es un \sphinxstylestrong{conjunto de instrucciones} que le dicen a una computadora COMO HACER una tarea específica. Las instrucciones están escritas en código y son diseñadas por programadores con conocimiento en un lenguaje de programación.

\sphinxstepscope


\section{Recursos}
\label{\detokenize{cursos/ofimatica11/recursos/index:recursos}}\label{\detokenize{cursos/ofimatica11/recursos/index::doc}}
\sphinxAtStartPar
En esta sección encuentra los recursos necesarios para el desarrollo del curso.


\bigskip\hrule\bigskip



\subsection{Contenidos}
\label{\detokenize{cursos/ofimatica11/recursos/index:contenidos}}
\sphinxstepscope


\subsubsection{T01: Conocimientos previos}
\label{\detokenize{cursos/ofimatica11/recursos/T01:t01-conocimientos-previos}}\label{\detokenize{cursos/ofimatica11/recursos/T01::doc}}
\sphinxAtStartPar
Esta sección muestra los conceptos de datos e información, jerarquía de datos, hardware y software. Estos conceptos son introductorios a la programación.


\bigskip\hrule\bigskip



\paragraph{Información y Datos}
\label{\detokenize{cursos/ofimatica11/recursos/T01:informacion-y-datos}}
\sphinxAtStartPar
%
\begin{footnote}[1]\sphinxAtStartFootnote
Trejos Buriticá, O. I. (2017). Lógica de programación. Ediciones de la U.
%
\end{footnote} En informática, es importante diferenciar entre los conceptos de datos e información:


\subparagraph{Datos}
\label{\detokenize{cursos/ofimatica11/recursos/T01:datos}}
\sphinxAtStartPar
Los datos son valores en bruto, hechos o símbolos que por sí solos no tienen significado. Por ejemplo:
\begin{itemize}
\item {} 
\sphinxAtStartPar
Números sueltos.

\item {} 
\sphinxAtStartPar
Caracteres sin contexto.

\item {} 
\sphinxAtStartPar
Valores sin procesar.

\end{itemize}


\subparagraph{Información}
\label{\detokenize{cursos/ofimatica11/recursos/T01:informacion}}
\sphinxAtStartPar
La información es el resultado de procesar, organizar e interpretar los datos de manera que tengan sentido y sean útiles. La información:
\begin{itemize}
\item {} 
\sphinxAtStartPar
Tiene significado y contexto.

\item {} 
\sphinxAtStartPar
Es comprensible para el usuario.

\item {} 
\sphinxAtStartPar
Ayuda en la toma de decisiones.

\end{itemize}

\sphinxAtStartPar
Por ejemplo:


\begin{savenotes}\sphinxattablestart
\sphinxthistablewithglobalstyle
\centering
\begin{tabulary}{\linewidth}[t]{TT}
\sphinxtoprule
\sphinxstyletheadfamily 
\sphinxAtStartPar
Datos
&\sphinxstyletheadfamily 
\sphinxAtStartPar
Información
\\
\sphinxmidrule
\sphinxtableatstartofbodyhook
\sphinxAtStartPar
37,5
&
\sphinxAtStartPar
La temperatura del paciente es 37.5°C, lo cual indica fiebre
\\
\sphinxhline
\sphinxAtStartPar
4,5
&
\sphinxAtStartPar
La nota final del estudiante es 4.5
\\
\sphinxbottomrule
\end{tabulary}
\sphinxtableafterendhook\par
\sphinxattableend\end{savenotes}

\sphinxAtStartPar
La transformación de datos en información es una parte fundamental del procesamiento computacional, ya que permite convertir valores sin sentido en conocimiento útil y aplicable.


\bigskip\hrule\bigskip



\paragraph{Valores cuantitativos y cualitativos}
\label{\detokenize{cursos/ofimatica11/recursos/T01:valores-cuantitativos-y-cualitativos}}
\sphinxAtStartPar
\sphinxfootnotemark[1] Los valores cualitativos y cuantitativos son dos tipos diferentes de información:
\begin{itemize}
\item {} 
\sphinxAtStartPar
\sphinxstylestrong{Valores cualitativos:} Describen cualidades o características que no se pueden medir numéricamente. Por ejemplo:
\begin{itemize}
\item {} 
\sphinxAtStartPar
Color de ojos

\item {} 
\sphinxAtStartPar
Estado de ánimo

\item {} 
\sphinxAtStartPar
Nivel de satisfacción (excelente, bueno, regular)

\end{itemize}

\item {} 
\sphinxAtStartPar
\sphinxstylestrong{Valores cuantitativos:} Son aquellos que se pueden medir y expresar en números. Por ejemplo:
\begin{itemize}
\item {} 
\sphinxAtStartPar
Edad: 12 años

\item {} 
\sphinxAtStartPar
Peso: 99 Kg

\item {} 
\sphinxAtStartPar
Temperatura: 40°C

\item {} 
\sphinxAtStartPar
Calificaciones numéricas: 3,0

\end{itemize}

\end{itemize}

\sphinxAtStartPar
La principal diferencia es que los valores cualitativos describen atributos o cualidades, mientras que los cuantitativos expresan cantidades medibles.


\bigskip\hrule\bigskip



\paragraph{Jerarquía de datos}
\label{\detokenize{cursos/ofimatica11/recursos/T01:jerarquia-de-datos}}
\sphinxAtStartPar
%
\begin{footnote}[2]\sphinxAtStartFootnote
Deitel, P. J., \& Deitel, H. (2014). C++: cómo programar (A. V. Romero Elizondo, Trans.). Pearson Educación.
%
\end{footnote} La jerarquía de datos es la organización lógica de los datos en niveles sucesivos de complejidad, desde la unidad más pequeña hasta la más completa. En informática, se utiliza para entender cómo se estructura la información dentro de un sistema o base de datos.

\sphinxAtStartPar
Los datos en la computadora se clasifican de MENOR a MAYOR en el siguiente ranking:
\begin{enumerate}
\sphinxsetlistlabels{\arabic}{enumi}{enumii}{}{.}%
\item {} 
\sphinxAtStartPar
\sphinxstylestrong{bit:} unidad lógica puede ser \sphinxcode{\sphinxupquote{1}} o \sphinxcode{\sphinxupquote{0}}.

\item {} 
\sphinxAtStartPar
\sphinxstylestrong{BYTE:} Son 8 bits. Ejemplo: \sphinxcode{\sphinxupquote{10010010}}

\item {} 
\sphinxAtStartPar
\sphinxstylestrong{Caracteres:} Representaciones de Bytes. Ejemplo: \sphinxstylestrong{\sphinxcode{\sphinxupquote{A}}}: \sphinxcode{\sphinxupquote{10001001}}, \sphinxstylestrong{\sphinxcode{\sphinxupquote{B}}:} \sphinxcode{\sphinxupquote{10001001}}

\item {} 
\sphinxAtStartPar
\sphinxstylestrong{Dato:} Se componen de representaciones de Caracteres. Ejemplo: \sphinxcode{\sphinxupquote{'CAMILO'}}, \sphinxcode{\sphinxupquote{'ROJO'}}, \sphinxcode{\sphinxupquote{'FRIO'}}, \sphinxcode{\sphinxupquote{23}} , \sphinxcode{\sphinxupquote{0.0999}}, \sphinxguilabel{True},

\item {} 
\sphinxAtStartPar
\sphinxstylestrong{Campo:} Nombre que se le coloca a un dato para identificarlo. Ejemplo: \sphinxcode{\sphinxupquote{nombre}}, \sphinxcode{\sphinxupquote{color}}, \sphinxcode{\sphinxupquote{clima}}, \sphinxcode{\sphinxupquote{edad}}, \sphinxcode{\sphinxupquote{temperatura}}, \sphinxcode{\sphinxupquote{estado}}

\item {} 
\sphinxAtStartPar
\sphinxstylestrong{Registro:} Conjunto de campos que conforman un solo objeto. Ejemplo: El observador de un estudiante, la hoja técnica de una moto.

\item {} 
\sphinxAtStartPar
\sphinxstylestrong{Tablas:} Conjunto de registros que tienen la misma estructura y que pertenecen a una organización. Ejemplo:
\begin{quote}

\sphinxAtStartPar
“La \sphinxstylestrong{tabla} de los \sphinxstyleemphasis{estudiantes} de un colegio”.
\end{quote}

\item {} 
\sphinxAtStartPar
\sphinxstylestrong{Bases de datos:} Conjunto de tablas relacionadas técnicamente y organizadas. Ejemplo:
\begin{quote}

\sphinxAtStartPar
“La \sphinxstylestrong{base de datos} donde están relacionadas las tablas de los estudiantes junto con las tablas de las notas”.
\end{quote}

\end{enumerate}


\bigskip\hrule\bigskip



\paragraph{Hardware}
\label{\detokenize{cursos/ofimatica11/recursos/T01:hardware}}
\sphinxAtStartPar
\sphinxfootnotemark[2] El hardware es el conjunto de componentes de un sistema informático o computacional que se puede ver o tocar. Por ejemplo:
\begin{itemize}
\item {} 
\sphinxAtStartPar
\sphinxstylestrong{Dispositivos de entrada:} Teclados, mouse, micrófonos.

\item {} 
\sphinxAtStartPar
\sphinxstylestrong{Dispositivos de salida:} Monitores, impresoras, bafles.

\item {} 
\sphinxAtStartPar
\sphinxstylestrong{Procesamiento:} CPU.

\item {} 
\sphinxAtStartPar
\sphinxstylestrong{Memoria:} memoria RAM, memoria CACHE.

\item {} 
\sphinxAtStartPar
\sphinxstylestrong{Almacenamiento:} Discos duros, memorias USB.

\end{itemize}


\subparagraph{La memoria RAM (acceso rápido a los datos)}
\label{\detokenize{cursos/ofimatica11/recursos/T01:la-memoria-ram-acceso-rapido-a-los-datos}}
\sphinxAtStartPar
La \sphinxstylestrong{memoria RAM} (\sphinxstyleemphasis{Random Access Memory o Memoria de Acceso Aleatorio}) es un componente fundamental del hardware que:
\begin{itemize}
\item {} 
\sphinxAtStartPar
Almacena temporalmente los datos y programas que el computador está usando activamente.

\item {} 
\sphinxAtStartPar
Permite acceder rápidamente a la información mientras el equipo está encendido.

\item {} 
\sphinxAtStartPar
Se borra cuando se apaga el computador (memoria volátil).

\end{itemize}

\sphinxAtStartPar
Características importantes de la RAM:
\begin{itemize}
\item {} 
\sphinxAtStartPar
A mayor cantidad de RAM, más programas pueden ejecutarse simultáneamente.

\item {} 
\sphinxAtStartPar
Su velocidad afecta directamente el rendimiento del sistema.

\item {} 
\sphinxAtStartPar
Es diferente al almacenamiento permanente (disco duro) porque es temporal.

\end{itemize}


\subparagraph{La CPU (procesa los datos)}
\label{\detokenize{cursos/ofimatica11/recursos/T01:la-cpu-procesa-los-datos}}
\sphinxAtStartPar
La \sphinxstylestrong{CPU (Unidad Central de Procesamiento)} es el “cerebro” del computador. Es el componente principal que se encarga de:
\begin{itemize}
\item {} 
\sphinxAtStartPar
Procesar todas las instrucciones y cálculos del sistema

\item {} 
\sphinxAtStartPar
Coordinar el funcionamiento de los demás componentes

\item {} 
\sphinxAtStartPar
Ejecutar los programas y aplicaciones

\end{itemize}

\sphinxAtStartPar
La \sphinxstylestrong{CPU} está compuesta principalmente por:
\begin{itemize}
\item {} 
\sphinxAtStartPar
\sphinxstylestrong{Unidad de Control:} Coordina y dirige todas las operaciones

\item {} 
\sphinxAtStartPar
\sphinxstylestrong{Unidad Aritmético\sphinxhyphen{}Lógica (ALU):} Realiza cálculos matemáticos y operaciones lógicas

\item {} 
\sphinxAtStartPar
\sphinxstylestrong{Registros:} Pequeñas unidades de memoria de alta velocidad

\end{itemize}

\sphinxAtStartPar
La velocidad de la \sphinxstylestrong{CPU} se mide en Hertz (Hz) y determina qué tan rápido puede procesar las instrucciones. A mayor velocidad, mejor rendimiento del sistema.


\bigskip\hrule\bigskip



\paragraph{Software}
\label{\detokenize{cursos/ofimatica11/recursos/T01:software}}
\sphinxAtStartPar
\sphinxfootnotemark[2] Un \sphinxstylestrong{Software} es el conjunto de \sphinxstylestrong{programas}, bases de datos, archivos y documentación que permiten realizar diferentes tareas en un sistema computacional.

\sphinxAtStartPar
Existen dos tipos de software: Software de sistema y software de aplicación.
\begin{itemize}
\item {} 
\sphinxAtStartPar
\sphinxstylestrong{Software de sistema:} Ejecutan todas las funciones básicas de una computadora (por ejemplo encender la máquina, mostrar imagen en pantalla, monitorear temperatura, comunicarse con los diferentes periféricos, etc.). El \sphinxstylestrong{software de sistema} controla e intercambia datos con el hardware. Un ejemplo de software de sistema es el sistema operativo Windows. Inclusive los controladores o drivers también hacen parte del software de sistema.

\item {} 
\sphinxAtStartPar
\sphinxstylestrong{Software de aplicación:} Permite al usuario trabajar en una tarea específica (procesar texto, edición de video, edición de imágenes, navegación en INTERNET, desarrollo de software, etc.). Por lo general, el software de aplicación trabaja en función del software de sistema.

\end{itemize}

\sphinxAtStartPar
A diferencia del hardware que es físico y tangible, el software no se puede tocar \sphinxhyphen{} solo puede verse su funcionamiento en la pantalla.


\bigskip\hrule\bigskip



\paragraph{Software Vs Programa}
\label{\detokenize{cursos/ofimatica11/recursos/T01:software-vs-programa}}
\sphinxAtStartPar
Un programa es un \sphinxstylestrong{conjunto de instrucciones} que le dicen a una \sphinxstylestrong{computadora} COMO HACER una tarea específica. Las instrucciones están escritas en código y son diseñadas por programadores con conocimiento en un lenguaje de programación.

\sphinxAtStartPar
A diferencia de un \sphinxstylestrong{algoritmo}, las instrucciones de un programa \sphinxstylestrong{NO pueden ser escritas en lenguaje natural}.
\begin{quote}

\sphinxAtStartPar
“Aunque se usa el término \sphinxstylestrong{software} como sinónimo, la realidad es que un \sphinxstylestrong{programa} es solo una parte del \sphinxstylestrong{software}”.
\end{quote}


\bigskip\hrule\bigskip


\sphinxstepscope


\subsubsection{T02: Diagramas de flujo}
\label{\detokenize{cursos/ofimatica11/recursos/T02:t02-diagramas-de-flujo}}\label{\detokenize{cursos/ofimatica11/recursos/T02::doc}}
\sphinxAtStartPar
%
\begin{footnote}[1]\sphinxAtStartFootnote
Deitel, P. J., \& Deitel, H. (2014). C++: cómo programar (A. V. Romero Elizondo, Trans.). Pearson Educación.
%
\end{footnote} Un diagrama de flujo es un \sphinxstylestrong{esquema gráfico} donde se muestra el paso a paso para cumplir un objetivo.

\sphinxAtStartPar
Los diagramas de flujo usados en informática, son útiles porque ayudan a los programadores a organizar y visualizar el código haciendo más fácil encontrar y corregir errores.

\sphinxAtStartPar
A continuación se explican los símbolos que componen los diagramas de flujo del estándar (ANSI/ISO):


\bigskip\hrule\bigskip



\paragraph{Operación}
\label{\detokenize{cursos/ofimatica11/recursos/T02:operacion}}

\begin{savenotes}\sphinxattablestart
\sphinxthistablewithglobalstyle
\centering
\begin{tabulary}{\linewidth}[t]{T}
\sphinxtoprule
\sphinxtableatstartofbodyhook

\begin{sphinxfigure-in-table}
\centering
\capstart
\noindent\sphinxincludegraphics[width=250\sphinxpxdimen]{{image016}.png}
\sphinxfigcaption{Acción o Instrucción}\label{\detokenize{cursos/ofimatica11/recursos/T02:id3}}\end{sphinxfigure-in-table}\relax
\\
\sphinxbottomrule
\end{tabulary}
\sphinxtableafterendhook\par
\sphinxattableend\end{savenotes}

\sphinxAtStartPar
Se usa para describir cualquier acción o instrucción. En el interior se escribe una descripción de alguna acción o instrucción a ejecutar.

\begin{sphinxadmonition}{note}{¿Qué es una instrucción?}

\sphinxAtStartPar
Una \sphinxstylestrong{instrucción} es una orden o indicación que se da para realizar una acción específica. Puede estar presente en distintos ámbitos. Por ejemplo:
\begin{itemize}
\item {} 
\sphinxAtStartPar
\sphinxstylestrong{Una instrucción de educación:} Indicaciones que un profesor da a los estudiantes sobre cómo realizar una tarea o actividad.

\item {} 
\sphinxAtStartPar
\sphinxstylestrong{Instrucciones de programación:} Son comandos que le indican a una computadora qué hacer. Las instrucciones son parte del código fuente de un programa.

\item {} 
\sphinxAtStartPar
\sphinxstylestrong{Instrucciones de cocina:} Son pasos a seguir para una receta.

\item {} 
\sphinxAtStartPar
\sphinxstylestrong{Instrucciones técnicas:} Guías que explican cómo usar o ensamblar un producto. También aplica en el software.

\end{itemize}
\subsubsection*{Reglas para dar buenas instrucciones}
\begin{itemize}
\item {} 
\sphinxAtStartPar
\sphinxstylestrong{CLARIDAD:} Ayudan a entender qué se espera hacer.

\item {} 
\sphinxAtStartPar
\sphinxstylestrong{EFICIENCIA:} Permiten realizar tareas de manera correcta y rápida.

\item {} 
\sphinxAtStartPar
\sphinxstylestrong{SEGURIDAD:} En ciertos contextos, como el uso de maquinaria, las instrucciones son vitales para evitar accidentes.

\end{itemize}
\end{sphinxadmonition}


\bigskip\hrule\bigskip



\paragraph{Límites de proceso}
\label{\detokenize{cursos/ofimatica11/recursos/T02:limites-de-proceso}}

\begin{savenotes}\sphinxattablestart
\sphinxthistablewithglobalstyle
\centering
\begin{tabulary}{\linewidth}[t]{T}
\sphinxtoprule
\sphinxtableatstartofbodyhook

\begin{sphinxfigure-in-table}
\centering
\capstart
\noindent\sphinxincludegraphics[width=250\sphinxpxdimen]{{image025}.png}
\sphinxfigcaption{Inicio \sphinxhyphen{} Fin}\label{\detokenize{cursos/ofimatica11/recursos/T02:id4}}\end{sphinxfigure-in-table}\relax
\\
\sphinxbottomrule
\end{tabulary}
\sphinxtableafterendhook\par
\sphinxattableend\end{savenotes}

\sphinxAtStartPar
Indica el inicio (círculo blanco) y final (círculo negro) del algoritmo.


\bigskip\hrule\bigskip



\paragraph{Punto de decisión}
\label{\detokenize{cursos/ofimatica11/recursos/T02:punto-de-decision}}

\begin{savenotes}\sphinxattablestart
\sphinxthistablewithglobalstyle
\centering
\begin{tabulary}{\linewidth}[t]{T}
\sphinxtoprule
\sphinxtableatstartofbodyhook

\begin{sphinxfigure-in-table}
\centering
\capstart
\noindent\sphinxincludegraphics[width=150\sphinxpxdimen]{{image034}.png}
\sphinxfigcaption{Condicional}\label{\detokenize{cursos/ofimatica11/recursos/T02:id5}}\end{sphinxfigure-in-table}\relax
\\
\sphinxbottomrule
\end{tabulary}
\sphinxtableafterendhook\par
\sphinxattableend\end{savenotes}

\sphinxAtStartPar
Denota que en este punto se toma una decisión. Las únicas posibles opciones son SI/NO.


\bigskip\hrule\bigskip



\paragraph{Conector}
\label{\detokenize{cursos/ofimatica11/recursos/T02:conector}}

\begin{savenotes}\sphinxattablestart
\sphinxthistablewithglobalstyle
\centering
\begin{tabulary}{\linewidth}[t]{T}
\sphinxtoprule
\sphinxtableatstartofbodyhook

\begin{sphinxfigure-in-table}
\centering
\capstart
\noindent\sphinxincludegraphics[width=150\sphinxpxdimen]{{image044}.png}
\sphinxfigcaption{Conexión}\label{\detokenize{cursos/ofimatica11/recursos/T02:id6}}\end{sphinxfigure-in-table}\relax
\\
\sphinxbottomrule
\end{tabulary}
\sphinxtableafterendhook\par
\sphinxattableend\end{savenotes}

\sphinxAtStartPar
Señala que la salida de un proceso/algoritmo, puede ser la entrada de otro proceso.


\bigskip\hrule\bigskip



\paragraph{Dirección de flujo}
\label{\detokenize{cursos/ofimatica11/recursos/T02:direccion-de-flujo}}

\begin{savenotes}\sphinxattablestart
\sphinxthistablewithglobalstyle
\centering
\begin{tabulary}{\linewidth}[t]{T}
\sphinxtoprule
\sphinxtableatstartofbodyhook

\begin{sphinxfigure-in-table}
\centering
\capstart
\noindent\sphinxincludegraphics[width=150\sphinxpxdimen]{{image053}.png}
\sphinxfigcaption{Dirección}\label{\detokenize{cursos/ofimatica11/recursos/T02:id7}}\end{sphinxfigure-in-table}\relax
\\
\sphinxbottomrule
\end{tabulary}
\sphinxtableafterendhook\par
\sphinxattableend\end{savenotes}

\sphinxAtStartPar
Denota la dirección y el orden de los pasos del proceso/algoritmo.


\bigskip\hrule\bigskip



\paragraph{Entrada manual}
\label{\detokenize{cursos/ofimatica11/recursos/T02:entrada-manual}}

\begin{savenotes}\sphinxattablestart
\sphinxthistablewithglobalstyle
\centering
\begin{tabulary}{\linewidth}[t]{T}
\sphinxtoprule
\sphinxtableatstartofbodyhook

\begin{sphinxfigure-in-table}
\centering
\capstart
\noindent\sphinxincludegraphics[width=150\sphinxpxdimen]{{image063}.png}
\sphinxfigcaption{Entrada}\label{\detokenize{cursos/ofimatica11/recursos/T02:id8}}\end{sphinxfigure-in-table}\relax
\\
\sphinxbottomrule
\end{tabulary}
\sphinxtableafterendhook\par
\sphinxattableend\end{savenotes}

\sphinxAtStartPar
Permite una entrada manual del usuario para ser procesada en el algoritmo.


\bigskip\hrule\bigskip



\paragraph{Documento}
\label{\detokenize{cursos/ofimatica11/recursos/T02:documento}}

\begin{savenotes}\sphinxattablestart
\sphinxthistablewithglobalstyle
\centering
\begin{tabulary}{\linewidth}[t]{T}
\sphinxtoprule
\sphinxtableatstartofbodyhook

\begin{sphinxfigure-in-table}
\centering
\capstart
\noindent\sphinxincludegraphics[width=150\sphinxpxdimen]{{image073}.png}
\sphinxfigcaption{Documento}\label{\detokenize{cursos/ofimatica11/recursos/T02:id9}}\end{sphinxfigure-in-table}\relax
\\
\sphinxbottomrule
\end{tabulary}
\sphinxtableafterendhook\par
\sphinxattableend\end{savenotes}

\sphinxAtStartPar
Documento o registro generado en el proceso/algoritmo.


\bigskip\hrule\bigskip



\paragraph{Base de datos}
\label{\detokenize{cursos/ofimatica11/recursos/T02:base-de-datos}}

\begin{savenotes}\sphinxattablestart
\sphinxthistablewithglobalstyle
\centering
\begin{tabulary}{\linewidth}[t]{T}
\sphinxtoprule
\sphinxtableatstartofbodyhook

\begin{sphinxfigure-in-table}
\centering
\capstart
\noindent\sphinxincludegraphics[width=120\sphinxpxdimen]{{image082}.png}
\sphinxfigcaption{Base de datos}\label{\detokenize{cursos/ofimatica11/recursos/T02:id10}}\end{sphinxfigure-in-table}\relax
\\
\sphinxbottomrule
\end{tabulary}
\sphinxtableafterendhook\par
\sphinxattableend\end{savenotes}

\sphinxAtStartPar
Conjunto de tablas donde se retiene la información. En espera que se cumplan otras condiciones para continuar con el proceso/algoritmo.


\bigskip\hrule\bigskip


\sphinxstepscope


\subsubsection{T03: Algoritmos}
\label{\detokenize{cursos/ofimatica11/recursos/T03:t03-algoritmos}}\label{\detokenize{cursos/ofimatica11/recursos/T03::doc}}
\sphinxAtStartPar
%
\begin{footnote}[1]\sphinxAtStartFootnote
Trejos Buriticá, O. I. (2017). Lógica de programación. Ediciones de la U.
%
\end{footnote} Un algoritmo es un conjunto de \sphinxstylestrong{instrucciones} que le dicen a un ser humano COMO REALIZAR una tarea específica, dichas instrucciones pueden estar escritas en código o lenguaje natural.

\sphinxAtStartPar
Un algoritmo puede ser expresado con palabras o números.


\bigskip\hrule\bigskip



\paragraph{Tipos de algoritmos}
\label{\detokenize{cursos/ofimatica11/recursos/T03:tipos-de-algoritmos}}
\sphinxAtStartPar
Un algoritmo es cualquier actividad o evento con inicio y final, que tenga una serie de pasos lógicos para lograr su ejecución. Y estos suelen ser representados mediante diagramas de flujo.

\sphinxAtStartPar
Existen dos tipos de algoritmos:
\begin{itemize}
\item {} 
\sphinxAtStartPar
Algoritmos cualitativos (informales)

\item {} 
\sphinxAtStartPar
Algoritmos cuantitativos (formales)

\end{itemize}


\subparagraph{Algoritmos cualitativos (informales)}
\label{\detokenize{cursos/ofimatica11/recursos/T03:algoritmos-cualitativos-informales}}
\sphinxAtStartPar
No requieren de cálculos numéricos para su ejecución. En su lugar, se deben ejecutar secuencias lógicas. Por ejemplo: las instrucciones para ensamblar un artefacto.


\subparagraph{Algoritmos cuantitativos (formales)}
\label{\detokenize{cursos/ofimatica11/recursos/T03:algoritmos-cuantitativos-formales}}
\sphinxAtStartPar
Requieren de cálculos numéricos (por ejemplo la solución de una ecuación). Si el cálculo numérico puede ser procesado por una persona, entonces es un \sphinxstylestrong{algoritmo cuantitativo no computacional}.

\sphinxAtStartPar
Si los cálculos numéricos son muy complejos entonces requieren poder computacional. A esta clase de algoritmos se les dice que son \sphinxstylestrong{algoritmos cuantitativos computacionales}.
\begin{itemize}
\item {} 
\sphinxAtStartPar
\sphinxstylestrong{Algoritmo cuantitativo no computacional:} Estos algoritmos requieren de operaciones numéricas simples, que un ser humano puede resolver sin necesidad de usar un sistema computacional. Por ejemplo una receta de cocina, un cálculo numérico sencillo o las instrucciones para llegar a un lugar.

\item {} 
\sphinxAtStartPar
\sphinxstylestrong{Algoritmo cuantitativo computacional:} Estos algoritmos requieren de operaciones numéricas complejas que deben resolverse con dispositivos de cálculo, como una computadora o calculadora. Ecuaciones de mucha complejidad o códigos que solo pueden ser interpretados por una máquina, son ejemplos de este tipo de algoritmo.

\end{itemize}


\bigskip\hrule\bigskip



\paragraph{Metodología para desarrollar un algoritmo}
\label{\detokenize{cursos/ofimatica11/recursos/T03:metodologia-para-desarrollar-un-algoritmo}}
\sphinxAtStartPar
Cuando se desarrolla un algoritmo, cualquier usuario podría ejecutarlo siempre y cuando cumpla unos requisitos iniciales.

\sphinxAtStartPar
Para desarrollar un algoritmo, se tienen las siguientes fases:
\begin{itemize}
\item {} 
\sphinxAtStartPar
\sphinxstylestrong{FASE 01: Análisis de requisitos (objetivo y condiciones)}

\sphinxAtStartPar
En esta fase se identifica lo siguiente:
\begin{itemize}
\item {} 
\sphinxAtStartPar
\sphinxstylestrong{Objetivo:} Identificar el objetivo claro del algoritmo.

\item {} 
\sphinxAtStartPar
\sphinxstylestrong{Condiciones iniciales:} Identificar los requisitos que debe cumplir el usuario para poder desarrollar el algoritmo.

\item {} 
\sphinxAtStartPar
\sphinxstylestrong{Condiciones de incumplimiento:} Identificar todos los condicionales del algoritmo.

\item {} 
\sphinxAtStartPar
\sphinxstylestrong{Condiciones de cumplimiento:} Identificar el estado del usuario cuando se cumple el algoritmo.

\end{itemize}



\bigskip\hrule\bigskip

\item {} 
\sphinxAtStartPar
\sphinxstylestrong{FASE 02: Diseño del algoritmo (diagrama de flujo)}

\sphinxAtStartPar
Realizar un diagrama de flujo donde muestre el paso a paso lógico por medio de las siguientes figuras:
\begin{itemize}
\item {} 
\sphinxAtStartPar
\sphinxstylestrong{Círculos:} INICIO y FIN

\item {} 
\sphinxAtStartPar
\sphinxstylestrong{Rectángulos:} Acciones

\item {} 
\sphinxAtStartPar
\sphinxstylestrong{Rombos:} Condicionales

\item {} 
\sphinxAtStartPar
\sphinxstylestrong{Rombos anidados:} Bucles

\item {} 
\sphinxAtStartPar
\sphinxstylestrong{Flechas:} Orden de secuencia.

\end{itemize}

\sphinxAtStartPar
Después de tener el diagrama de flujo, realizar un análisis.



\bigskip\hrule\bigskip

\item {} 
\sphinxAtStartPar
\sphinxstylestrong{FASE 03: Construcción del algoritmo (código fuente)}

\sphinxAtStartPar
Conversión del \sphinxstylestrong{diagrama de flujo} a \sphinxstylestrong{CÓDIGO FUENTE}.
El código fuente es cualquier código que usa un lenguaje de programación que entienda una PC.
En este caso sería PYTHON pero de momento se usará PSEUDOCÓDIGO.

\sphinxAtStartPar
Se debe identificar el algoritmo. La identificación debe tener:
\begin{itemize}
\item {} 
\sphinxAtStartPar
\sphinxstylestrong{TÍTULO:} “Nombre del algoritmo”

\item {} 
\sphinxAtStartPar
\sphinxstylestrong{AUTOR:} “Creador del algoritmo”

\item {} 
\sphinxAtStartPar
\sphinxstylestrong{FECHA:} “Fecha de creación o última modificación del algoritmo”

\end{itemize}

\end{itemize}


\bigskip\hrule\bigskip



\paragraph{Algoritmo para ir al colegio (ejemplo)}
\label{\detokenize{cursos/ofimatica11/recursos/T03:algoritmo-para-ir-al-colegio-ejemplo}}
\sphinxAtStartPar
Desarrolle el siguiente ejercicio:
\subsubsection*{Algoritmo para ir al colegio}

\sphinxAtStartPar
En su cuaderno desarrolle un \sphinxstylestrong{algoritmo que permita a un estudiante ir al colegio}. Tenga en cuenta lo siguiente:
\begin{quote}

\sphinxAtStartPar
El estudiante inicialmente está en su casa. Luego debe ponerse el uniforme e ir al garaje para intentar prender la moto.

\sphinxAtStartPar
Si la moto NO prende, entonces salir del garaje, luego salir de la casa para llegar al colegio a pie.

\sphinxAtStartPar
Si la moto SI prende, entonces abrir el garaje, luego salir del garaje para llegar al colegio en moto.
\end{quote}

\sphinxAtStartPar
A continuación se muestra la solución del desarrollo del algoritmo.


\subparagraph{FASE 01: Análisis de requisitos}
\label{\detokenize{cursos/ofimatica11/recursos/T03:fase-01-analisis-de-requisitos}}
\begin{sphinxadmonition}{note}{(Objetivo y condiciones)}
\subsubsection*{OBJETIVO}
\begin{itemize}
\item {} 
\sphinxAtStartPar
Llegar al colegio.

\end{itemize}
\subsubsection*{CONDICIONES INICIALES}
\begin{itemize}
\item {} 
\sphinxAtStartPar
El estudiante debe estar en su casa

\item {} 
\sphinxAtStartPar
La casa del estudiante debe tener un garaje.

\item {} 
\sphinxAtStartPar
El estudiante no debe tener puesto el uniforme.

\item {} 
\sphinxAtStartPar
El estudiante debe tener una moto en el garaje.

\end{itemize}
\subsubsection*{CONDICIONES DE INCUMPLIMIENTO}
\begin{itemize}
\item {} 
\sphinxAtStartPar
La moto no prende.

\end{itemize}
\subsubsection*{CONDICIONES DE CUMPLIMIENTO}
\begin{itemize}
\item {} 
\sphinxAtStartPar
El estudiante está en el colegio.

\item {} 
\sphinxAtStartPar
El estudiante tiene el uniforme puesto.

\item {} 
\sphinxAtStartPar
La moto puede estar en el colegio o en el garaje.

\end{itemize}
\end{sphinxadmonition}


\subparagraph{FASE 02: Diseño del algoritmo}
\label{\detokenize{cursos/ofimatica11/recursos/T03:fase-02-diseno-del-algoritmo}}
\begin{sphinxadmonition}{note}{(Diagrama de flujo)}

\begin{figure}[H]
\centering

\noindent\sphinxincludegraphics[width=400\sphinxpxdimen]{{image092}.png}
\end{figure}
\end{sphinxadmonition}


\subparagraph{FASE 03: Construcción del algoritmo}
\label{\detokenize{cursos/ofimatica11/recursos/T03:fase-03-construccion-del-algoritmo}}
\begin{sphinxadmonition}{note}{(Código fuente)}


\begin{savenotes}\sphinxattablestart
\sphinxthistablewithglobalstyle
\centering
\begin{tabulary}{\linewidth}[t]{TT}
\sphinxtoprule
\sphinxtableatstartofbodyhook
\sphinxAtStartPar
TÍTULO
&
\sphinxAtStartPar
Algoritmo para que un estudiante pueda ir al colegio
\\
\sphinxhline
\sphinxAtStartPar
AUTOR
&
\sphinxAtStartPar
Camilo Fonseca
\\
\sphinxhline
\sphinxAtStartPar
FECHA
&
\sphinxAtStartPar
Domingo, 12 de Abril de 2026 \sphinxhyphen{} 03:38pm
\\
\sphinxbottomrule
\end{tabulary}
\sphinxtableafterendhook\par
\sphinxattableend\end{savenotes}

\begin{sphinxVerbatim}[commandchars=\\\{\},numbers=left,firstnumber=1,stepnumber=1]
 \PYG{n}{INICIO}

 \PYG{l+m+mf}{1.} \PYG{n}{Ponerse} \PYG{n}{el} \PYG{n}{uniforme}
 \PYG{l+m+mf}{3.} \PYG{n}{Ir} \PYG{n}{al} \PYG{n}{garaje}
 \PYG{l+m+mf}{4.} \PYG{n}{Intentar} \PYG{n}{prender} \PYG{n}{la} \PYG{n}{moto}
 \PYG{l+m+mf}{5.} \PYG{n}{Si} \PYG{p}{(}\PYG{n}{la} \PYG{n}{moto} \PYG{n}{SI} \PYG{n}{prende}\PYG{p}{)}
     \PYG{l+m+mf}{5.1} \PYG{n}{Abrir} \PYG{n}{el} \PYG{n}{garaje}
     \PYG{l+m+mf}{5.2} \PYG{n}{Sacar} \PYG{n}{la} \PYG{n}{moto}
     \PYG{l+m+mf}{5.3} \PYG{n}{Cerrar} \PYG{n}{el} \PYG{n}{garaje}
     \PYG{l+m+mf}{5.4} \PYG{n}{Conducir} \PYG{n}{hasta} \PYG{n}{el} \PYG{n}{colegio}
 \PYG{l+m+mf}{6.} \PYG{n}{Si} \PYG{p}{(}\PYG{n}{la} \PYG{n}{moto} \PYG{n}{NO} \PYG{n}{prende}\PYG{p}{)}
     \PYG{l+m+mf}{6.1} \PYG{n}{Salir} \PYG{k}{del} \PYG{n}{garage}
     \PYG{l+m+mf}{6.2} \PYG{n}{Salir} \PYG{n}{de} \PYG{n}{la} \PYG{n}{casa}
     \PYG{l+m+mf}{6.3} \PYG{n}{Cerrar} \PYG{n}{la} \PYG{n}{puerta}
     \PYG{l+m+mf}{6.4} \PYG{n}{Caminar} \PYG{n}{hasta} \PYG{n}{el} \PYG{n}{colegio}
 \PYG{l+m+mf}{7.} \PYG{n}{Llegar} \PYG{n}{al} \PYG{n}{colegio}
 
 \PYG{n}{FIN}
\end{sphinxVerbatim}
\end{sphinxadmonition}


\bigskip\hrule\bigskip


\sphinxstepscope


\chapter{Referencias}
\label{\detokenize{referencias:referencias}}\label{\detokenize{referencias::doc}}\begin{itemize}
\item {} 
\sphinxAtStartPar
Antonio Fernández, B. (2025, May 9). Guía de iniciación 24.8 \sphinxhyphen{} Prefacio. Libreoffice.org. https://books.libreoffice.org/es/GS248/GS24800\sphinxhyphen{}Prefacio.html

\item {} 
\sphinxAtStartPar
Andrés, R. (2019, September 1). Raspberry Pi: la historia de un mini\sphinxhyphen{}PC que se ha convertido en un éxito de masas. Computer Hoy. https://computerhoy.20minutos.es/reportajes/tecnologia/historia\sphinxhyphen{}raspberry\sphinxhyphen{}pi\sphinxhyphen{}482477

\item {} 
\sphinxAtStartPar
Andreu, A. (2024, September 17). Qué es Mozilla Firefox, diferencias con otros navegadores y extensiones imprescindibles. Computer Hoy. https://computerhoy.20minutos.es/tecnologia/mozilla\sphinxhyphen{}firefox\sphinxhyphen{}diferencias\sphinxhyphen{}otros\sphinxhyphen{}navegadores\sphinxhyphen{}extensiones\sphinxhyphen{}imprescindibles\sphinxhyphen{}1404587

\item {} 
\sphinxAtStartPar
ARPANET. (2005, Feb 10). Sistemas.com. Retrieved April 25, 2026, from https://sistemas.com/arpanet.php

\item {} 
\sphinxAtStartPar
Audacity: una forma diferente de estar en la onda. (2023, November 20). INTEF. https://intef.es/observatorio\_tecno/audacity\sphinxhyphen{}una\sphinxhyphen{}forma\sphinxhyphen{}diferente\sphinxhyphen{}de\sphinxhyphen{}estar\sphinxhyphen{}en\sphinxhyphen{}la\sphinxhyphen{}onda/

\item {} 
\sphinxAtStartPar
Catucci, A. (2020, September 9). La historia de Apple: repaso por este caso de éxito tecnológico. Marketing Insider Review. https://marketinginsiderreview.com/historia\sphinxhyphen{}apple\sphinxhyphen{}exito\sphinxhyphen{}tecnologico/

\item {} 
\sphinxAtStartPar
Corrales, R. (2024, August 19). Microsoft, el gigante de la informática: historia, Windows, todos sus servicios y la IA de Copilot. Computer Hoy. https://computerhoy.20minutos.es/industria/microsoft\sphinxhyphen{}gigante\sphinxhyphen{}informatica\sphinxhyphen{}historia\sphinxhyphen{}windows\sphinxhyphen{}todos\sphinxhyphen{}servicios\sphinxhyphen{}ia\sphinxhyphen{}copilot\sphinxhyphen{}1399819

\item {} 
\sphinxAtStartPar
Deitel, P. J., \& Deitel, H. (2014). C++: cómo programar (A. V. Romero Elizondo, Trans.). Pearson Educación.

\item {} 
\sphinxAtStartPar
Delgado, J. M. (2024, September 15). Razones por las que GIMP es la mejor alternativa a Photoshop para editar tus fotos. Computer Hoy. https://computerhoy.20minutos.es/apps/razones\sphinxhyphen{}gimp\sphinxhyphen{}mejor\sphinxhyphen{}alternativa\sphinxhyphen{}photoshop\sphinxhyphen{}editar\sphinxhyphen{}fotos\sphinxhyphen{}1400576

\item {} 
\sphinxAtStartPar
Delgado, J. M. (2024, December 15). Funciones ocultas de Microsoft Office que deberías usar ahora mismo para mejorar tu productividad. Computer Hoy. https://computerhoy.20minutos.es/apps/funciones\sphinxhyphen{}ocultas\sphinxhyphen{}microsoft\sphinxhyphen{}office\sphinxhyphen{}deberias\sphinxhyphen{}usar\sphinxhyphen{}ahora\sphinxhyphen{}mismo\sphinxhyphen{}mejorar\sphinxhyphen{}productividad\sphinxhyphen{}1428291

\item {} 
\sphinxAtStartPar
Delgado, J. M. (2025, June 14). VLC sigue muy vivo: cosas sorprendentes que puedes hacer con el mejor reproductor de vídeo para PC. Computer Hoy. https://computerhoy.20minutos.es/tecnologia/vlc\sphinxhyphen{}sigue\sphinxhyphen{}muy\sphinxhyphen{}vivo\sphinxhyphen{}cosas\sphinxhyphen{}sorprendentes\sphinxhyphen{}puedes\sphinxhyphen{}hacer\sphinxhyphen{}mejor\sphinxhyphen{}reproductor\sphinxhyphen{}video\sphinxhyphen{}pc\sphinxhyphen{}1463759

\item {} 
\sphinxAtStartPar
de Luis, E. R. (2026, March 15). En 1990 un estudiante tuvo un problema mostrando imágenes en su Mac. La solución fue crear uno de los softwares más exitosos de la historia. Xataka.com; Xataka. https://www.xataka.com/aplicaciones/1987\sphinxhyphen{}tuvo\sphinxhyphen{}problema\sphinxhyphen{}mostrando\sphinxhyphen{}imagenes\sphinxhyphen{}su\sphinxhyphen{}mac\sphinxhyphen{}asi\sphinxhyphen{}que\sphinxhyphen{}creo\sphinxhyphen{}app\sphinxhyphen{}hoy\sphinxhyphen{}editor\sphinxhyphen{}imagenes\sphinxhyphen{}usado\sphinxhyphen{}historia

\item {} 
\sphinxAtStartPar
Eadicicco, L. (2025, May 11). La tecnologí­a que definió el internet moderno está cambiando y Silicon Valley finalmente lo admite. CNN en Español. https://cnnespanol.cnn.com/2025/05/11/ciencia/tecnologia\sphinxhyphen{}internet\sphinxhyphen{}moderno\sphinxhyphen{}silicon\sphinxhyphen{}valley\sphinxhyphen{}trax

\item {} 
\sphinxAtStartPar
Hill, S. (2024, March 24). Qué es Google One y cómo saber si lo necesitas. WIRED. https://es.wired.com/articulos/que\sphinxhyphen{}es\sphinxhyphen{}google\sphinxhyphen{}one\sphinxhyphen{}y\sphinxhyphen{}como\sphinxhyphen{}saber\sphinxhyphen{}si\sphinxhyphen{}lo\sphinxhyphen{}necesitas

\item {} 
\sphinxAtStartPar
Historia de Microsoft. (2012, November 10). Blog de tecnología; Sistemas Ibertrónica. https://ibertronica.es/blog/productos/microsoft\sphinxhyphen{}pasado\sphinxhyphen{}presente\sphinxhyphen{}y\sphinxhyphen{}futuro/

\item {} 
\sphinxAtStartPar
Huertos, A. A. (2019, May 28). La historia de los lenguajes de programación. Computer Hoy. https://computerhoy.20minutos.es/reportajes/tecnologia/historia\sphinxhyphen{}lenguajes\sphinxhyphen{}programacion\sphinxhyphen{}428041

\item {} 
\sphinxAtStartPar
Inkscape, tu ilustrador de cabecera. (2021, July 19). INTEF. https://intef.es/observatorio\_tecno/inkscape\sphinxhyphen{}tu\sphinxhyphen{}ilustrador\sphinxhyphen{}de\sphinxhyphen{}cabecera/

\item {} 
\sphinxAtStartPar
Lautrec, T. (n.d.). ¿Qué es AutoCAD y para qué sirve? Toulouse Lautrec. Retrieved May 7, 2026, from https://www.toulouselautrec.edu.pe/blogs/que\sphinxhyphen{}es\sphinxhyphen{}autocad

\item {} 
\sphinxAtStartPar
Maturana, J. (2023, June 2). 25 aniversario de Computer Hoy: La historia de Windows, cómo ha sido la evolución del servicio más popular de Microsoft. Computer Hoy. https://computerhoy.20minutos.es/windows/historia\sphinxhyphen{}windows\sphinxhyphen{}como\sphinxhyphen{}ha\sphinxhyphen{}sido\sphinxhyphen{}evolucion\sphinxhyphen{}servicio\sphinxhyphen{}popular\sphinxhyphen{}microsoft\sphinxhyphen{}1237508

\item {} 
\sphinxAtStartPar
Nava, J. (2026, February 19). ¿Qué es blender? sus usos y beneficios. IMSED Business \& Tech School. https://imsed.com/postgrados/tecnologia/que\sphinxhyphen{}es\sphinxhyphen{}blender\sphinxhyphen{}usos\sphinxhyphen{}beneficios/

\item {} 
\sphinxAtStartPar
(N.d.). Libreoffice.org. Retrieved May 17, 2026, from https://es.libreoffice.org/who\sphinxhyphen{}uses\sphinxhyphen{}libreoffice/

\item {} 
\sphinxAtStartPar
Onieva, D. (2020, February 17). LibreOffice, la suite gratuita que compite cara a cara con Office. SoftZone. https://www.softzone.es/programas/oficina/libreoffice/

\item {} 
\sphinxAtStartPar
Pastor, J. (2015, January 27). Altair 8800: todo comenzó hace ya 40 años. Xataka.com; Xataka. https://www.xataka.com/historia\sphinxhyphen{}tecnologica/altair\sphinxhyphen{}8800\sphinxhyphen{}todo\sphinxhyphen{}comenzo\sphinxhyphen{}hace\sphinxhyphen{}ya\sphinxhyphen{}40\sphinxhyphen{}anos

\item {} 
\sphinxAtStartPar
Pastor, J. (2018, April 28). Cuando Mosaic dominaba el mundo (de los navegadores). Xataka.com; Xataka. https://www.xataka.com/historia\sphinxhyphen{}tecnologica/cuando\sphinxhyphen{}mosaic\sphinxhyphen{}dominaba\sphinxhyphen{}el\sphinxhyphen{}mundo\sphinxhyphen{}de\sphinxhyphen{}los\sphinxhyphen{}navegadores

\item {} 
\sphinxAtStartPar
Rohaut, S. Linux: preparación a la certificación LPIC 1 (exámenes LPI 101 y LPI 102), 3ª Edición. Barcelona: ENI ediciones, 2015.

\item {} 
\sphinxAtStartPar
Rojas, J. (2022, May 17). El primer sitio web de la historia. Un hito que cambió vidas. Telefónica. https://www.telefonica.com/es/sala\sphinxhyphen{}comunicacion/blog/el\sphinxhyphen{}primer\sphinxhyphen{}sitio\sphinxhyphen{}web\sphinxhyphen{}de\sphinxhyphen{}la\sphinxhyphen{}historia\sphinxhyphen{}un\sphinxhyphen{}hito\sphinxhyphen{}que\sphinxhyphen{}cambio\sphinxhyphen{}vidas/

\item {} 
\sphinxAtStartPar
Superior, E. F. P. (2023, August 31). ¿Qué es la ofimática y para qué sirve?: herramientas. Esic.edu; ESIC. https://www.esic.edu/business/que\sphinxhyphen{}es\sphinxhyphen{}la\sphinxhyphen{}ofimatica\sphinxhyphen{}para\sphinxhyphen{}que\sphinxhyphen{}sirve\sphinxhyphen{}c

\item {} 
\sphinxAtStartPar
Textos: Jorge Eliécer Lozano / Ilustraciones: Gissell Andreína García Rodríguez. (n.d.). El software pirata representa un riesgo inminente. Edu.co. Retrieved May 7, 2026, from https://notired.univalle.edu.co/el\sphinxhyphen{}software\sphinxhyphen{}pirata\sphinxhyphen{}representa\sphinxhyphen{}un\sphinxhyphen{}riesgo\sphinxhyphen{}inminente/

\item {} 
\sphinxAtStartPar
The Conversation. (2024, January 22). La Mac cumple 40 años: debutó con este antológico aviso dirigido por Ridley Scott para mostrar cómo había otra forma de usar una computadora. LA NACION. https://www.lanacion.com.ar/tecnologia/la\sphinxhyphen{}mac\sphinxhyphen{}cumple\sphinxhyphen{}40\sphinxhyphen{}anos\sphinxhyphen{}debuto\sphinxhyphen{}con\sphinxhyphen{}este\sphinxhyphen{}antologico\sphinxhyphen{}aviso\sphinxhyphen{}dirigido\sphinxhyphen{}por\sphinxhyphen{}ridley\sphinxhyphen{}scott\sphinxhyphen{}para\sphinxhyphen{}mostrar\sphinxhyphen{}como\sphinxhyphen{}nid22012024/

\item {} 
\sphinxAtStartPar
Tipos de licencias de software: software libre, propietario y demás. (2022, July 19). Velneo.com. https://velneo.com/blog/licencias\sphinxhyphen{}software\sphinxhyphen{}libre\sphinxhyphen{}propietario\sphinxhyphen{}otros/

\item {} 
\sphinxAtStartPar
Trejos Buriticá, O. I. (2017). Lógica de programación. Ediciones de la U.

\item {} 
\sphinxAtStartPar
Universidad Nacional del Nordeste. (n.d.). Informática. Conceptos fundamentales. https://exa.unne.edu.ar/ingenieria/computacion/Tema1.pdf

\item {} 
\sphinxAtStartPar
Velásquez, W. (2022, April 1). 3 historias del software libre que deberías conocer. Wajari. https://wajari.com/blog/historias\sphinxhyphen{}software\sphinxhyphen{}libre\sphinxhyphen{}open\sphinxhyphen{}source/

\item {} 
\sphinxAtStartPar
(Visnuh), M. G. (2024, January 24). Con Microsoft Office 365 dominarás la productividad y con el precio de esta licencia en oferta no tendrás que empeñarte para ello. Xataka.com; Xataka. https://www.xataka.com/seleccion/microsoft\sphinxhyphen{}office\sphinxhyphen{}365\sphinxhyphen{}dominaras\sphinxhyphen{}productividad\sphinxhyphen{}precio\sphinxhyphen{}esta\sphinxhyphen{}licencia\sphinxhyphen{}oferta\sphinxhyphen{}no\sphinxhyphen{}tendras\sphinxhyphen{}que\sphinxhyphen{}empenarte\sphinxhyphen{}para\sphinxhyphen{}ello

\end{itemize}

\sphinxstepscope


\chapter{Descargas}
\label{\detokenize{descargas:descargas}}\label{\detokenize{descargas::doc}}
\sphinxAtStartPar
Identifique el recurso y PULSE el botón \sphinxstylestrong{Descargar}.
\begin{itemize}
\item {} 
\sphinxAtStartPar
{\hyperref[\detokenize{descargas::doc}]{\sphinxcrossref{\DUrole{std}{\DUrole{std-doc}{\DUrole{sd-sphinx-override}{\DUrole{sd-badge}{\DUrole{sd-bg-primary}{\DUrole{sd-bg-text-primary}{Descargar}}}}}}}}} Servidor principal

\item {} 
\sphinxAtStartPar
{\hyperref[\detokenize{descargas::doc}]{\sphinxcrossref{\DUrole{std}{\DUrole{std-doc}{\DUrole{sd-sphinx-override}{\DUrole{sd-badge}{\DUrole{sd-bg-secondary}{\DUrole{sd-bg-text-secondary}{Descargar}}}}}}}}} Servidor secundario

\end{itemize}


\bigskip\hrule\bigskip



\section{Institucional}
\label{\detokenize{descargas:institucional}}\begin{itemize}
\item {} 
\sphinxAtStartPar
MANUAL DE CONVIVENCIA \DUrole{xref}{\DUrole{download}{\DUrole{myst}{\DUrole{sd-sphinx-override}{\DUrole{sd-badge}{\DUrole{sd-bg-primary}{\DUrole{sd-bg-text-primary}{Descargar}}}}}}} \sphinxhref{https://drive.google.com/file/d/101ObkDOMInftIxeh4Q\_yyKYM9gwauDN7/view?usp=drive\_link}{\DUrole{sd-sphinx-override}{\DUrole{sd-badge}{\DUrole{sd-bg-secondary}{\DUrole{sd-bg-text-secondary}{Descargar}}}}}

\item {} 
\sphinxAtStartPar
PORTAL: aula 115 (PDF) \DUrole{xref}{\DUrole{download}{\DUrole{myst}{\DUrole{sd-sphinx-override}{\DUrole{sd-badge}{\DUrole{sd-bg-primary}{\DUrole{sd-bg-text-primary}{Descargar}}}}}}}

\end{itemize}


\bigskip\hrule\bigskip



\section{Proyecto de grado}
\label{\detokenize{descargas:proyecto-de-grado}}\begin{itemize}
\item {} 

\begin{savenotes}\sphinxattablestart
\sphinxthistablewithglobalstyle
\raggedright
\begin{tabulary}{\linewidth}[t]{TTTTT}
\sphinxtoprule
\sphinxstyletheadfamily 
\sphinxAtStartPar
Plantilla
&\sphinxstyletheadfamily 
\sphinxAtStartPar
Microsoft Word
&\sphinxstyletheadfamily 
\sphinxAtStartPar
LibreOffice Writer
&\sphinxstyletheadfamily 
\sphinxAtStartPar
Google Docs
&\sphinxstyletheadfamily 
\sphinxAtStartPar
LaTeX
\\
\sphinxmidrule
\sphinxtableatstartofbodyhook
\sphinxAtStartPar
Marco teórico
&&
\sphinxAtStartPar
\DUrole{xref}{\DUrole{download}{\DUrole{myst}{\DUrole{sd-sphinx-override}{\DUrole{sd-badge}{\DUrole{sd-bg-primary}{\DUrole{sd-bg-text-primary}{Descargar}}}}}}} \sphinxhref{https://drive.google.com/file/d/1CgOI-NqmOQrc3K\_oq9lFJrJ5Q9hqbKQH/view?usp=drive\_link}{\DUrole{sd-sphinx-override}{\DUrole{sd-badge}{\DUrole{sd-bg-secondary}{\DUrole{sd-bg-text-secondary}{Descargar}}}}}
&&\\
\sphinxhline
\sphinxAtStartPar
Anteproyecto
&&
\sphinxAtStartPar
\DUrole{xref}{\DUrole{download}{\DUrole{myst}{\DUrole{sd-sphinx-override}{\DUrole{sd-badge}{\DUrole{sd-bg-primary}{\DUrole{sd-bg-text-primary}{Descargar}}}}}}} \sphinxhref{https://drive.google.com/file/d/1klPSl7wJwT1KyA956JvUN11-gPWVMJKR/view?usp=drive\_link}{\DUrole{sd-sphinx-override}{\DUrole{sd-badge}{\DUrole{sd-bg-secondary}{\DUrole{sd-bg-text-secondary}{Descargar}}}}}
&&\\
\sphinxhline
\sphinxAtStartPar
Proyecto
&
\sphinxAtStartPar
\DUrole{xref}{\DUrole{download}{\DUrole{myst}{\DUrole{sd-sphinx-override}{\DUrole{sd-badge}{\DUrole{sd-bg-primary}{\DUrole{sd-bg-text-primary}{Descargar}}}}}}} \sphinxhref{https://docs.google.com/document/d/1Ela1bpCYS60FAh82kLc8D9H6ZBJNU8nZ/edit?usp=drive\_link\&amp;ouid=112729153804362979178\&amp;rtpof=true\&amp;sd=true}{\DUrole{sd-sphinx-override}{\DUrole{sd-badge}{\DUrole{sd-bg-secondary}{\DUrole{sd-bg-text-secondary}{Descargar}}}}}
&&&\\
\sphinxbottomrule
\end{tabulary}
\sphinxtableafterendhook\par
\sphinxattableend\end{savenotes}

\end{itemize}



\renewcommand{\indexname}{Índice}
\printindex
\end{document}